\chapter{Binary \& Variable Stars}

\section*{Theory problems}

\probhead{8.1}{Binary System}{2007}{Chiang Mai, Thailand}{TH Q3}{10}
% source id: 2007_TH_Q3   tag: binary mass transfer, period change
A binary star system consists of $M_1$ and $M_2$ separated by a distance $D$.
$M_1$ and $M_2$ are revolving with an angular velocity $\omega$ in circular
orbits about their common centre of mass. Mass is continuously being
transferred from one star to the other. This transfer of mass causes their
orbital period and their separation to change slowly with time.

In order to simplify the analysis, we will assume that the stars are like
point particles and that the effects of the rotation about their own axes are
negligible.

\begin{description}
\item[(a)] What is the total angular momentum and kinetic energy of the
  system? \hfill(2 points)
\item[(b)] Find the relation between the angular velocity $\omega$ and the
  distance $D$ between the stars. \hfill(2 points)
\item[(c)] In a time duration $\Delta t$, a mass transfer between the two
  stars results in a change of mass $\Delta M_1$ in star $M_1$; find the
  quantity $\Delta\omega$ in terms of $\omega$, $M_1$, $M_2$ and
  $\Delta M_1$. \hfill(3 points)
\item[(d)] In a certain binary system, $M_1 = 2.9\,M_\odot$,
  $M_2 = 1.4\,M_\odot$ and the orbital period, $T = 2.49$ days. After 100
  years, the period $T$ has increased by 20\,s. Find the value of
  $\dfrac{\Delta M_1}{M_1 \Delta t}$ (in the unit ``per year'').
  \hfill(1.5 points)
\item[(e)] In which direction is mass flowing, from $M_1$ to $M_2$, or $M_2$
  to $M_1$? \hfill(0.5 point)
\item[(f)] Find also the value of $\dfrac{\Delta D}{D \Delta t}$ (in the
  unit ``per year''). \hfill(1 point)
\end{description}

You may use these approximations:
\[ (1+x)^n \approx 1 + nx, \quad \mbox{when } x \ll 1; \]
\[ (1+x)(1+y) \approx 1 + x + y, \quad \mbox{when } x \ll 1,\ y \ll 1. \]

\probhead{8.2}{}{2008}{Bandung, Indonesia}{TH L1}{100}
% source id: 2008_TH_L1   tag: eclipsing binary radii, masses
An eclipsing binary star system has a period of 30 days. The light curve in
the figure below shows that the secondary star eclipses the primary star
(from point A to point D) in 8 hours (measured from the time of first contact
to final contact), whereas from point B to point C, the total eclipse period
is 1 hour and 18 minutes. The spectral analysis yields the maximum radial
velocity of the primary star to be 30\,km/s and of the secondary star to be
40\,km/s. If we assume that the orbits are circular and have an inclination
of $i = 90^\circ$, determine the radii and the masses of both stars in units
of solar radius and solar mass.
% sic: the source prints "has an inclination" and "in unit of solar radius";
% grammar normalized in transcription.
\ioaafig{2008_TH_L1-f1.pdf}

\probhead{8.3}{}{2009}{Tehran, Iran}{TH P15}{10}
% source id: 2009_TH_P15   tag: prolate-spheroid eclipsing binary amplitude
An eclipsing close binary system consists of two giant stars with the same
sizes. As a result of mutual gravitational force, stars are deformed from
perfect sphere to the prolate spheroid with $a = 2b$, where $a$ and $b$ are
semi-major and semi-minor axes (the major axes are always co-linear). The
inclination of the orbital plane to the plane of sky is $90^\circ$. Calculate
the amplitude of light variation in magnitude ($\Delta m$) as a result of the
orbital motion of two stars. Ignore temperature variation due to tidal
deformation and limb darkening on the surface of the stars.

Hint: A prolate spheroid is a geometrical shape made by rotating of an
ellipse around its major axis, like rugby ball or melon.

\probhead{8.4}{}{2010}{Beijing, China}{TH S4}{10}
% source id: 2010_TH_S4   tag: visual binary component masses via Kepler
A binary system is 10\,pc away, the largest angular separation between the
components is $7.0''$, the smallest is $1.0''$. Assume that the orbital
period is 100 years, and that the orbital plane is perpendicular to the line
of sight. If the semi-major axis of the orbit of one component corresponds to
$3.0''$, that is $a_1 = 3.0''$, estimate the mass of each component of the
binary system, in terms of solar mass.

\probhead{8.5}{}{2011}{Katowice, Poland}{TH L3}{30}
% source id: 2011_TH_L3   tag: two stars gravitationally bound?
The planetarium program `Guide' gives the following data for two solar mass
stars:

\begin{center}
\begin{tabular}{lll}
Star & 1 & 2 \\
Right Ascension & $14^\mathrm{h}\,29^\mathrm{m}\,44.95^\mathrm{s}$
  & $14^\mathrm{h}\,39^\mathrm{m}\,39.39^\mathrm{s}$ \\
Declination & $-62^\circ\,40'\,46.14''$ & $-60^\circ\,50'\,22.10''$ \\
Distance & 1.2953\,pc & 1.3475\,pc \\
Proper motion in R.A. & $-3.776$\,arcsec/year & $-3.600$\,arcsec/year \\
Proper motion in Dec. & 0.95\,arcsec/year & 0.77\,arcsec/year \\
\end{tabular}
\end{center}

Based on these data, determine whether these stars form a gravitationally
bound system. Assume the stars are on the main sequence. Write YES if bound
% sic: the source prints "Write YES of bound"
or NO if not bound alongside your final calculation.

Note: the proper motion in R.A. has been corrected for the declination of
the stars.

\probhead{8.6}{}{2013}{Volos, Greece}{TH S4}{10}
% source id: 2013_TH_S4   tag: Cepheid P-L distance with extinction
Figure 2 shows the relation between absolute magnitude and period for
classical cepheids. Figure 3 shows the light curve (apparent magnitude versus
time in days) of a classical cepheid in a local group galaxy.

\begin{description}
\item[(a)] Using these two figures estimate the distance of the cepheid from
  us.
\item[(b)] Revise your estimate assuming that the interstellar extinction
  towards the cepheid is $A = 0.25$\,mag.
\end{description}
\ioaafig{2013_TH_S4-f1.pdf}
\ioaafig{2013_TH_S4-f2.pdf}

\probhead{8.7}{Absolute magnitude of a cepheide}{2014}{Suceava, Romania}{TH P15}{}
% source id: 2014_TH_P15   tag: derive Cepheid period-luminosity relation
The cepheides are variable stars, whose luminosities vary due to volume
oscillations. The period of the oscillations of a cepheide star is:
\[ P = 2\pi R \sqrt{\frac{R}{KM}}, \]
where $R$ is the radius of the cepheide; $M$ is the mass of the cepheid
(constant during oscillation); $R = R(t)$; $P = P(t)$.

Demonstrate that the absolute magnitude of the cepheide, $M_{\mathrm{cef}}$,
depends on the period of the cepheide's pulsation $P$ according to the
following relation:
\[ M_{\mathrm{cef}} = -2.5^{\mathrm{m}} \cdot \log k
   - \left(\frac{10}{3}\right)^{\mathrm{m}} \cdot \log P, \]
where $k$ is a constant; $P = P(t)$;
$M_{\mathrm{cef}} = M_{\mathrm{cef}}(t)$.


\probhead{8.8}{Mass function of a visual binary stellar system}{2014}{Suceava, Romania}{TH P4}{}
% source id: 2014_TH_P4   tag: mass function from RV amplitude
For a visual binary stellar system consisting of the stars $\sigma_1$ and
$\sigma_2$, the following relation represents the mass function of the
system:
\[ f(M_1; M_2) = \frac{M_2^3 \sin^3 i}{(M_1 + M_2)^2}, \]
where $M_1$ is the mass of star $\sigma_1$, $M_2$ is the mass of star
$\sigma_2$ and $i$ is the angle between the plane of the stars' orbits and a
plane perpendicular on the direction of observation.

The recorded spectrum of radiations emitted by the star $\sigma_1$, during
several months, reveals a sinusoidal variation of radiation wavelength, with
the period $T = 7$ days and a shift factor
$z = (\Delta\lambda)/\lambda = 0.001$.

\begin{description}
\item[(a)] Prove that the mass function of the system is:
  \[ f(M_1; M_2) = \frac{M_2^3 \sin^3 i}{(M_1 + M_2)^2}
     = \frac{T}{2\pi K} \left( v_1 \cdot \sin i \right)^3, \]
  where $v_1 \cdot \sin i$ is the maximum speed of star $\sigma_1$ relatively
  to the observer; $K$ is the gravitational constant; $i$ is the angle
  between the plane of the orbits and the plane normal to the observation
  direction. Assumptions: the orbits of the stars are circular.
\item[(b)] Derive an expression for the mass function of the system. The
  following values are known: $c = 3 \times 10^8$\,m/s;
  $K = 6.67 \times 10^{-11}\,\mathrm{N\,m^2\,kg^{-2}}$.
\end{description}


\probhead{8.9}{}{2015}{Magelang, Indonesia}{TH L2}{}
% source id: 2015_TH_L2   tag: binary RV curve, mass function, BH condition
Two massive stars A and B with masses $m_A$ and $m_B$ are separated by a
distance $d$. Both stars orbit around their center of mass under
gravitational force. Assume their orbits are circular and lie on the
$X$--$Y$ plane whose origin is at the stars' center of mass (see Figure 2).

\begin{description}
\item[(a)] Find the expressions for the tangential and angular speeds of
  star A.
\end{description}

An observer standing on the $Y$--$Z$ plane (see Figure 2) sees the stars from
a large distance with an angle $\theta$ relatively to the $Z$-axis. He
measures that the velocity component of star A to his line of sight has the
form $K\cos(\omega t + \varepsilon)$, where $K$ and $\varepsilon$ are
positive.

\begin{description}
\item[(b)] Express $K^3/\omega G$ in terms of $m_A$, $m_B$, and $\theta$,
  where $G$ is the universal gravitational constant.
\end{description}

Assume that the observer then identifies that star A has mass equal to
$30\,M_S$ where $M_S$ is the Sun's mass. In addition, he observes that star B
produces X-rays and then realizes that it could be a neutron star or a black
hole. This conclusion would depend on $m_B$, i.e.: i) if $m_B < 2M_S$, then B
is a neutron star; ii) if $m_B > 2M_S$, then B is a black hole.

\begin{description}
\item[(c)] A measurement by the observer shows that
  $\dfrac{K^3}{\omega G} = \dfrac{1}{250}\,M_S$. In practice, the value of
  $\theta$ is usually not known. What is the condition on $\theta$ for star B
  to be a black hole?
\end{description}
\ioaafig{2015_TH_L2-f1.pdf}

\probhead{8.10}{Cepheid Pulsations}{2016}{Bhubaneswar, India}{TH T6}{20}
% source id: 2016_TH_T6   tag: beta Dor radii from Wien peaks, Doppler
The star $\beta$-Doradus is a Cepheid variable star with a pulsation period
of 9.84 days. We make a simplifying assumption that the star is brightest
when it is most contracted (radius being $R_1$) and it is faintest when it is
most expanded (radius being $R_2$). For simplicity, assume that the star
maintains its spherical shape and behaves as a perfect black body at every
instant during the entire cycle. The bolometric magnitude of the star varies
from 3.46 to 4.08. From Doppler measurements, we know that during pulsation
the stellar surface expands or contracts at an average radial speed of
$12.8\,\mathrm{km\,s^{-1}}$. Over the period of pulsation, the peak of
thermal radiation (intrinsic) of the star varies from 531.0\,nm to
649.1\,nm.

\begin{description}
\item[(T6.1)] Find the ratio of radii of the star in its most contracted and
  most expanded states ($R_1/R_2$). \hfill[7]
\item[(T6.2)] Find the radii of the star (in metres) in its most contracted
  and most expanded states ($R_1$ and $R_2$). \hfill[3]
\item[(T6.3)] Calculate the flux of the star, $F_2$, when it is in its most
  expanded state. \hfill[5]
\item[(T6.4)] Find the distance to the star, $D_{\mathrm{star}}$, in
  parsecs. \hfill[5]
\end{description}

\probhead{8.11}{Identification of light curves of types of selected variable stars}{2019}{Keszthely, Hungary}{TH 11}{25}
% source id: 2019_TH_11   tag: light-curve type identification
TESS (Transiting Exoplanet Survey Satellite) is NASA's most recent exoplanet
hunter mission. It has been surveying the southern sky to locate exoplanets
around the brightest and closest stars, along with a large number of
time-variable phenomena (pulsating and eclipsing variable stars, supernovae,
stellar flares, and asteroids among others).

On the separate sheet you will find graphs with light curve plots of 8
periodic variable stars (intrinsic, eclipsing and rotational; namely FO Eri,
RW Dor, 24 Eri, TIC 147272181, ST Pic, UY Eri, VV Ori, AH Col) from the TESS
target list, numbered from 1 to 8. The horizontal axis is
$\mathrm{BJD} - 2\,400\,000$ in days, and the vertical axis is $V$ in
magnitudes, where BJD is the Barycentric Julian Date, and $V$ is the visual
brightness.

\begin{description}
\item[(a)] Below is a list of variable star types. Match each light curve of
  the star with its variability type by writing its number into the rectangle
  on the answer sheet. \hfill[8\,p]
  \begin{itemize}
  \item Heart-beat star
  \item RR Lyrae type (RRab subclass) pulsating variable star
  \item Eclipsing binary of Algol type (semi-detached) with a pulsating
    component
  \item $\alpha^2$ CVn pulsating variable star
  \item W Vir type (Population II) Cepheid pulsating variable star
  \item Detached eclipsing binary with strong reflection effect
  \item Contact eclipsing binary of W UMa type
  \item Rotationally variable (spotted) star
  \end{itemize}
\item[(b)] Based on the light curve plots, estimate the periods of each
  variable star in days. Give your answers on the answer sheet up to 2
  decimal places. Periods in the range $\pm5$\% of the true periods are
  acceptable. \hfill[16\,p]
\item[(c)] What type of astronomical object produced the TESS light curve
  shown in the figure below? Mark the letter of your answer with $\times$ on
  the answer sheet. \hfill[1\,p]
  \begin{description}
  \item[(A)] Microlensing event
  \item[(B)] Saturated galaxy due to the proximity of Mars
  \item[(C)] Comet's coma close to the edge of the camera field
  \item[(D)] Supernova in a distant galaxy
  \item[(E)] Superflare on a supergiant star
  \end{description}
\end{description}
\ioaafig{2019_TH_11-f1.pdf}

\probhead{8.12}{Amplitude variation of RR Lyrae stars}{2019}{Keszthely, Hungary}{TH 8}{20}
% source id: 2019_TH_8   tag: Blazhko effect, pulsation amplitude
Hungarian astronomers significantly contributed to the study of pulsating
variable stars of RR Lyrae type, which show cyclic amplitude modulation in
their light variation (Blazhko effect).

The light variation of an RR Lyrae star is observed at two different
wavelengths: $\lambda_1 = 500$\,nm and $\lambda_2 = 2000$\,nm. At each
wavelength, we see into the star at a different depth. We refer to the depths
as layer 1 and layer 2. One can approximate the light intensity of a star at
a given wavelength by the radius and the black-body radiation intensity at
the appropriate depth. Moreover, the Wien approximation can be used for the
black-body radiation to calculate the emitted power from unit surface area:
\[ F(\lambda, T) \propto \frac{1}{\lambda^5}
   \exp\left(-\frac{hc}{k\lambda T}\right), \]
where $h$ is the Planck constant, $k$ is the Boltzmann constant, and $c$ is
the speed of light. To simplify the calculations we introduce a new constant
$C_b = hc/k \approx 0.0144$\,m\,K.

\begin{description}
\item[(a)] Assume that the temperature varies between $T_1 = 6000$\,K and
  $T_2 = 7400$\,K in each of the layers and ignoring radius variation, what
  is the ratio of the amplitudes of variation in magnitudes at the two
  wavelengths? \hfill[5\,p]
\item[(b)] What is the peak to peak amplitude of the light curve at
  $\lambda_1$? Use the magnitude scale. \hfill[3\,p]
\item[(c)] Ignoring temperature variation, what is the contribution of the
  radius variation to the light-curve amplitude for a given wavelength, if
  $R_{\min} = 0.9\,\langle R\rangle$ and
  $R_{\max} = 1.05\,\langle R\rangle$? $\langle R\rangle$ is the mean radius
  of the given layer. \hfill[3\,p]
\item[(d)] Recent observations and models show that the radius of the
  photosphere is only minimally modulated during the Blazhko cycle; however,
  the temperature variation is significant. As a result, the amplitude of
  light curve itself keeps varying. Let us assume that during the minimum
  amplitude of pulsation the temperature variation is reduced to
  $T_{\min} = 6100$\,K and $T_{\max} = 6900$\,K. What is the modulation
  amplitude at the two wavelengths? For the maximum amplitude, use the
  temperature values given in part (a). \hfill[5\,p]
\item[(e)] Which statement is correct? (There can be more than one correct
  selection.) Mark the letter of your answer with $\times$ on the answer
  sheet. \hfill[4\,p]
  \begin{description}
  \item[(A)] It is easier to observe the Blazhko effect in the infrared
    band.
  \item[(B)] The temperature variation dominates the visual light curve.
  \item[(C)] If we neglect the radius variation, the amplitude is inversely
    proportional to the wavelength.
  \item[(D)] Multi-colour observations are not useful in understanding the
    Blazhko effect.
  \end{description}
\end{description}

\probhead{8.13}{Light Curves}{2020}{GeCAA (Estonia, remote)}{TH 4}{20}
% source id: 2020_TH_4   tag: eclipsing-binary light-curve changes
The light curve A shown below, shows a fictional edge-on eclipsing binary
system containing stars X (radius $r_X$, luminosity $L_X$) and Y (radius
$r_Y$, luminosity $L_Y$). Assume that star X is brighter, but star Y is
hotter.

\begin{description}
\item[(a)] (1 point) Which of the two stars is likely to be on the main
  sequence? (Write ``X'' or ``Y'')
\item[(b)] Based on light curve A, estimate:
  \begin{description}
  \item[(I)] (2 points) $\dfrac{r_X}{r_Y}$, the ratio of the radii of the
    two stars.
  \item[(II)] (2 points) $\dfrac{L_X}{L_Y}$, the ratio of the luminosity of
    the two stars.
  \end{description}
\item[(c)] (15 points) For light curves B to F, in each case only one
  parameter of the binary system has been changed from the case in light
  curve A. For each case, choose the description from the following list
  that best corresponds to the change (write the appropriate roman numeral
  in the answer sheet).
  \begin{description}
  \item[(i)] Star X increased in size.
  \item[(ii)] Star X increased in luminosity.
  \item[(iii)] Star X decreased in size.
  \item[(iv)] Star X decreased in luminosity.
  \item[(v)] Star Y increased in size.
  \item[(vi)] Star Y increased in luminosity.
  \item[(vii)] Star Y decreased in size.
  \item[(viii)] Star Y decreased in luminosity.
  \item[(ix)] Star X is a variable star.
  \item[(x)] Star Y is a variable star.
  \item[(xi)] The inclination of the system relative to the Earth has
    changed.
  \item[(xii)] The distance of the system from the Earth has decreased.
  \item[(xiii)] The distance of the system from the Earth has increased.
  \item[(xiv)] The orbital period of the system increased.
  \item[(xv)] The orbital period of the system decreased.
  \end{description}
\end{description}
\ioaafig{2020_TH_4-f1.pdf}
\ioaafig{2020_TH_4-f2.pdf}
\ioaafig{2020_TH_4-f3.pdf}

\probhead{8.14}{Menkalinan (beta Aurigae)}{2021}{Bogot\'a (remote)}{TH Q7}{13}
% source id: 2021_TH_Q7   tag: spectroscopic binary masses, luminosities
Almost half of the stars that we see are either binary or multiple star
systems. A well-known example of this is Menkalinan (Beta Aurigae), which was
initially thought to be a single star, but today recognised as a binary
system comprising two stars that we will refer to as Menkalinan A and B. In
the following figure, a spectrum of the system (obtained by the observatory
of the Universidad de los Andes, in Bogot\'a) is shown: a spectrum of the
Menkalinan binary system in the region of $H_\alpha$; the $y$-axis is for the
relative flux, and the $x$-axis measures wavelengths. Menkalinan A is marked
as A in the graph, and Menkalinan B is marked as B.

Answer the following questions using the plot and noting that the wavelength
of the $H_\alpha$ line in the laboratory frame is 656.28\,nm. Assume circular
orbits, and assume that the binary system as a whole is at rest with respect
to the observer.

\begin{description}
\item[(7.1)] In the spectrum, we can see the $H_\alpha$ line for each star
  in the system. Calculate the line-of-sight velocity of each star (km/s)
  and determine, at the time of this observation, which of the two stars is
  moving towards us. \hfill[5.0\,pt]
\item[(7.2)] The binary system is located 81.1 light years from Earth and
  has an orbital period of 3.96 days. The semi-major axis for Menkalinan B
  (smaller star) was measured to be 3.35 milliarcseconds. If the mass ratio
  of the two components is 1.026, find the total mass of the system (in
  solar masses). \hfill[4.0\,pt]
\item[(7.3)] Calculate the individual masses of Menkalinan A and B in solar
  masses. \hfill[2.0\,pt]
\item[(7.4)] Since Menkalinan A and B are main sequence stars, use the
  relation:
  \[ \frac{L}{L_\odot} = \left(\frac{M}{M_\odot}\right)^{3.5} \]
  to estimate the luminosity of each star (in solar luminosity).
  \hfill[2.0\,pt]
\end{description}
\ioaafig{2021_TH_Q7-f1.pdf}

\probhead{8.15}{Photometry of Binary stars}{2022}{Kutaisi, Georgia}{TH 6}{20}
% source id: 2022_TH_6   tag: binary photometry, masses, extinction
We observe a binary system, at a distance $d = 89$\,pc, with a circular,
edge-on orbit around the common center of mass and period of 100 days. Our
space-based telescope is of diameter $D = 10$\,m and operates at a wavelength
$\lambda = 364$\,nm. We notice that for a total of 38 days during each full
period, the two stars cannot be resolved by this telescope as separate
objects. The wavelength of peak emission for star A is
$\lambda_A = 500$\,nm and that for star B is $\lambda_B = 600$\,nm.

\begin{description}
\item[(a)] What are the temperatures of the stars A and B, ($T_A$, $T_B$)?
  \hfill(2pt)
\item[(b)] Calculate the distance $l$ between the stars. \hfill(6pt)
\item[(c)] Calculate the sum of the masses of the stars ($M_T$). \hfill(2pt)
\end{description}

Combined photometry of the system:

\begin{center}
\begin{tabular}{clcccc}
 & Configuration & $U_0$ & $(U-B)$ & $(B-V)$ & BC \\
1 & Stars next to each other & 6.39 & 0.2 & 0.1 & 0.1 \\
2 & B transiting in front of A & 6.86 & 0.25 & 0.12 & 0.175 \\
\end{tabular}
\end{center}

The interstellar extinction per kpc of $U$ is $a_U = 1.4$\,mag/kpc. The
ratio of the densities of the stars is $\rho_A/\rho_B = 0.7$.

Note: for bolometric correction we use the convention:
\[ \mathrm{BC} = m_{\mathrm{bol}} - m_V \]

\begin{description}
\item[(d)] Calculate the masses of both the stars ($M_A$, $M_B$).
  \hfill(10pt)
\end{description}

\probhead{8.16}{Second eclipse}{2023}{Katowice, Poland}{TH Q9}{20}
% source id: 2023_TH_Q9   tag: predict secondary eclipse light curves
For each of two eclipsing binary systems, Bolek and Lolek, the primary
eclipses were observed with very high cadence as depicted in the figures
below (Figure 1: observed lightcurve for system Bolek; Figure 2: observed
lightcurve for system Lolek).

In the figures, $t$ is the time in hours relative to the moment of minimum
and $V$ is the brightness in the $V$ (visible) band in magnitudes. The points
are the measurements and the line is the fitted model of the shape of the
eclipse.

You can assume that in both cases the eclipses are central ($i = 90^\circ$)
and last for a very small fraction of the orbital period, limb darkening is
negligible, and the orbits have low eccentricity.

On the Answer Sheet, draw the predicted shape of the light curve for each of
the secondary eclipses. Write down the equations and calculations leading to
your predictions.
\ioaafig{2023_TH_Q9-f1.pdf}
\ioaafig{2023_TH_Q9-f2.pdf}
\ioaafig{2023_TH_Q9-f3.pdf}

\probhead{8.17}{Binary Hardening}{2024}{Rio de Janeiro, Brazil}{TH T8}{25}
% source id: 2024_TH_T8   tag: BH binary hardening by stellar slingshots
Consider a binary system of black holes, both of equal mass $M$, separated by
a distance $a$, and revolving around their common centre of mass (CM) in
circular orbits. This binary system moves against, and interacts with, a very
large, uniform field of stars (each of mass $m \ll M$) with number density
$n$. Consider a star that approaches the system from infinity with speed $v$
and impact parameter $b$, in the reference frame of the CM (as shown in the
figure below). Its closest approach distance to the CM is
$r_p \approx \frac{1}{2}a$. For tasks (a) and (c), you should make use of the
fact that $v^2 \ll \frac{GM}{a}$.

\begin{description}
\item[(a)] (5 points) Obtain an expression for $b$, in terms of $M$, $a$,
  $v$, and physical constants. In this task, assume that the star interacts
  with the binary as if its total mass was fixed at the CM.
\end{description}

After a complex interaction with the binary, the star is slingshot from the
system. The exact calculation of its ejection speed is complex, but the
result can be estimated by considering that the star only interacts with one
of the components when near the system. As such, consider, in part (b),
\emph{only} the gravitational interaction between the star and \emph{one of
the components} in the binary.

\begin{description}
\item[(b)] (6 points) The star approaches the component with an initial
  speed negligible compared to the component's orbital speed, and both are
  moving directly towards each other. After interacting with the system,
  when the star is again far away from the black hole, we find that the
  direction of its velocity vector is reversed and the final speed is
  $v_f$. Determine $v_f$, in terms of $M$, $a$, and physical constants.
  Assume that linear momentum and mechanical energy are conserved during
  this interaction and that it takes place in a timescale much smaller than
  the binary's period. Recall that $m \ll M$.
\end{description}

For the following task, assume that all stars approaching the system from
infinity with an impact parameter $\frac{1}{2}b_0 \le b \le \frac{3}{2}b_0$
(where $b_0$ is the impact parameter of a star whose speed at infinity is
$v_0$) attain a closest-approach distance $r_p \approx \frac{1}{2}a$ to the
CM. Also, assume that all stars exit the system with the speed found in (b).

\begin{description}
\item[(c)] (14 points) Upon each encounter, part of the total energy of the
  binary is transferred to the kinetic energy of the star. Assume that the
  binary orbit remains circular. Knowing this, using your results from
  previous tasks, and taking into account \emph{only} encounters with the
  stars within the specified range of impact parameters, show that the
  reciprocal of the binary's separation increases at a constant rate:
  \[ \frac{\mathrm{d}}{\mathrm{d}t}\left(\frac{1}{a}\right)
     = H\,\frac{G\rho}{v_0}. \]
  Here, $\rho = nm$ is the mass density of the star field, and $G$ is the
  universal gravitational constant. Find the dimensionless constant $H$,
  which refers to hardening.
\end{description}
\ioaafig{2024_TH_T8-f1.pdf}

\section*{Data-analysis problems}

\probhead{8.18}{Stellar masses in a visual binary system}{2008}{Bandung, Indonesia}{DA D2}{100}
% source id: 2008_DA_D2   tag: alpha Cen visual binary masses
The star $\alpha$-Centauri (Rigel Kentaurus) is a triple star which consists
of two main-sequence stars $\alpha$-Centauri A and $\alpha$-Centauri B
representing a visual binary system, and the third star, called Proxima
Centauri, which is smaller and fainter than the other two stars. The angular
distance between $\alpha$-Centauri A and $\alpha$-Centauri B is $17.59''$.
The binary system has an orbital period of 79.24 years. The visual magnitudes
of $\alpha$-Centauri A and $\alpha$-Centauri B are $-0.01$ and 1.34
respectively. Their color indices are 0.65 and 0.85 respectively. Use the
data below to answer the following questions.

\begin{center}
Data for main-sequence stars

\medskip
\begin{tabular}{ccc}
$(B-V)_0$ & $T_{\mathrm{eff}}$ & BC \\
$-0.25$ & 24500 & 2.30 \\
$-0.23$ & 21000 & 2.15 \\
$-0.20$ & 17700 & 1.80 \\
$-0.15$ & 14000 & 1.20 \\
$-0.10$ & 11800 & 0.61 \\
$-0.05$ & 10500 & 0.33 \\
0.00 & 9480 & 0.15 \\
0.10 & 8530 & 0.04 \\
0.20 & 7910 & 0 \\
0.30 & 7450 & 0 \\
0.40 & 6800 & 0 \\
0.50 & 6310 & 0.03 \\
0.60 & 5910 & 0.07 \\
0.70 & 5540 & 0.12 \\
0.80 & 5330 & 0.19 \\
0.90 & 5090 & 0.28 \\
1.00 & 4840 & 0.40 \\
1.20 & 4350 & 0.75 \\
\end{tabular}

\medskip
BC = Bolometric Correction, $(B-V)_0$ = Intrinsic Color
\end{center}

Questions:
\begin{enumerate}
\item Plot the curve BC versus $(B-V)_0$.
\item Determine the apparent bolometric magnitudes of $\alpha$-Centauri A
  and $\alpha$-Centauri B using the corresponding curve.
\item Calculate the mass of each star.
\end{enumerate}

Notes:
\begin{enumerate}
\item Bolometric correction (BC) is a correction that must be made to the
  apparent magnitude of an object in order to convert an object's visible
  magnitude to its bolometric magnitude:
  \[ \mathrm{BC} = m_v - m_{\mathrm{bol}} \quad \mbox{or} \quad
     \mathrm{BC} = M_v - M_{\mathrm{bol}} \]
\item Luminosity mass relation:
  \[ M_{\mathrm{bol}} = -10.2\log\left(\frac{M}{M_\odot}\right) + 4.9 \]
\end{enumerate}
\ioaafig[0.45]{2008_DA_D2-f1.pdf}

\probhead{8.19}{Light curves of stars}{2010}{Beijing, China}{DA D2}{35}
% source id: 2010_DA_D2   tag: KZ Hya pulsation period, amplitude
A pulsating variable star KZ Hydrae was observed with a telescope equipped
with a CCD camera. Figure 1 shows a CCD image of KZ Hya marked together with
the comparison star and the check star. Table 1 lists the observation time in
Heliocentric Julian dates, the magnitude differences of KZ Hya and the check
star relative to the comparison star in $V$ and $R$ band.

The questions are:
\begin{enumerate}
\item Draw the light curves of KZ Hya relative to the comparison star in $V$
  and $R$ band, respectively.
\item What are the average magnitude differences of KZ Hya relative to the
  comparison star in $V$ and $R$, respectively?
\item What are the photometry precisions in $V$ and $R$, respectively?
\item Estimate the pulsation periods of KZ Hya in $V$ and $R$.
\item Give the estimation of the pulsation amplitudes of KZ Hya in $V$ and
  $R$.
\item What is the phase delay between the $V$ and $R$ bands, in terms of the
  pulsation period?
\end{enumerate}

Table 1: Data for the light curves of KZ Hya in $V$ and $R$. $\Delta V$ and
$\Delta R$ are KZ Hya relative to the comparison in $V$ and $R$.
$\Delta V_{\mathrm{chk}}$ and $\Delta R_{\mathrm{chk}}$ are the check star
relative to the comparison in $V$ and $R$.

\begin{center}
\begin{tabular}{cccccc}
HJD$-2453800$($t$) & $\Delta V$(mag) & $\Delta V_{\mathrm{chk}}$ &
HJD$-2453800$($t$) & $\Delta R$(mag) & $\Delta R_{\mathrm{chk}}$ \\
3.162 & 0.068 & 4.434 & 3.1679 & 0.260 & 2.789 \\
3.1643 & 0.029 & 4.445 & 3.1702 & 0.185 & 2.802 \\
3.1667 & $-0.011$ & 4.287 & 3.1725 & $-0.010$ & 2.789 \\
3.1691 & $-0.100$ & 4.437 & 3.1749 & $-0.147$ & 2.809 \\
3.1714 & $-0.310$ & 4.468 & 3.1772 & $-0.152$ & 2.809 \\
\end{tabular}
\end{center}
% TABLE ABRIDGED — full data (36 rows per band) in crop edition
\ioaafig{2010_DA_D2-f1.pdf}

\probhead{8.20}{Analysis of times of minima}{2011}{Katowice, Poland}{DA D1}{25}
% source id: 2011_DA_D1   tag: V1107 Cas O-C diagram, period
Figure 1 shows the light curve of the eclipsing binary V1107 Cas, classified
as a W Ursae Majoris type. Table 1 contains a list of observed minima of the
light variation. The columns contain: the number of the minimum, the date on
which the minimum was observed, the heliocentric time of minimum expressed in
Julian days and an error (in fractions of a day).

Using these data:
\begin{description}
\item[(a)] Determine an initial period of V1107 Cas, assuming that the
  period of the star is constant during the interval of observations. Assume
  that observations during one night are continuous. Duration of the transit
  is negligible.
\item[(b)] Make what is known as an (O$-$C) diagram (for ``observed $-$
  calculated'') of the times of minima, as follows: on the $x$-axis put the
  number of periods elapsed (the ``epoch'') since a chosen initial moment
  $M_o$; on the $y$-axis the difference between the observed moment of
  minimum $M_{\mathrm{obs}}$ and the moment of minimum calculated using the
  formula (``ephemeris''):
  \[ M_{\mathrm{calc}} = M_o + P \times E \]
  where $E$, the epoch, is exactly an integer or half-integer, and $P$ is
  the period in days.
\item[(c)] Using this (O$-$C) diagram, improve the determination of the
  initial moment $M_o$ and the period $P$, and estimate the errors in their
  values.
\item[(d)] Calculate the predicted times of minima of V1107 Cas given in
  heliocentric JD occurring between 19h, 1 September 2011 UT and 02h,
  2 September 2011 UT.
\end{description}

\begin{center}
\begin{tabular}{cccc}
No. & Date of minimum (UT) & Time of minimum (Heliocentric JD) & Error \\
1 & 22 December 2006 & 2\,454\,092.4111 & 0.0004 \\
2 & 23 December 2006 & 2\,454\,092.5478 & 0.0002 \\
3 & 23 September 2007 & 2\,454\,367.3284 & 0.0005 \\
4 & 23 September 2007 & 2\,454\,367.4656 & 0.0005 \\
5 & 15 October 2007 & 2\,454\,388.5175 & 0.0009 \\
6 & 15 October 2007 & 2\,454\,388.6539 & 0.0011 \\
7 & 26 August 2008 & 2\,454\,704.8561 & 0.0002 \\
8 & 5 November 2008 & 2\,454\,776.4901 & 0.0007 \\
9 & 3 January 2009 & 2\,454\,835.2734 & 0.0007 \\
10 & 15 January 2009 & 2\,454\,847.3039 & 0.0004 \\
11 & 15 January 2009 & 2\,454\,847.4412 & 0.0001 \\
12 & 16 January 2009 & 2\,454\,847.5771 & 0.0004 \\
\end{tabular}

\medskip
Table 1: Observed times of minima of V1107 Cassiopeae
\end{center}
\ioaafig{2011_DA_D1-f1.pdf}

\probhead{8.21}{Cepheids}{2012}{Rio de Janeiro, Brazil}{DA D2}{}
% source id: 2012_DA_D2   tag: Cepheid P-L fit, two distances
Cepheids are very bright variable stars whose mean absolute magnitudes are
functions of their pulsation periods. This allows astrophysicists to easily
determine their intrinsic luminosities from the variation in their observed,
apparent luminosities.

Below is a table with Cepheid data. $P_0$ is the pulsation period in days and
$\langle M_V\rangle$ is the mean absolute visual magnitude.

% The sixth Cepheid's name is printed with a corrupted glyph in the official
% paper (it appears as a CJK character); it is rendered here as alpha UMi.
\begin{center}
\begin{tabular}{ccc}
Cepheid & $P_0$ (days) & $\langle M_V\rangle$ \\
SU Cas & 1.95 & $-1.99$ \\
V1726 Cyg & 4.24 & $-3.04$ \\
SZ Tau & 4.48 & $-3.09$ \\
CV Mon & 5.38 & $-3.37$ \\
QZ Nor & 5.46 & $-3.32$ \\
$\alpha$ UMi & 5.75 & $-3.42$ \\
V367 Sct & 6.30 & $-3.58$ \\
U Sgr & 6.75 & $-3.64$ \\
DL Cas & 8.00 & $-3.80$ \\
S Nor & 9.75 & $-3.95$ \\
$\zeta$ Gem & 10.14 & $-4.10$ \\
X Cyg & 16.41 & $-4.69$ \\
WZ Sgr & 21.83 & $-5.06$ \\
SW Vel & 23.44 & $-5.09$ \\
SV Vul & 44.98 & $-6.04$ \\
\end{tabular}
\end{center}

\begin{enumerate}
\item Plot all Cepheids in a scatter diagram. $\log_{10}(P_0)$ should be the
  abscissa and $\langle M_V\rangle$ should be the ordinate.
\item Fit, using least squares, a straight line to the $\langle M_V\rangle$
  vs $\log_{10}(P_0)$ plot. This equation allows one to obtain the absolute
  magnitude from the pulsation period for any Cepheid.
\item Figures 1 and 2 show the light curves of two Cepheids. Use the
  available data to estimate the distances to each of these two Cepheids.
  Also estimate the uncertainty of the distance determination (exact
  formulae are not necessary).
\item Comparing the difference between the distances of the two stars with
  the typical size of a galaxy, would it be possible for these two stars to
  be in the same galaxy? Please mark ``YES'' (~) or ``NO'' (~).
\end{enumerate}

\paragraph{Appendix A: extinction coefficients.}
The magnitude outside the atmosphere is given by:
\[ M = m - A\cdot X + B \]
where $A$ is the extinction coefficient and $B$ is the zero point coefficient
for the night, $X$ is the airmass of the observation and $m$ is the
instrumental magnitude (that is, the magnitude obtained directly from the
image).

\paragraph{Appendix B: calibrated differential magnitude.}
If the standard star is in the same frame as the object, the calibrated
magnitude of the object can be obtained by:
\[ M_{\mathrm{ast}} = m_{\mathrm{ast}} - m_{\mathrm{star}}
   + M_{\mathrm{star}} \]

\paragraph{Appendix C: least squares estimation of angular coefficients.}
Given straight lines described by
\[ y_i = \alpha + \beta x_i,\ i = 1..n \]
the least squares estimator of $\beta$ is given by
\[ \beta = \frac{\sum_i^n x_i y_i - \frac{1}{n}\sum_i^n x_i \sum_i^n y_i}
   {\sum_i^n x_i^2 - \frac{1}{n}\left(\sum_i^n x_i\right)^2} \]
\ioaafig{2012_DA_D2-f1.pdf}
\ioaafig{2012_DA_D2-f2.pdf}
\ioaafig{2012_DA_D2-f3.pdf}

\probhead{8.22}{Cepheid HV2257}{2015}{Magelang, Indonesia}{DA D1}{}
% source id: 2015_DA_D1   tag: Cepheid curves, Baade-Wesselink distance
Photometry and radial velocity data for the Cepheid type star HV2257 are
given in Tables 1--3, based on observations by Gieren (MNRAS vol 265, 1993).
The pulsation period of the star is $P = 39.294$ days. A reference graph for
the temperature--color relation and the bolometric correction tables are
given in Figure 1 (Houdashelt et al., 2000) and Table 4
(http://xoomer.virgilio.it/hrtrace/Straizys.htm). Given that the solar
luminosity is $L_\odot = 3.96 \times 10^{26}\,\mathrm{J\,s^{-1}}$ and its
bolometric magnitude $M_{\odot\,\mathrm{bol}} = 4.72$. Please do not use the
period--luminosity relation from the second question for this question.

\begin{description}
\item[(a)] Plot the light curve based on Table 1, between phases 0.6 and 1.
\item[(b)] Plot the color in Table 2, between phases 0.6 and 1.
\item[(c)] Plot the radial velocity curve from Table 3, between phases 0.6
  and 1.
\item[(d)] Calculate the average radial velocity of the star.
\item[(e)] Calculate the distance to this pulsating star using the observed
  data and supplementary data given in Table 4 and Figure 1. Assume that
  there is no extinction in this direction.
\end{description}

\begin{center}
\begin{tabular}{cc@{\qquad}cc@{\qquad}cc}
\multicolumn{2}{c}{Table 1} & \multicolumn{2}{c}{Table 2}
  & \multicolumn{2}{c}{Table 3} \\
Phase & $V$ mag & Phase & $V-R$ & Phase & RadVel (km/s) \\
0.11 & 12.81 & 0.22 & 0.71 & 0.03 & 232 \\
0.13 & 12.84 & 0.24 & 0.73 & 0.05 & 234 \\
0.14 & 12.87 & 0.25 & 0.74 & 0.08 & 234 \\
0.16 & 12.88 & 0.27 & 0.75 & 0.08 & 237 \\
0.19 & 12.90 & 0.29 & 0.75 & 0.13 & 242 \\
\end{tabular}
\end{center}
% TABLE ABRIDGED — Tables 1 (36 rows), 2 (36 rows) and 3 (46 rows) are given
% in full in the crop edition

\begin{center}
Table 4. Bolometric correction

\medskip
\begin{tabular}{cc@{\qquad\qquad}cc}
$T_{\mathrm{eff}}$, K & BC, mag & $T_{\mathrm{eff}}$, K & BC, mag \\
9600 & $-0.25$ & 5250 & $-0.22$ \\
9400 & $-0.16$ & 5050 & $-0.29$ \\
9150 & $-0.10$ & 4950 & $-0.35$ \\
8900 & $-0.03$ & 4850 & $-0.42$ \\
8400 & 0.05 & 4700 & $-0.57$ \\
8000 & 0.09 & 4600 & $-0.75$ \\
7300 & 0.13 & 4400 & $-1.17$ \\
7100 & 0.11 & 3900 & $-1.25$ \\
6500 & 0.08 & 3750 & $-1.40$ \\
6150 & 0.03 & 3550 & $-1.60$ \\
5950 & 0.00 & 3400 & $-2.00$ \\
5800 & $-0.05$ & & \\
5500 & $-0.13$ & & \\
\end{tabular}
\end{center}
\ioaafig{2015_DA_D1-f1.pdf}
\ioaafig[0.52]{2015_DA_D1-f2.pdf}
\ioaafig[0.45]{2015_DA_D1-f3.pdf}

\probhead{8.23}{Binary Pulsar}{2016}{Bhubaneswar, India}{DA D1}{50}
% source id: 2016_DA_D1   tag: binary pulsar orbit determination
Through systematic searches during the past decades, astronomers have found a
large number of millisecond pulsars (spin period $< 10$\,ms). The majority of
these pulsars are found in binaries, with nearly circular orbits.

For a pulsar in a binary orbit, the measured pulsar spin period ($P$) and the
measured line-of-sight acceleration ($a$) both vary systematically due to
orbital motion. For circular orbits, this variation can be described
mathematically in terms of orbital phase $\phi$ ($0 \le \phi \le 2\pi$) as,
\[ P(\phi) = P_0 + P_{\mathrm{t}}\cos\phi \qquad \mbox{where }
   P_{\mathrm{t}} = \frac{2\pi P_0 r}{c P_{\mathrm{B}}}, \]
\[ a(\phi) = -a_{\mathrm{t}}\sin\phi \qquad \mbox{where }
   a_{\mathrm{t}} = \frac{4\pi^2 r}{P_{\mathrm{B}}^2}, \]
where $P_{\mathrm{B}}$ is the orbital period of the binary, $P_0$ is the
intrinsic spin period of the pulsar and $r$ is the radius of the orbit.

The following table gives one such set of measurements of $P$ and $a$ at
different heliocentric epochs, $T$, expressed in truncated Modified Julian
Days (tMJD), i.e.\ number of days since $\mathrm{MJD} = 2\,440\,000$.

\begin{center}
\begin{tabular}{cccc}
No. & $T$ & $P$ & $a$ \\
 & (tMJD) & ($\mu$s) & (m\,s$^{-2}$) \\
1 & 5740.654 & 7587.8889 & $-0.92 \pm 0.08$ \\
2 & 5740.703 & 7587.8334 & $-0.24 \pm 0.08$ \\
3 & 5746.100 & 7588.4100 & $-1.68 \pm 0.04$ \\
4 & 5746.675 & 7588.5810 & $+1.67 \pm 0.06$ \\
5 & 5981.811 & 7587.8836 & $+0.72 \pm 0.06$ \\
6 & 5983.932 & 7587.8552 & $-0.44 \pm 0.08$ \\
7 & 6005.893 & 7589.1029 & $+0.52 \pm 0.08$ \\
8 & 6040.857 & 7589.1350 & $+0.00 \pm 0.04$ \\
9 & 6335.904 & 7589.1358 & $+0.00 \pm 0.02$ \\
\end{tabular}
\end{center}

By plotting $a(\phi)$ as a function of $P(\phi)$, we can obtain a parametric
curve. As evident from the relations above, this curve in the
period-acceleration plane is an ellipse.

In this problem, we estimate the intrinsic spin period, $P_0$, the orbital
period, $P_{\mathrm{B}}$, and the orbital radius, $r$, by an analysis of this
data set, assuming a circular orbit.

\begin{description}
\item[(D1.1)] Plot the data, including error bars, in the
  period-acceleration plane (mark your graph as ``D1.1'').
\item[(D1.2)] Draw an ellipse that appears to be a best fit to the data (on
  the same graph ``D1.1'').
\item[(D1.3)] From the plot, estimate $P_0$, $P_{\mathrm{t}}$ and
  $a_{\mathrm{t}}$, including error margins.
\item[(D1.4)] Write expressions for $P_{\mathrm{B}}$ and $r$ in terms of
  $P_0$, $P_{\mathrm{t}}$, $a_{\mathrm{t}}$.
\item[(D1.5)] Calculate approximate value of $P_{\mathrm{B}}$ and $r$ based
  on your estimations made in (D1.3), including error margins.
\item[(D1.6)] Calculate orbital phase, $\phi$, corresponding to the epochs
  of the following five observations in the above table: data rows 1, 4, 6,
  8, 9.
\item[(D1.7)] Refine the estimate of the orbital period, $P_{\mathrm{B}}$,
  using the results in part (D1.6) in the following way:
  \begin{description}
  \item[(D1.7a)] First determine the initial epoch, $T_0$, which corresponds
    to the nearest epoch of zero orbital phase before the first observation.
  \item[(D1.7b)] The expected time, $T_{\mathrm{calc}}$, of the estimated
    orbital phase angle of each observation is given by,
    \[ T_{\mathrm{calc}} = T_0
       + \left(n + \frac{\phi}{360^\circ}\right) P_{\mathrm{B}}, \]
    where $n$ is the number of full cycles of orbital phases that may have
    elapsed between $T_0$ and $T$ (or $T_{\mathrm{calc}}$). Estimate $n$ and
    $T_{\mathrm{calc}}$ for each of the five observations in part (D1.6).
    Note down the difference $T_{\mathrm{O-C}}$ between observed $T$ and
    $T_{\mathrm{calc}}$. Enter these calculations in the table given in the
    Summary Answersheet.
  \item[(D1.7c)] Plot $T_{\mathrm{O-C}}$ against $n$ (mark your graph as
    ``D1.7'').
  \item[(D1.7d)] Determine the refined values of the initial epoch,
    $T_{0,\mathrm{r}}$, and the orbital period, $P_{\mathrm{B,r}}$.
  \end{description}
\end{description}

\probhead{8.24}{Compact Object in a Binary System}{2018}{Beijing, China}{DA D2}{75}
% source id: 2018_DA_D2   tag: RV curve, mass function, compact companion
Astronomers discovered an extraordinary binary system in the constellation of
Auriga during the course of the Apache Point Observatory Galactic Evolution
Experiment (APOGEE). In these questions, you will try to analyse the data and
recreate their discovery for yourself.

The research team is aiming to find compact stars in binary systems using the
radial velocity (RV) technique. They examined archival APOGEE spectra of
``single'' stars and measured the apparent variation of their RV within this
data. Among ${\sim}200$ stars with the highest accelerations, researchers
searched for periodic photometric variations in data from the All-Sky
Automated Survey for Supernovae (ASAS-SN) that might be indicative of
transits, ellipsoidal variations or starspots. After this process, they
spotted a star named 2M05215658+4359220, with a large variation in RV and
photometric variability.

2.1. The following table presents the radial velocity measurements of
2M05215658+4359220 during three epochs of APOGEE spectroscopic observation.
Here we assume the variation of its RV is due to the existence of an unseen
companion. The proper motion of the stars can be ignored.

\begin{center}
Table 3. APOGEE Radial Velocity Measurements of 2M05215658+4359220

\medskip
\begin{tabular}{cccc}
Observation & MJD & RV (km\,s$^{-1}$) & Uncertainty (km\,s$^{-1}$) \\
1 & 56204.9537 & $-37.417$ & 0.011 \\
2 & 56229.9213 & 34.846 & 0.010 \\
3 & 56233.8732 & 42.567 & 0.010 \\
\end{tabular}
\end{center}

\begin{description}
\item[(D2.1.1)] (6 points) Use the data and a simple linear model to obtain
  an initial estimate of the apparent \textbf{maximum acceleration} of the
  star:
  \[ a_{\max} = \left.\frac{\Delta RV}{\Delta t}\right|_{\max}, \quad
     \mbox{unit: } \mathrm{km\,s^{-1}\,day^{-1}} \]
\item[(D2.1.2)] (9 points) Now use the data to obtain an initial estimate of
  the \textbf{mass} of its unseen companion.
\end{description}

2.2. After discovering this peculiar star, astronomers conducted follow-up
observations using the 1.5-m Tillinghast Reflector Echelle Spectrograph
(TRES) at the Fred Lawrence Whipple Observatory (FLWO) located on Mt.\
Hopkins in Arizona, USA. The following table presents the RV measurements
using this instrument:

\begin{center}
Table 4. TRES Radial Velocity Measurements of 2M05215658+4359220

\medskip
\begin{tabular}{ccc}
MJD & RV (km/s) & Uncertainty (km/s) \\
58006.9760 & 0 & 0.075 \\
58023.9823 & $-43.313$ & 0.075 \\
58039.9004 & $-27.963$ & 0.045 \\
58051.9851 & 10.928 & 0.118 \\
58070.9964 & 43.782 & 0.075 \\
58099.8073 & $-30.033$ & 0.054 \\
58106.9178 & $-42.872$ & 0.135 \\
58112.8188 & $-44.863$ & 0.088 \\
58123.7971 & $-25.810$ & 0.115 \\
58136.6004 & 15.691 & 0.146 \\
58143.7844 & 34.281 & 0.087 \\
\end{tabular}
\end{center}

\begin{description}
\item[(D2.2.1)] (14 points) \textbf{Plot} the diagram of RV variation
  (measured with TRES) versus time on your graph paper and label it as
  Figure 4. Draw a suitable \textbf{sinusoidal function} to fit the given
  data. \textbf{Estimate} the orbital period ($P_{\mathrm{orb}}$) and radial
  velocity semi-amplitude ($K$) from your plot.
\item[(D2.2.2)] (4 points) If the star is moving in a circular orbit,
  calculate the \textbf{minimum value} of the orbital radius
  ($r_{\mathrm{orb}}$) of the star in units of both $R_\odot$ and au.
\item[(D2.2.3)] (7 points) The mass function of a binary system is defined
  as:
  \[ f(M_1, M_2) = \frac{(M_2 \sin i_{\mathrm{orb}})^3}{(M_1 + M_2)^2} \]
  where the subscript ``1'' represents the primary star and ``2'' represents
  its companion. The parameter $i_{\mathrm{orb}}$ is the orbital inclination
  of the binary system. This mass function can also be expressed in terms of
  observable parameters. Calculate the \textbf{mass function of this system}
  in units of $M_\odot$.
\end{description}

2.3. Based on a detailed analysis of APOGEE, TRES spectra and GAIA parallax
measurements, astronomers derived the following stellar parameters:

\begin{center}
Table 5. Selected Physical Properties of 2M05215658+4359220

\medskip
\begin{tabular}{ccccc}
Effective & Surface Gravity & Parallax & Measured Rotation & Bolometric \\
Temperature & $\log g$ (cm\,s$^{-2}$) & $\pi$ (mas) & Velocity & Flux \\
$T_{\mathrm{eff}}$ (K) & & &
  $v_{\mathrm{rot}}\sin i$ (km\,s$^{-1}$) &
  $F$ (J\,s$^{-1}$\,m$^{-2}$) \\
$4890 \pm 130$ & $2.2 \pm 0.1$ & $0.272 \pm 0.049$ & $14.1 \pm 0.6$ &
  $(1.1 \pm 0.1) \times 10^{-12}$ \\
\end{tabular}
\end{center}

Photometric observations indicate that the period of its light curve is
identical to its orbital period, thus we may assume that the rotation period
satisfies $P_{\mathrm{rot}} = P_{\mathrm{orb}} \equiv P$, and the inclination
satisfies $i_{\mathrm{orb}} = i_{\mathrm{rot}} \equiv i$.

\begin{description}
\item[(D2.3.1)] (16 points) \textbf{Calculate} the luminosity ($L_1$, in
  unit of $L_\odot$), radius ($R_1$, in unit of $R_\odot$), sine of the
  inclination angle ($\sin i$), as well as mass ($M_1$, in unit of
  $M_\odot$) of the visible star. Please \textbf{include} the uncertainty in
  your results.
\item[(D2.3.2)] (4 points) \textbf{Choose} the correct type of this star
  from the following options: (1) Blue Giant (2) Yellow main sequence star
  (3) Red Giant (4) Red main sequence star (5) White Dwarf.
\item[(D2.3.3)] (10 points) Based on the mass function $f(M_1, M_2)$ of the
  binary system, \textbf{plot} the rough relationship of $M_2$ (as vertical
  axis) and $M_1$ (as horizontal axis) on your graph paper and label it as
  Figure 5. Plot the most probable relation (by using $\sin i$), upper limit
  (with $\sin i + \Delta\sin i$) and lower limit (with
  $\sin i - \Delta\sin i$) derived in (D2.3.1).
\item[(D2.3.4)] (5 points) Draw a vertical \textbf{shadowed region} of
  $[M_1 - \Delta M_1, M_1 + \Delta M_1]$, as well as two horizontal
  \textbf{dashed lines} showing the maximum mass of the white dwarf and
  neutron star, on your Figure 5. What is the possible mass of the invisible
  companion, and what kind of celestial object could it be?
\end{description}

\probhead{8.25}{Photometry and spectroscopy of Nova Del 2013}{2019}{Keszthely, Hungary}{DA 1}{60}
% source id: 2019_DA_1   tag: nova decline, MMRD distance, P Cygni
Classical nova V339 Del (Nova Delphini 2013) was discovered by Koichi Itagaki
at 6.8 magnitude on 14 August 2013 at 14:01 UT ($\mathrm{MJD} = 56518.584$).
Both professional and amateur astronomers analysed the photometry and
spectroscopy of the nova. Less than 10 hours after the alert, when the night
falls at the Piszk\'estet\H{o} Mountain Station of the Konkoly Observatory of
the Hungarian Academy of Sciences, Hungarian astronomers took the first
spectrum data of the nova using the eShel echelle spectrograph in the Gothard
Astrophysical Observatory of Lor\'and E\"otv\"os University attached to the
1 meter telescope of the Konkoly Observatory.

Refer to Fig.~1.1 and Fig.~1.2 to complete the questions. The larger versions
of Fig.~1.1 and Fig.~1.2 are found on separate A3 papers.

Fig.~1.1 shows the visual light curve of the nova based on the data
downloaded from the website of AAVSO (American Association of Variable Star
Observers). On the horizontal and vertical axes, the Modified Julian Date
($\mathrm{MJD} = \mathrm{JD} - 2\,400\,000.5$) of the observations and the
Johnson $V$ magnitudes are plotted, respectively. The grey circles (about
38000 data points) represent the measured values, while the continuous black
line is the result of smoothing the data with a Gaussian filter (Full Width
at Half Maximum $= 0.5$\,day) to define an ``average'' light curve from the
data points.

The rate of decline can be characterized by the values $t_2$ and $t_3$, which
show the time interval in days in which a nova fades from its maximum
brightness by 2 and 3 magnitudes.

A few empirical formulae between the peak of the absolute magnitude in the
$V$ band ($M_0$) and $t_2$, $t_3$ values can be found in the following
literature:
\begin{description}
\item[(a)] $M_0 = -7.92 - 0.81\arctan\dfrac{1.32 - \log t_2}{0.23}$ \\
  (Della Valle, M. \& Livio, M.: 1995, ApJ 452, 704)
\item[(b)] $M_0 = -11.32 + 2.55\log t_2$ \\
  (Downes, R.A. \& Durbeck, H.W.: 2000, AJ 120, 2007)
\item[(c)] $M_0 = -11.99 + 2.54\log t_3$ \\
  (Downes, R.A. \& Durbeck, H.W.: 2000, AJ 120, 2007)
\end{description}

The $E(B-V)$ color excess of Nova Del 2013 (Chochol, D. et al.: 2014,
Contrib. Astron. Obs. Skalnat\'e Pleso 43, 330) is:
\[ E(B-V) = 0.184 \pm 0.035 \]

Fig.~1.2 shows the nova spectra taken in the wavelength region around the
H$\alpha$ line on six consecutive nights before and after the time of the
maximum brightness ($t_0$). The individual spectra have been shifted
vertically for clarity. The Modified Julian Dates (MJD) of the observations
are listed on the right hand side of each spectrum slice.

The H$\alpha$ line shows the so-called P Cygni profile with very broad wings,
which is typical not of novae only, but is present in almost all spectral
types and is a reliable sign of a massive radial motion of matter ejected
from the star. The P Cygni profile is composed of a strong, broad emission
peak which is considered to be centered at the rest wavelength in air
$\lambda_0$ of the line -- for H$\alpha$ $\lambda_0 = 6562.82$\,\AA{} -- and
a usually weaker, blueshifted absorption component. The expansion (radial)
velocity of the shell can be approximated from the measured wavelength
$\lambda$ of the absorption peak using the well known Doppler formula
connecting the displacement $\Delta\lambda = \lambda - \lambda_0$, the
radial velocity $v_{\mathrm{r}}$, and $c$, the speed of light.

Assume that the H$\alpha$ line showing P Cygni profile is excited in the
outermost part of the spherically expanding shell, and its extent at the
moment of taking the first spectrum was still negligible.

\begin{description}
\item[(a)] From Fig.~1.1, estimate the Modified Julian Date of the peak
  magnitude ($\mathrm{MJD}_0$) and the value of the peak magnitude itself.
  Consider the error of this brightness value to be $0.05^{\mathrm{m}}$.
  \hfill[3\,p]
\item[(b)] Estimate the Modified Julian Dates based on the time interval
  (days) in which the nova has faded by 2 and 3 magnitudes, then calculate
  $t_2$ and $t_3$ values. \hfill[6\,p]
\item[(c)] With reference to $t_2$ and $t_3$ from (b), determine the peak
  absolute magnitude of the nova using all three empirical formulae listed
  earlier, calculate their mean ($M_0$) and their standard deviation, and
  consider this latter as the uncertainty of $M_0$. \hfill[5\,p] \\
  The formula for standard deviation:
  \[ \sigma = \sqrt{\frac{\sum_{i=1}^{n}(x_i - \bar{x})^2}{n-1}} \]
\item[(d)] Using the value of the color excess $E(B-V)$, determine the
  interstellar extinction $A_V$ and its uncertainty in the direction of the
  nova. Use $R_V = 3.1$, without error. \hfill[4\,p]
\item[(e)] Estimate the distance to the nova and its uncertainty. Give the
  result in kpc. \hfill[11\,p]
\item[(f)] Measure the central wavelengths of the P Cygni absorption
  features plotted in Fig.~1.2 (refer to magnified version), and calculate
  the corresponding radial velocities. No error estimation is needed.
  \hfill[14\,p]
\item[(g)] Plot these radial velocities against the Modified Julian Dates of
  the observations. \hfill[6\,p]
\item[(h)] From the graph in (g), estimate the physical radius of the
  envelope at the end of the time interval. Give the answer in astronomical
  units (au). \hfill[7\,p]
\item[(i)] Knowing the distance to the nova and the physical radius of the
  spherical envelope 5 days after the discovery, estimate the apparent
  angular diameter of the envelope then. \hfill[4\,p]
\end{description}
\ioaafig{2019_DA_1-f1.pdf}

\probhead{8.26}{Triply eclipsing hierarchical triple stellar system}{2019}{Keszthely, Hungary}{DA 2}{90}
% source id: 2019_DA_2   tag: triple-star eclipses, RV+ETV masses
HD 181068 was one of the brightest targets which was continuously observed
during the almost 4-year-long primary mission of NASA's exoplanet-hunter
Kepler space telescope. The spacecraft observed ${\approx}3-4 \times 10^{-3}$
magnitude dimmings every 0.453 days. (Note: the even dimmings were slightly
smaller amplitude than the odd ones.) Furthermore, additional 0.007
magnitude, 2.3-day-long dimmings were detected every 22.7 days.

The correct explanation of this very unusual photometric behaviour was given
by Hungarian astronomers. They found that HD 181068 is a compact hierarchical
triple stellar system seen almost edge-on.

Hierarchical triple star systems consist of three stars; A, B, and C. Two of
these stars (B and C) form an inner or close stellar binary system, whilst
the outer component (star A) orbits at a distance from the inner system
significantly larger (usually orders of magnitude) than the semi-major axis
of the inner system. The schematic view of an example of a triple star system
is illustrated in Fig.~2.1 (the schematic pole-on view of a hypothetical
hierarchical triple stellar system; the black arrow is directed towards the
Earth; the thick segments of the three orbits represent the stars' orbital
arcs during an outer eclipse).

Mathematically, the motion of a hierarchical triple system can be well
approximated with two unperturbed Keplerian two-body motions; (1) Keplerian
motion of the inner binary. (2) The centre of mass of this close binary and
the third star revolves on a second Keplerian orbit, ``outer binary''.

In this problem, stars B and C form a $P_1 = 0.9056768$-day-period eclipsing
binary, while the centre of mass of these stars with star A forms the
$P_2 = 45.4711$-day-period outer binary. As the orbital plane of this outer
orbit is seen almost edge-on from the Kepler spacecraft (and from the Earth),
during their revolution on the outer orbit, stars B and C eclipse not only
each other, but also star A or, a half outer revolution later are eclipsed by
it, causing the extra dimmings.

\begin{description}
\item[i.] \emph{Determination of the physical stellar sizes (and other
  quantities) from the geometry of the eclipses}

  These assumptions are used throughout this section: (1) both the inner and
  outer orbits are exactly circular, (2) the orbital planes of the inner and
  outer orbits are identical, and (3) this plane is seen exactly edge-on
  (i.e.\ $i_1 = i_2 = 90^\circ$ and $i_{\mathrm{rel}} = 0^\circ$). Let's
  consider the extra dimmings, which are central eclipses (i.e.\ either
  occultations or transits -- annular eclipses), therefore, these events
  have four contacts. In the case of an ordinary eclipsing binary (or of a
  transiting exoplanet) at the times of the outer contacts the sky-projected
  disks of the two objects connect with each other at one point from
  outside, while at the inner contacts they approach each other from inside.
  While this last statement is also valid for outer eclipses, the situation
  becomes more complex, because instead of two, three stars are involved
  into the eclipses. However, despite this fact, we can certainly define the
  times of each contact from the light curve, and furthermore, we can also
  decide explicitly which member of the inner binary is involved in a given
  contact. (The other member, of course, is always star A.)

  In the table below, the accurate times of some contacts of different
  eclipses observed by the Kepler spacecraft, the contact types and the
  stars are documented. Time is expressed in barycentric Julian Days (BJD).

  \begin{center}
  \begin{tabular}{cccccc}
  event no. & contact & stars & BJD & $\varphi_1$ & $\varphi_2$ \\
  1 & I & A, B & 2455476.1096 & & \\
    & II & A, C & 2455476.4245 & & \\
    & III & A, B & 2455477.9677 & & \\
    & IV & A, B & 2455478.4722 & & \\
  2 & I & A, B & 2455521.5217 & & \\
  3 & III & A, C & 2455568.9434 & & \\
  4 & I & A, C & 2455612.4733 & & \\
    & III & A, C & 2455614.3571 & & \\
  5 & III & A, B & 2455659.9241 & & \\
    & IV & A, C & 2455660.2422 & & \\
  \end{tabular}
  \end{center}

  \begin{description}
  \item[(a)] Given that $T_{01} = 2455051.2361$ and
    $T_{02} = 2455522.7318$ denote the time of an inferior conjunction of
    the inner and outer binaries respectively (i.e.\ that time, when, from
    the perspective of the observer, star C eclipses star B, and when star A
    eclipses the centre of mass of stars B and C.)

    Define
    \[ \varphi_1(t) = \{(t - T_{01})/P_1\}, \quad \mbox{and} \quad
       \varphi_2(t) = \{(t - T_{02})/P_2\} \]
    as the photometric phases of the inner and outer binaries respectively.
    $\{x\}$ denotes the decimal part of the real number $x$. If
    $\{x\} < 0$, use $\{x\} + 1$ instead. Calculate the phases for the times
    of the tabulated contact events and write the answers in the appropriate
    columns of the table on the answer sheet. Round your answers to four
    decimal places. \hfill[10\,p]
  \item[(b)] Determine, whether star A, or the close binary (i.e.\ stars B
    and C) were closer to the observer during each eclipsing event. Write
    your answer in the table on the answer sheet. \hfill[5\,p]
  \item[(c)] Using the table above, calculate (1) the dimensionless radius
    of each star relative to the semi-major axis of the outer orbit
    ($R_{\mathrm{A,B,C}}/a_2$), (2) the ratio of the semi-major axes of the
    two orbits ($a_1/a_2$) and (3) the mass ratio of stars B and C
    ($q_1 = m_{\mathrm{C}}/m_{\mathrm{B}}$). Hint: use at least four decimal
    place accuracy in your calculations. Be cautious, it may not be possible
    to use all theoretically possible contact combinations with a given
    limited accuracy of time data. \hfill[30\,p]
  \item[(d)] Based on the results obtained above, calculate the outer mass
    ratio ($q_2 = m_{\mathrm{BC}}/m_{\mathrm{A}}$). \hfill[8\,p]
  \end{description}

\item[ii.] \emph{Dynamical determination of the stellar masses using radial
  velocity (RV) and eclipse timing variations (ETV) measurements}

  To obtain RV data, ground-based spectroscopic follow up observations were
  carried out with four different instruments. Only the lines of star A were
  detectable in all spectra. Plotting all the measurements against time, the
  RV curve was nicely fitted in the following form:
  \[ V_{\mathrm{rad,A}} = V_\gamma + K_{\mathrm{A}}\sin\phi_{\mathrm{RV}}, \]
  where $V_\gamma$ is the systemic velocity and $K_{\mathrm{A}}$ is the
  velocity amplitude:
  \[ V_\gamma = 6.993 \pm 0.011\,\mathrm{km\,s^{-1}}, \quad
     K_{\mathrm{A}} = 37.195 \pm 0.053\,\mathrm{km\,s^{-1}}, \]
  \[ P_2 = 45.4711 \pm 0.0002\,\mathrm{d}, \quad
     \phi_{\mathrm{RV}} = \frac{2\pi}{P_2}
     \left[t - (2455522.7318 \pm 0.0095)\right]. \]
  Furthermore, the researchers determined the mid-times of the regular
  eclipses of the close binary (formed by stars B and C), and found that the
  occurrence of for example the eclipsing minima belonging to the $N$th
  orbital revolution can be described by the simple expression:
  \[ T_N = T_0 + P_1 N + A_{\mathrm{ETV}}
     \sin\left(\frac{2\pi}{P_2}P_1 N + \phi_0\right), \]
  where
  \[ T_0 = \mathrm{BJD}\ 2455051.23607 \pm 5\times10^{-5}, \quad
     P_1 = 0.9056768 \pm 3\times10^{-7}\,\mathrm{d}, \]
  \[ A_{\mathrm{ETV}} = 0.001446 \pm 0.000110\,\mathrm{d}, \quad
     \phi_0 = -0.76779 \pm 0.01937\,\mathrm{rad}. \]
  In this expression $A_{\mathrm{ETV}}$ is the amplitude of the eclipse
  timing variation, $T_0$ denotes the mid-eclipse time of the reference
  (zeroth) primary eclipse, and $N$ is the cycle number, which is an integer
  for primary eclipses (i.e.\ when the slightly fainter star C eclipses star
  B), and half-integer for secondary ones (i.e.\ when star B eclipses star
  C).

  Determine (1) again the mass ratio
  ($q_2 = m_{\mathrm{BC}}/m_{\mathrm{A}}$) of the centre of mass of the
  inner binary and star A using only the results obtained in point ii., (2)
  the mass of component A ($m_{\mathrm{A}}$) and (3) the total mass of the
  inner, close binary ($m_{\mathrm{BC}}$). Calculate the errors for (1),
  (2), and (3) in masses. Hint: you can save much time by expressing the
  masses in solar mass and the orbital separations either in solar radius or
  au. \hfill[22\,p]

\item[iii.] Using results obtained in questions 1 and 2, determine the
  masses of stars B and C respectively and calculate the physical dimensions
  of all three stars (i.e.\ stellar radii in physical units). \hfill[15\,p]
\end{description}
\ioaafig{2019_DA_2-f1.pdf}
\ioaafig{2019_DA_2-f2.pdf}

\section*{Observation --- telescope}

\probhead{8.27}{Colours of the binary epsilon Trianguli Australis}{2012}{Rio de Janeiro, Brazil}{OBS-TEL 3}{}
% source id: 2012_OBS-TEL_3   tag: visual colour estimation of a binary
Point your telescope to the binary star $\varepsilon$-Trianguli Australis
using star chart-3 as a guide. That pair's components are magnitude 4.1 and
9.3, separated by $82''$. Choose the best option for the correct color of
each star:

\begin{description}
\item[Brighter:] White/blue (\quad) \qquad Yellow (\quad) \qquad
  Red (\quad)
\item[Dimmer:] White/blue (\quad) \qquad Yellow (\quad) \qquad
  Red (\quad)
\end{description}

\clearpage
\section*{Solutions to Chapter 8}

\solhead{8.1}{Binary System}
% source id: 2007_TH_Q3
\begin{description}
\item[a)] The total angular momentum of the system is
\[ L = I\omega = (M_1 r_1^2 + M_2 r_2^2)\,\omega \]
\hfill(0.5 points)

We have also $M_1 r_1 = M_2 r_2$ and $D = r_1 + r_2$, which yield
\[ L = \frac{M_1 M_2}{M_1 + M_2} D^2 \omega \qquad (1) \]
\hfill(0.5 points)

The kinetic energy of the system is
\begin{align*}
K.E. &= \frac{1}{2} M_1 (r_1\omega)^2 + \frac{1}{2} M_2 (r_2\omega)^2
  = \frac{1}{2}\left(M_1 r_1^2 + M_2 r_2^2\right)\omega^2
  && \text{(0.5 points)} \\
&= \frac{1}{2}\,\frac{M_1 M_2}{(M_1 + M_2)}\, D^2\omega^2
  \qquad (2)
  && \text{(0.5 points)}
\end{align*}

\item[b)] From Newton's laws of motion we have
\[ M_1\omega^2 r_1 = M_2\omega^2 r_2 = \frac{G M_1 M_2}{D^2} \]
\hfill(1 point)

These equations together with those in a) yield
\[ \omega^2 = \frac{G(M_1 + M_2)}{D^3} \qquad (3) \]
\hfill(1 point)

\item[c)] In order to find the quantity $\Delta\omega$, we must also realize
that
\[ M_1 + M_2 = \text{constant} \qquad (4) \]
\hfill(0.5 points)

And that, since there is no external torque acting on the system, the total
angular momentum must be conserved.
\[ L = \frac{M_1 M_2}{M_1 + M_2} D^2\omega = \text{constant} \]
that is,
\[ M_1 M_2 D^2 \omega = \text{constant} \]
\hfill(0.5 points)

Now, after the mass transfer,
\begin{align*}
\omega &\to \omega + \Delta\omega \\
M_1 &\to M_1 + \Delta M_1 \\
M_2 &\to M_2 - \Delta M_1 \\
D &\to D + \Delta D
\end{align*}
Hence,
\[ M_1 M_2 D^2\omega
   = \left(M_1 + \Delta M_1\right)\left(M_2 - \Delta M_1\right)
     \left(D + \Delta D\right)^2\left(\omega + \Delta\omega\right) \]

After using the approximation $(1+x)^n \sim 1 + nx$ and rearranging, we get
\[ \frac{\Delta\omega}{\omega} + 2\,\frac{\Delta D}{D}
   = \left(\frac{M_1 - M_2}{M_1 M_2}\right)\Delta M_1 \qquad (5) \]
\hfill(1 point)

From equation (3), $\omega^2 D^3$ is also constant. That is,
\[ \omega^2 D^3 = (\omega + \Delta\omega)^2 (D + \Delta D)^3 \]
This gives,
\[ \frac{\Delta D}{D} = -\frac{2}{3}\,\frac{\Delta\omega}{\omega}
   \qquad (6) \]
\hfill(0.5 points)

Hence
\[ \Delta\omega = -\frac{3\left(M_1 - M_2\right)}{M_1 M_2}\,
   \omega\,\Delta M_1 \qquad (7) \]
\hfill(0.5 points)

\item[d)] Given that
\begin{gather*}
M_1 = 2.9\,\mathrm{M}_\odot\,, \quad M_2 = 1.4\,\mathrm{M}_\odot \\
\text{orbital period,}\quad T = 2.49\ \text{days}
\end{gather*}
and that $T$ has increased by 20\,s in the past 100 years, we have
\[ \omega = \frac{2\pi}{T} \quad\text{and}\quad
   \Delta\omega = -\frac{\omega}{T}\,\Delta T \qquad (8) \]
\hfill(0.5 points)

\[ \Delta M_1 = +\frac{1}{3}\left(\frac{M_1 M_2}{M_1 - M_2}\right)
   \frac{\Delta T}{T} \]
\hfill(0.5 points)
\begin{align*}
\frac{\Delta M_1}{M_1 \Delta t}
  &= \frac{1}{3}\left(\frac{1.4}{(2.9 - 1.4)}\right)
     \left(\frac{20}{2.49 \times 24 \times 3600}\right)
     \left(\frac{1}{100}\right) \\
  &= 2.89 \times 10^{-7}\ \text{per year}
  && \text{(0.5 points)}
\end{align*}

\item[e)] Mass is flowing from $M_2$ to $M_1$. \hfill(0.5 points)

\item[f)] From equations (6) and (8):
\[ \frac{\Delta D}{D\Delta t} = -\frac{2}{3}\,\frac{\Delta\omega}{\omega}
   = +\frac{2}{3}\,\frac{\Delta T}{T}
   = 6.20 \times 10^{-7}\ \text{per year} \]
\hfill(1.0 points)
\end{description}

\solhead{8.2}{}
% source id: 2008_TH_L1
\ioaafig{2008_TH_L1-s1.pdf}

Binary period: $P = 30$ days $= 0.0821$ years.

Time of eclipse: $t_e = 8$ hours $= 0.0009$ years.

Time of total eclipse: $t_t = 1$ hour 18 minutes $= 0.0001$ years.

Radial velocity of primary star: $V_{r1} = 30\,\mathrm{km/s}
= 6.3285$\,AU/year.

The velocity of secondary star: $V_{r2} = 40\,\mathrm{km/s}
= 8.4380$\,AU/year.

For circular orbit
\[ V_r = \frac{2\pi a}{P} \quad\Longrightarrow\quad
   a = \frac{P V_r}{2\pi} \quad (a = \text{radius of orbit}) \]

For primary star: \hfill(10\%)
\[ a_1 = \frac{P V_{r1}}{2\pi}
   = \frac{(0.0821)(6.3285)}{2\pi}
   = 0.0827\ \mathrm{AU} = 1.24 \times 10^7\ \mathrm{km} \]

For secondary star: \hfill(10\%)
\[ a_2 = \frac{P V_{r2}}{2\pi}
   = \frac{(0.0821)(8.4380)}{2\pi}
   = 0.1103\ \mathrm{AU} = 1.65 \times 10^7\ \mathrm{km} \]
\[ a = a_1 + a_2 = 0.0827 + 0.1103
   = 0.1930\ \mathrm{AU} = 2.89 \times 10^7\ \mathrm{km} \]

The radius of the primary star can be determined from the equation:
\hfill(10\%)
\[ \frac{(t_e + t_t)}{P} = \frac{4R_1}{2\pi a}
   \qquad\text{or}\qquad
   R_1 = \frac{2\pi a}{4P}(t_e + t_t)
   = \frac{2\pi(0.1930)}{4(0.0821)}(0.0009 + 0.0001) \]
\[ = 0.0039\ \mathrm{AU} = 5.86 \times 10^5\ \mathrm{km}
   = 0.84\ R_\odot \]
\hfill(20\%)

The radius of the secondary star can be determined from the equation:
\hfill(10\%)
\begin{align*}
\frac{(t_e - t_t)}{P} = \frac{4R_2}{2\pi a}
  \qquad\text{or}\qquad
  R_2 &= \frac{2\pi a}{4P}(t_e - t_t) \\
  &= \frac{2\pi(0.1930)}{4(0.0821)}(0.0009 - 0.0001) \\
  &= 0.0028\ \mathrm{AU} = 4.22 \times 10^5\ \mathrm{km} = 0.61\ R_\odot
  && \text{(15\%)}
\end{align*}

From the Kepler's third law:
\[ \frac{a^3 \sin^3 i}{P^2} = \left(\mathcal{M}_1 + \mathcal{M}_2\right)
   \qquad\text{and}\qquad
   \frac{\mathcal{M}_1}{\mathcal{M}_2} = \frac{a_2}{a_1} \]

Since $i = 90^\circ$, $\sin i = \sin 90^\circ = 1$, and then
\[ \mathcal{M}_1 = \frac{a^3}{P^2(1 + a_1/a_2)}
   = \frac{(0.1930)^3}{(0.0821)^2(1 + 0.0827/0.1103)}
   = 0.61\ \mathcal{M}_\odot \]
\hfill(15\%)
\[ \mathcal{M}_2 = \frac{a^3}{P^2(1 + a_2/a_1)}
   = \frac{(0.1930)^3}{(0.0821)^2(1 + 0.1103/0.0827)}
   = 0.46\ \mathcal{M}_\odot \]
\hfill(10\%)

\solhead{8.3}{}
% source id: 2009_TH_P15
% NOTE: the official solution refers to "figure 1" (maximum-brightness
% configuration) and "Figure 2" (total eclipse, projected circle of radius b),
% but no such figures are printed in the source solutions PDF
% (2009_th_sol.pdf, pp. 54--55).

Ignoring temperature variation on the stars' surface, the brightness of the
star system will be proportional to the projected surface of both stars on the
plane of the sky. Maximum brightness will occur when the two stars are seen
like figure 1 and minimum light will happen when one of the stars is in total
eclipse and the projected surface of the other is a circle with radius $b$
(Figure 2). In maximum brightness
\[ I_{max} \propto 2\pi a b \]
\hfill(3 points)

In minimum brightness
\[ I_{min} \propto \pi b^2 \]
\hfill(3 points)

So
\[ \Delta m = -2.5\log\frac{I_{max}}{I_{min}}
   = -2.5\log\frac{2\pi ab}{\pi b^2}
   = -2.5\log 4 \]
\hfill(2 points)
\[ \Delta m = -1.5 \]
\hfill(2 points)

\solhead{8.4}{}
% source id: 2010_TH_S4
The semi-major axis is
\[ a = 1/2 \times (7 + 1) \times 10 = 40\ \mathrm{AU} \]
\hfill(2 points)

From Kepler's 3rd law,
\[ M_1 + M_2 = \frac{a^3}{p^2} = \frac{(40)^3}{(100)^2}
   = 6.4\,M_{\mathrm{sun}} \]
\hfill(4 points)

Since $a_1 = 3''$, $a_2 = 1''$, then
\[ \frac{m_1}{m_2} = \frac{a_2}{a_1} \]
\hfill(2 points)
\[ m_1 = 1.6\,M_{\mathrm{sun}}\,,\qquad m_2 = 4.8\,M_{\mathrm{sun}} \]
\hfill(2 points)

\solhead{8.5}{}
% source id: 2011_TH_L3
% The official 2011 solution for this long problem is published only as a
% worked marking scheme (max 30 points); it is transcribed as printed.

Difference in radial distance is about $1.58 \times 10^{15}$\,m
\hfill[2]\\
Difference in right ascension is about $2.4767^\circ$ \hfill[2]\\
Difference in declination is about $1.84^\circ$ \hfill[2]

Thus the angular distance on the sky is about $3.09^\circ$ \hfill[3]

Calculating the linear distance perpendicular to the line of sight:
using trigonometric functions and the distance from the Sun, we obtain a
distance of $2.1 \times 10^{15}$\,m. \hfill[2]

The distance between the stars is thus $2.63 \times 10^{15}$\,m. \hfill[3]

Calculate the difference in angular velocity in R.A.\ and Dec. Calculations
(same method) for R.A.\ and Dec.\ give a difference of 0.18\ arcsec/year in
both. \hfill[4]

Conversion of angular velocity on the sky to transverse linear velocity
(works out around 15). \hfill[4]

Note that transverse velocity is a lower limit to total velocity. \hfill[2]

Compare obtained velocity to velocity expected for the distances between
stars. At distances of tens of thousands of AU the velocity is too great for
the stars to form a bound system even with very massive stars. Even for a
pair of one light and one heavy star, the heavy one would have a mass of
thousands of solar masses. \hfill[4]

Final answer and conclusion: \hfill[2]

\solhead{8.6}{}
% source id: 2013_TH_S4
\begin{description}
\item[(a)] From Figure 2, % sic: the period is read off the light curve, i.e. Figure 3 of the statement
  the period of the cepheid is $P \sim 11$ days and
  its average apparent magnitude is $\sim(14.8 + 14.1)/2$ mag, i.e.\
  $m = 14.45$ mag. \hfill(2 points)

  [A careful student will notice that the graph is not upside-down symmetrical
  so he/she choose some value closer to the bottom; that is $m = 14.5$ mag.
  \hfill(1 point)]

  From Figure 1, % sic: the period-luminosity relation is Figure 2 of the statement
  we derive that for a period of 11 days the expected absolute
  magnitude of the cepheid is $M \approx -4.2$. \hfill(1 point)

  [A careful student will notice that the graph is logarithmic so he/she
  choose a value closer to 4.3] \hfill(1 point)

  Using the formula $m - M = -5 + 5\log r$, where $r$ is the distance of the
  cepheid, we get
  \[ \log r = (14.45 + 4.2 + 5)/5 = 4.73, \]
  thence $r = 10^{4.73} \approx 57500$\,pc or $r = 57.5$\,kpc.
  \hfill(3 points)
\item[(b)] Assuming $A = 0.25$, then
  \[ \log r = (14.45 + 4.2 + 5 + 0.25)/5 = 4.78, \]
  % sic: the extinction term should be subtracted (m - M - A = -5 + 5 log r),
  % which gives log r = 4.68 and r ~ 48 kpc.
  thence $r = 10^{4.78} \sim 53000$\,pc or $r = 64.5$\,kpc.
  % sic: 10^{4.78} = 60300 pc = 60.3 kpc; neither printed value matches, and
  % 53000 pc is inconsistent with 64.5 kpc.
  \hfill(2 points)
\end{description}

\solhead{8.7}{Absolute magnitude of a cepheide}
% source id: 2014_TH_P15
From
\[ P = 2\pi R\sqrt{\frac{R}{KM}} \]
results
\[ P^2 = \frac{4\pi^2 R^3}{KM}\,;\qquad
   R = \sqrt[3]{\frac{KMP^2}{4\pi^2}}
     = \left(\frac{KM}{4\pi^2}\right)^{1/3}\cdot P^{2/3}\,;
\]
\[ R^2 = \left(\frac{KM}{4\pi^2}\right)^{2/3}\cdot P^{4/3}. \]

The absolute brightness is:
\[ L_{\mathrm{cef}} = \sigma T_{\mathrm{cef}}^4 \cdot 4\pi R^2, \]
and the apparent brightness:
\[ E_{\mathrm{cef}} = \frac{L_{\mathrm{cef}}}{4\pi d_{\mathrm{P,cef}}^2}
   = \frac{\sigma T_{\mathrm{cef}}^4 \cdot 4\pi R^2}
          {4\pi d_{\mathrm{P,cef}}^2}, \]
where $d_{\mathrm{P,cef}}$ is the distance between the observer on Earth and
the cepheid,
\[ E_{\mathrm{cef}} = \frac{\sigma T_{\mathrm{cef}}^4\cdot 4\pi\cdot
   \left(\dfrac{KM}{4\pi^2}\right)^{2/3}\cdot P^{4/3}}
   {4\pi d_{\mathrm{P,cef}}^2}. \]

Similarly for the Sun
\[ E_{\mathrm{S}} = \frac{L_{\mathrm{S}}}{4\pi d_{\mathrm{PS}}^2}
   = \frac{\sigma T_{\mathrm{S}}^4 \cdot 4\pi R_{\mathrm{S}}^2}
          {4\pi d_{\mathrm{PS}}^2}. \]

By using the Pogson formula:
\[ \log\frac{E_{\mathrm{cef}}}{E_{\mathrm{S}}}
   = -0.4\left(m_{\mathrm{cef}} - m_{\mathrm{S}}\right); \]
\[ M_{\mathrm{cef}} = M_{\mathrm{S}}
   - 5^{\mathrm{m}}\log\frac{\left|d_{\mathrm{P,cef}}\right|}
                            {\left|d_{\mathrm{PS}}\right|}
   - 2.5\cdot\log\frac{E_{\mathrm{cef}}}{E_{\mathrm{S}}}\,; \]
% sic: the source prints the exponent 3/2; consistency with the earlier
% R^2 = (KM/4pi^2)^{2/3} P^{4/3} requires 2/3.
\[ \frac{T_{\mathrm{cef}}^4\cdot\left(\dfrac{KM}{4\pi^2}\right)^{3/2}
   \cdot d_{\mathrm{PS}}^2}
   {T_{\mathrm{S}}^4\cdot R_{\mathrm{S}}^2\cdot d_{\mathrm{P,cef}}^2}
   = k_1 = \text{constant}; \]
\begin{gather*}
M_{\mathrm{cef}} = M_{\mathrm{S}}
  - 5^{\mathrm{m}}\log\frac{\left|d_{\mathrm{P,cef}}\right|}
                           {\left|d_{\mathrm{PS}}\right|}
  - 2.5\cdot\log k_1 - \frac{10}{3}\cdot\log P\,; \\
M_{\mathrm{S}}
  - 5^{\mathrm{m}}\log\frac{\left|d_{\mathrm{P,cef}}\right|}
                           {\left|d_{\mathrm{PS}}\right|}
  - 2.5\cdot\log k_1 = -2.5\cdot\log k\,; \\
k = \text{constant}; \\
M_{\mathrm{cef}} = -2.5\cdot\log k - \frac{10}{3}\cdot\log P.
\end{gather*}

\solhead{8.8}{Mass function of a visual binary stellar system}
% source id: 2014_TH_P4
\textit{The official 2014 solutions file contains no solution for this problem (its marking-scheme index skips Problem 4).}

\solhead{8.9}{}
% source id: 2015_TH_L2
\begin{description}
\item[a.] The center of gravity of the stars is relatively to the star A given
  by
  \[ r_A = \frac{m_B}{m_A + m_B}\,d \]
  and since the orbit of A is a circle, then
  \[ F_{AX} = \frac{G m_A m_B}{d^2} = \frac{m_A v_A^2}{r_A} \]
  So, we get
  \[ v_A = m_B\sqrt{\frac{G}{(m_A + m_B)d}} \]
  The angular velocity of A is given by
  \[ \omega = \frac{v_A}{r_A} = \sqrt{\frac{G(m_A + m_B)}{d^3}} \]
  \hfill(30)
\item[b.] In Cartesian coordinate system, the velocity of A is
  \[ \vec{v}_A(t) = v_A\left(-\sin(\omega t + \varepsilon)\,\hat{\imath}
     + \cos(\omega t + \varepsilon)\,\hat{\jmath}\right) \]
  Unit vector of the observer is
  \[ \hat{r}_P = \cos\theta\,\hat{k} + \sin\theta\,\hat{\jmath} \]
  so the component of $\vec{v}_A$ in the line of the observer sight is given
  by
  \[ \vec{v}_A\cdot\hat{r}_P = v_A\sin\theta\,\cos(\omega t + \varepsilon) \]
  Since the component of $\vec{v}_A$ in the line of the observer sight is
  $K\cos(\omega t + \varepsilon)$, then
  \[ K = v_A\sin\theta \]
  Finally, we have
  \[ \frac{K^3}{\omega G} = \frac{m_B^3}{(m_A + m_B)^2}\sin^3\theta
     \qquad (**) \]
  \hfill(30)
\item[c.] From the result in b., namely eq.\ (**), we get
  \[ \sin^3\theta = \frac{K^3}{\omega G}\,
     \frac{(m_A + m_B)^2}{m_B^3}
     < \frac{1}{250}\,\frac{32^2}{8} = \frac{64}{125} \]
  Since $\theta \in [0, \pi]$, then $\sin\theta < 0.8$. Thus, the probability
  of B is a black hole is the same as the probability of $\sin\theta < 0.8$
  for $\theta \in [0, \pi]$. So $\theta$ is less than $53^\circ$ or greater
  than $127^\circ$.
  \hfill(40)
\end{description}

\solhead{8.10}{Cepheid Pulsations}
% source id: 2016_TH_T6
\begin{description}
\item[(T6.1)]
We first find the flux ratio and then use Stefan's law to compare the fluxes.
\[ m_1 - m_2 = -2.5\log\left(\frac{F_1}{F_2}\right) \]
\hfill(1.0 points)
\[ \Rightarrow\ \frac{F_1}{F_2} = 10^{-0.4(m_1 - m_2)}
   = 10^{-0.4(3.46 - 4.08)} = 1.77 \]
\hfill(1.0 points)
\[ L_i = 4\pi R_i^2\sigma T_i^4 \]
\hfill(1.0 points)
\[ \Rightarrow\quad
   F_1 = \frac{4\pi R_1^2\sigma T_1^4}{4\pi D^2}\,,\qquad
   F_2 = \frac{4\pi R_2^2\sigma T_2^4}{4\pi D^2} \]
\hfill(1.0 points)
\begin{align*}
\frac{F_1}{F_2} &= \frac{R_1^2}{R_2^2}\times\frac{T_1^4}{T_2^4} \\
\frac{R_1}{R_2} &= \sqrt{\frac{F_1}{F_2}}\times
  \left(\frac{T_2}{T_1}\right)^2
\end{align*}
From Wien's displacement law, $\dfrac{T_2}{T_1} = \dfrac{\lambda_1}{\lambda_2}$.
\hfill(1.0 points)
\begin{align*}
\Rightarrow\ \frac{R_1}{R_2}
  &= \sqrt{\frac{F_1}{F_2}}\times\left(\frac{\lambda_1}{\lambda_2}\right)^2
  && \text{(1.0 points)} \\
  &= \sqrt{1.77}\times\left(\frac{531.0}{649.1}\right)^2 \\
\frac{R_1}{R_2} &= 0.890
  && \text{(1.0 points)}
\end{align*}
Acceptable range: $\pm 0.010$.

\item[(T6.2)]
\[ R_2 - R_1 = v \times P/2 \]
\hfill(2.0 points)
\begin{gather*}
R_2 - R_1 = \frac{12.8\times 10^3\times 86\,400\times 9.84}{2}\ \mathrm{m} \\
(1 - 0.890)R_2 = 5.441\times 10^9\ \mathrm{m}
\end{gather*}
\[ \Rightarrow R_2 = 4.95\times 10^{10}\ \mathrm{m} \]
\hfill(0.5 points)
\[ R_1 = 4.41\times 10^{10}\ \mathrm{m} \]
\hfill(0.5 points)

Acceptable range: $\pm 0.02\times 10^{10}$\,m for both.

\item[(T6.3)]
To get the absolute value of flux ($F_2$) we must compare it with the observed
flux of the Sun.
\begin{gather*}
m_2 - m_\odot = -2.5\log\left(\frac{F_2}{F_\odot}\right) \\
\Rightarrow\ F_2 = F_\odot 10^{-0.4(m_2 - m_\odot)}
\end{gather*}
\hfill(3.0 points)
\begin{align*}
&= \frac{L_\odot}{4\pi a_\oplus^2}\times 10^{-0.4(4.08 + 26.72)} \\
&= \frac{3.826\times 10^{26}\times 4.7863\times 10^{-13}}
        {4\pi(1.496\times 10^{11})^2}\ \mathrm{W\,m^{-2}} \\
F_2 &= 6.51\times 10^{-10}\ \mathrm{W\,m^{-2}}
\end{align*}
\hfill(2.0 points)

Acceptable range: $\pm 0.04\times 10^{-10}\ \mathrm{W\,m^{-2}}$.

\item[(T6.4)]
\[ D_{\mathrm{star}} = \sqrt{\frac{L_2}{4\pi F_2}}
   = \sqrt{\frac{R_2^2\sigma T_2^4}{F_2}}
   = R_2 T_2^2\sqrt{\frac{\sigma}{F_2}} \]
\hfill(2.0 points)

Wien's law:
\[ T_2 = \frac{2.898\times 10^{-3}\ \mathrm{m\,K}}{\lambda_2} \]
\hfill(1.0 points)
\[ D_{\mathrm{star}} = 4.95\times 10^{10}\times
   \left(\frac{2.898\times 10^{-3}}{649.1\times 10^{-9}}\right)^2
   \sqrt{\frac{5.670\times 10^{-8}}{6.51\times 10^{-10}}} \]
\hfill(2.0 points)
\[ \Rightarrow\ D_{\mathrm{star}} = 9.208\times 10^{18}\ \mathrm{m}
   = 298\ \mathrm{pc} \]

Acceptable range: $298 \pm 2$\,pc (depends on truncation).
\end{description}

\solhead{8.11}{Identification of light curves of types of selected variable stars}
% source id: 2019_TH_11
\begin{description}
\item[a)] The score for each correct answer: 1 point. \hfill(8 p)

  \begin{center}
  \begin{tabular}{ll}
  Heart-beat star & 8 \\
  RR Lyrae type (RRab subclass) pulsating variable star & 4 \\
  Eclipsing binary of Algol type (semi-detached) with a pulsating component
    & 6 \\
  $\alpha^2$ CVn star & 5 \\
  W Vir type (Population II) Cepheid pulsating variable & 3 \\
  Detached eclipsing binary with strong reflection effect & 2 \\
  Contact eclipsing binary of W UMa type & 1 \\
  Rotationally variable (spotted) star & 7 \\
  \end{tabular}
  \end{center}
\item[b)] The values below have been determined by an AoV period searching
  algorithm and rounded to 2 decimal places. The range of $\pm 5$\,\% is also
  given for each value: the period in this range is acceptable. The score for
  each correct answer: 2 points. \hfill(16 p)

  \begin{center}
  \begin{tabular}{llll}
  1. & RW Dor & $P \approx 0.29$\,d
     & $0.28\,\mathrm{d} \le P \le 0.30\,\mathrm{d}$ \\
  2. & FO Eri & $P \approx 2.20$\,d
     & $2.09\,\mathrm{d} \le P \le 2.31\,\mathrm{d}$ \\
  3. & UY Eri & $P \approx 2.21$\,d
     & $2.10\,\mathrm{d} \le P \le 2.32\,\mathrm{d}$ \\
  4. & ST Pic & $P \approx 0.49$\,d
     & $0.47\,\mathrm{d} \le P \le 0.51\,\mathrm{d}$ \\
  5. & AH Col & $P \approx 1.10$\,d
     & $1.04\,\mathrm{d} \le P \le 1.16\,\mathrm{d}$ \\
  6. & VV Ori & $P \approx 1.49$\,d
     & $1.42\,\mathrm{d} \le P \le 1.56\,\mathrm{d}$ \\
  7. & TIC 147272181 & $P \approx 0.55$\,d
     & $0.52\,\mathrm{d} \le P \le 0.58\,\mathrm{d}$ \\
  8. & 24 Eri & $P \approx 8.12$\,d
     & $7.71\,\mathrm{d} \le P \le 8.53\,\mathrm{d}$ \\
  \end{tabular}
  \end{center}
\item[c)] (D) Supernova in a distant galaxy (ASASSN-18tb) \hfill(1 p)
\end{description}

\solhead{8.12}{Amplitude variation of RR Lyrae stars}
% source id: 2019_TH_8
\begin{description}
\item[a)] The total power of the star in the given lambda:
  \begin{equation*}
  I(\lambda, T, R) = C_1 R^2 F(\lambda, T)
    = C_1 R^2 \frac{1}{\lambda^5}
      \exp\left(-\frac{C_b}{\lambda T}\right) \tag{8.1}
  \end{equation*}
  \hfill(1 p)

  Converting to magnitudes:
  \begin{equation*}
  m(\lambda, T, R) = -2.5\log I(\lambda, T)
    = C_2 + 12.5\log\lambda - 5\log R
      - 2.5\log e\,\frac{C_b}{\lambda T} \tag{8.2}
  % sic: expanding -2.5 log I gives +2.5 log e C_b/(lambda T); the sign slip
  % is absorbed by the absolute value in (8.3).
  \end{equation*}
  \hfill(1 p)

  Then:
  \begin{equation*}
  A(\lambda) = \left|m(\lambda, T_{\mathrm{max}})
    - m(\lambda, T_{\mathrm{min}})\right|
    = \frac{2.5\log e\,C_b}{\lambda}
      \left(\frac{1}{T_{\mathrm{min}}}
      - \frac{1}{T_{\mathrm{max}}}\right) \tag{8.3}
  \end{equation*}
  \hfill(1 p)

  Thus:
  \begin{equation*}
  \frac{A(\lambda_1)}{A(\lambda_2)} = \frac{\lambda_2}{\lambda_1} \tag{8.4}
  \end{equation*}
  \hfill(1 p)

  With numerical values:
  \[ \frac{A(\lambda_1)}{A(\lambda_2)}
     = \frac{2\times 10^{-6}\ \mathrm{m}}{5\times 10^{-7}\ \mathrm{m}}
     = 4 \]
  \hfill(1 p)
\item[b)] From Eq.\ (8.3):
  \begin{equation*}
  A_{\mathrm{max}}(\lambda_1) = \frac{2.5\log e\,C_b}{\lambda_1}
    \left(\frac{1}{T_{\mathrm{min}}}
    - \frac{1}{T_{\mathrm{max}}}\right) \tag{8.5}
  \end{equation*}
  \hfill(2 p)

  With numerical values:
  \[ A_{\mathrm{max}}(\lambda_1)
     = \frac{2.5\log e\times 0.0144\ \mathrm{m\,K}}
            {5\times 10^{-7}\ \mathrm{m}}
       \left(\frac{1}{6000\ \mathrm{K}}
       - \frac{1}{7400\ \mathrm{K}}\right)
     = 1^{\mathrm{m}} \]
  \hfill(1 p)
\item[c)] If we ignore the temperature variation, then:
  \begin{equation*}
  A(\lambda) = \left|m(\lambda, R_{\mathrm{max}})
    - m(\lambda, R_{\mathrm{min}})\right|
    = 5\log\frac{R_{\mathrm{max}}}{R_{\mathrm{min}}} \tag{8.6}
  \end{equation*}
  \hfill(2 p)

  With numerical values:
  \[ A(\lambda) = 5\log\frac{1.05}{0.9} = 0.33^{\mathrm{m}} \]
  \hfill(1 p)
\item[d)] From Eq.\ (8.4):
  \begin{equation*}
  A_{\mathrm{max}}(\lambda_2)
    = A_{\mathrm{max}}(\lambda_1)\,\frac{\lambda_1}{\lambda_2}
    = 1^{\mathrm{m}}\times\frac{1}{4} = 0.25^{\mathrm{m}} \tag{8.7}
  \end{equation*}
  \hfill(1 p)

  Similarly:
  \[ A_{\mathrm{min}}(\lambda_1)
     = \frac{2.5\log e\times 0.0144\ \mathrm{m\,K}}
            {5\times 10^{-7}\ \mathrm{m}}
       \left(\frac{1}{6100\ \mathrm{K}}
       - \frac{1}{6900\ \mathrm{K}}\right)
     = 0.6^{\mathrm{m}} \]
  \hfill(2 p)

  and
  \[ A_{\mathrm{min}}(\lambda_2)
     = A_{\mathrm{min}}(\lambda_1)\,\frac{\lambda_1}{\lambda_2}
     = 0.6^{\mathrm{m}}\times\frac{1}{4} = 0.15^{\mathrm{m}} \]
  \hfill(2 p)

  At $\lambda_1 = 500$\,nm the amplitude changes by
  $1^{\mathrm{m}} - 0.6^{\mathrm{m}} = 0.4^{\mathrm{m}}$, while at
  $\lambda_2 = 2000$\,nm it is reduced to
  $0.25^{\mathrm{m}} - 0.15^{\mathrm{m}} = 0.1^{\mathrm{m}}$, that is a
  significant difference.
\item[e)] B and C \hfill(4 p)
\end{description}

\solhead{8.13}{Light Curves}
% source id: 2020_TH_4
\begin{description}
\item[(a)] As ``Y'' is smaller, yet hotter, it must be on the main sequence.
  \hfill(1.0)
\item[(b)]
  \begin{description}
  \item[(I)] For this we have to measure the total duration of eclipse and
    compare it with the bottom flat part (total eclipse phase) in the
    lightcurve. Total eclipse time should be $(1.25 \pm 0.10)$ days, flat
    period of eclipse $(0.25 \pm 0.05)$ days. \hfill(1.0)
    \[ \frac{r_X + r_Y}{r_X - r_Y} = \frac{1.25 \pm 0.10}{0.25 \pm 0.05}
       = 5.0 \pm 1.5 \]
    \[ \therefore\ \frac{r_X}{r_Y} = \frac{5 + 1}{5 - 1} = 1.5 \pm 0.2 \]
    \hfill(1.0)
  \item[(II)] Because star Y is hotter, star Y will be brighter than a part
    of star X with the same area. Thus, the deeper dip in the light curve is
    due to star Y passing behind star X. Thus, the flux from star X must be
    $10 \times 10^{-12}\ \mathrm{W/m^2}$. As the combined flux is
    $16 \times 10^{-12}\ \mathrm{W/m^2}$, the flux of star Y is
    $6 \times 10^{-12}\ \mathrm{W/m^2}$. Thus
    \[ \frac{L_X}{L_Y} = \frac{10 \times 10^{-12}}{6 \times 10^{-12}}
       = 1.67 \]
    \hfill(2.0)
  \end{description}
\item[(c)]
  \begin{itemize}
  \item Light Curve B: We see that the dips in brightness both become
    shallower and have become rounded at the bottom, while the maximum
    brightness is still $16 \times 10^{-12}\ \mathrm{W/m^2}$. If the radius
    of one of the stars had changed (rounded bottom), then we would expect
    the total eclipse period to also change. But that has not happened.
    Brightening of star X can explain a shallower dip but not a rounded
    trough. This means that the stars no longer fully eclipse each other.
    Hence,

    (xi): The inclination of the system relative to the Earth has changed.
    \hfill(4.0)
  \item Light Curve C: The y-axis shows that the overall flux has increased.
    Since the flux of both the stars has gone up by the same factor and there
    is no substantial change in the profile of eclipse,

    (xii): The distance of the system from the Earth has decreased.
    \hfill(4.0)
  \item Light Curve D: The regular fluctuations in brightness indicate that
    one of the components of the system is a variable star. Since the
    fluctuations persist through all the dips in flux, this indicates that
    the variable star is never completely eclipsed. Therefore,

    (ix): Star X is a variable star. \hfill(3.0)
  \item Light Curve E: The overall flux and dips in the flux are the same
    size, but occur much more frequently. This means that

    (xv): The orbital period of the system decreased. \hfill(2.0)
  \item Light Curve F: The regular fluctuations again indicate that one of
    the components of the system is a variable star. During the primary
    eclipse, the light curve is perfectly flat. This can only occur if,

    (x): Star Y is a variable star. \hfill(2.0)
  \end{itemize}
\end{description}

\solhead{8.14}{Menkalinan (beta Aurigae)}
% source id: 2021_TH_Q7
\ioaafig{2021_TH_Q7-s1.pdf}

\begin{description}
\item[(7.1)] Observing the plot, it is possible to estimate the wavelength of
  both H$\alpha$ lines and later find the velocity of each star by Doppler.
  \[ v_A = \frac{6561\,\text{\AA} - 6562.8\,\text{\AA}}{6562.8\,\text{\AA}}
     \times\left(2.998\cdot 10^5\ \mathrm{km\cdot s^{-1}}\right)
     = -82.23\ \mathrm{km\cdot s^{-1}} \]
  \hfill[2 point]
  \[ v_B = \frac{6565\,\text{\AA} - 6562.8\,\text{\AA}}{6562.8\,\text{\AA}}
     \times\left(2.998\cdot 10^5\ \mathrm{km\cdot s^{-1}}\right)
     = +100.50\ \mathrm{km\cdot s^{-1}} \]
  \hfill[2 point]
  So, Menkalinan A is approaching and Menkalinan B is moving away from us.
  \hfill[1 point]
\item[(7.2)]
  \[ m_A > m_B \ \Rightarrow\ \alpha_B > \alpha_A \]
  From the definition of CM:
  \[ \frac{\alpha_B}{\alpha_A} = 1.026 \]
  \hfill[0.5 points]
  \[ \alpha = (\alpha_A + \alpha_B)
     = \alpha_B\left(1 + \frac{\alpha_A}{\alpha_B}\right)
     = 1.975\,\alpha_B \]
  \hfill[0.5 points]
  Applying the 3rd Kepler's law, the total mass of the system is:
  \[ M_{tot} = \frac{4\pi^2}{G}\times\frac{\alpha^3\cdot d^3}{T^2} \]
  \hfill[1 point]
  \[ M_{tot} = \frac{4\pi^2}{G}\times
     \frac{\left(0.00662''\times\dfrac{1\,\mathrm{rad}}{206265''}\right)^3
       \cdot\left(81.1\,\mathrm{ly}\cdot
       \dfrac{9.46\cdot 10^{15}\,\mathrm{m}}{1\,\mathrm{ly}}\right)^3}
     {\left(3.96\,\mathrm{d}\cdot
       \dfrac{86400\,\mathrm{s}}{1\,\mathrm{d}}\right)^2} \]
  \hfill[2 point]
  \[ M_{tot} = 7.528\cdot 10^{31}\ \mathrm{kg} = 37.868\,M_\odot \]
\item[(7.3)]
  \[ m_A/m_B = 1.026 \]
  \[ M_{tot} = m_A + m_B = 1.026\cdot m_B + m_B = 2.026\,m_B
     = 37.868\,M_\odot \]
  \[ m_B = 18.691\,M_\odot \]
  \hfill[1 point]
  \[ m_A = 1.026\cdot m_B = 19.177\,M_\odot \]
  \hfill[1 point]
\item[(7.4)] Using the mass/luminosity relation given:
  \[ \frac{L_A}{L_\odot} = \left(\frac{m_A}{M_\odot}\right)^{3.5}
     = \left(\frac{19.177\,M_\odot}{M_\odot}\right)^{3.5}
     = 3.088\times 10^4 \]
  \[ L_A = 3.088\times 10^4\,L_\odot \]
  \hfill[1 point]
  \[ \frac{L_B}{L_\odot} = \left(\frac{m_B}{M_\odot}\right)^{3.5}
     = \left(\frac{18.691\,M_\odot}{M_\odot}\right)^{3.5}
     = 2.823\times 10^4 \]
  \[ L_B = 2.823\times 10^4\,L_\odot \]
  \hfill[1 point]
\end{description}

\solhead{8.15}{Photometry of Binary stars}
% source id: 2022_TH_6
\begin{description}
\item[(a)] From the Wien's displacement law the temperature of a black body is
  related to a peak wavelength of emission by the relation:
  \[ \lambda_{max} T = b \]
  \hfill(1.0)
  \begin{align*}
    T_A &= 5978\,\mathrm{K} \\
    T_B &= 4830\,\mathrm{K}
  \end{align*}
  \hfill(0.5 + 0.5)

\item[(b)] Due to the diffraction minimal angular separation between two
  celestial objects to distinguish them:
  \[ \delta = 1.22\,\frac{\lambda}{D} = 4.44 \times 10^{-8}\,\mathrm{rad}
     = 9.16\,\mathrm{mas} \]
  \hfill(1.0)

  Now let us derive the time dependence of the angular separation between two
  stars as seen from the earth.

  \ioaafig{2022_TH_6-s1.pdf}

  As it is seen from the figure, distance $s = l\cos\theta$, where $l$ is the
  distance between stars and $\theta$ is the angle of rotation of the stars
  about their center of mass after the stars are seen with maximum angular
  separation. The angular separation seen by an observer on the earth:
  \[ \gamma = \frac{s}{d} = \frac{l}{d}\cos\theta \]
  \hfill(1.0)

  The 38 days, during which objects are seen as a one object the system rotates
  by the angle $\beta = 2\pi\,\frac{38}{100}$. \hfill(1.0)

  Thus, the corresponding starting angle $\theta'$ at which stars are no longer
  distinguishable:
  \[ \theta' = \frac{\pi}{2} - \frac{\beta}{4}
     = \frac{\pi}{2} - \frac{19\pi}{100} = \frac{31\pi}{100} \]
  \hfill(1.5)

  and this is the moment, when stars stop to be distinguishable. Thus,
  corresponding angular separation on the sky is $\delta$:
  \[ \delta = \frac{l}{d}\cos\theta' \]
  From which one obtains
  \[ l = \frac{\delta d}{\cos\theta'} = 1.45\,\mathrm{au} \]
  \hfill(1.5)

\item[(c)] From the Kepler's 3rd law:
  \begin{align*}
    T^2 &= \frac{4\pi^2 l^3}{G(M_A + M_B)} \\
    M_A + M_B &= \frac{4\pi^2 l^3}{G T^2} = 40.7\,M_\odot
  \end{align*}
  \hfill(2.0)

\item[(d)] From given data for case 1:
  \begin{align*}
    (U - V)_1 &= (U - B)_1 + (B - V)_1 = 0.3 \\
    U_1 &= U_{0_1} + a_U d = 6.51 \\
    V_1 &= U_1 - (U - V)_1 = 6.21 \\
    m_{Bol_1} &= V_1 + BC_1 = 6.31
  \end{align*}
  \hfill(1.5)

  Similarly for second configuration:
  \begin{align*}
    (U - V)_2 &= (U - B)_2 + (B - V)_2 = 0.37 \\
    U_2 &= U_{0_2} + a_U d = 6.98 \\
    V_2 &= U_2 - (U - V)_2 = 6.61 \\
    m_{Bol_2} &= V_2 + BC_2 = 6.79
  \end{align*}
  \hfill(1.0)

  The difference in bolometric magnitudes:
  \begin{align*}
    m_{bol_2} - m_{bol_1}
      &= -2.5\log\left(\frac{L_2}{L_1}\right) \\
      &= -2.5\log\left(
         \frac{\pi(R_A^2 - R_B^2)\sigma T_A^4 + \pi R_B^2 \sigma T_B^4}
              {\pi R_A^2 \sigma T_A^4 + \pi R_B^2 \sigma T_B^4}\right)
  \end{align*}
  \hfill(1.5)
  \begin{align*}
    \therefore\ 6.79 - 6.31
      &= -2.5\log\left(
         \frac{(R_A^2 - R_B^2)T_A^4 + R_B^2 T_B^4}
              {R_A^2 T_A^4 + R_B^2 T_B^4}\right) \\
      &= -2.5\log\left(
         \frac{(R_A^2/R_B^2 - 1)T_A^4 + T_B^4}
              {T_A^4 R_A^2/R_B^2 + T_B^4}\right) \\
    \therefore\ 0.48
      &= -2.5\log\left[
         \frac{\left(\frac{R_A^2}{R_B^2} - 1\right) + \frac{T_B^4}{T_A^4}}
              {\frac{T_B^4}{T_A^4} + \frac{R_A^2}{R_B^2}}\right]
  \end{align*}
  \hfill(1.0)
  \[ \therefore\ 10^{\frac{-0.48}{2.5}}
     = \frac{\frac{R_A^2}{R_B^2} - 1 + \frac{T_B^4}{T_A^4}}
            {\frac{T_B^4}{T_A^4} + \frac{R_A^2}{R_B^2}}
     = 10^{-0.19} = 0.65 \]
  \hfill(1.0)
  \begin{align*}
    \frac{R_A^2}{R_B^2} - 1 + \frac{T_B^4}{T_A^4}
      &= 0.65\,\frac{T_B^4}{T_A^4} + 0.65\,\frac{R_A^2}{R_B^2} \\
    0.35\,\frac{R_A^2}{R_B^2}
      &= 1 - 0.35\,\frac{T_B^4}{T_A^4} \\
      &= 1 - 0.35 \times 0.482 = 0.83 \\
    \therefore\ \frac{R_A}{R_B}
      &= \sqrt{\frac{0.83}{0.35}} = 1.54
  \end{align*}
  \hfill(1.0)
  \[ \frac{M_A}{M_B} = \frac{\rho_A R_A^3}{\rho_B R_B^3}
     = 1.54^3 \cdot 0.7 = 2.56 \]
  \hfill(1.0)
  \[ \text{But, } M_A + M_B = 40.7\,M_\odot \]
  \hfill(1.0)
  \[ \therefore\ M_B = \frac{40.7\,M_\odot}{3.56} = 11.4\,M_\odot \]
  \hfill(1.0)
  \[ M_B = 29.3\,M_\odot \]
  % sic — the source writes M_B again; this last line is clearly M_A
  % (= 40.7 - 11.4 = 29.3 solar masses).
\end{description}

\solhead{8.16}{Second eclipse}
% source id: 2023_TH_Q9
% Official solution, 2023 theory solutions pp. 22-24. The interleaved
% marking-rubric tables ("Noticing that R1 = R2 ... 1 point", etc.) are
% omitted per transcription policy.

For each of the systems, since the eclipses are short we can assume that the
angle swept along the orbit during the eclipse is negligibly small, which means
that the tangential velocity in relative orbital motion is constant.

\medskip
\noindent\textbf{System A}

A sharp peak minimum in a central eclipse implies that $R_1 = R_2$. Therefore,
both eclipses will be total, with minima of brightnesses $m_1$ and $m_2$,
respectively.

Define the magnitude at minimum as $m_{\min} = m_1$ and the baseline magnitude
as $m_{1,2} = m_1 + m_2$.
% sic — the source writes "m_{1,2} = m_1 + m_2"; m_{1,2} is really the
% combined magnitude of both stars, not a literal sum of magnitudes.
From Pogson's law applied to the two cases:
$m_1 - m_{1,2} = -2.5\log_{10}\left(L_1/(L_1 + L_2)\right)$ and
$m_2 - m_1 = -2.5\log_{10}\left(L_2/L_1\right)$. For the first case, with
simple operations we obtain: $L_2/L_1 = 10^{(m_1 - m_{1,2})/2.5} - 1$.
Plugging this into the second case:
\[ m_2 = m_1 - 2.5\log_{10}\left(10^{(m_1 - m_{1,2})/2.5} - 1\right). \]

As the orbits are near-circular, timescales will be the same for both eclipses.
Therefore it is sufficient to `copy' the shape of the primary one while scaling
down the depth to $|m_{1,2} - m_2|$. Reading off $m_{1,2}$ and $m_{\min}$ from
the plot, we arrive at $L_2/L_1 = 2$, $m_2 = 11.44$, and the following figure:

\ioaafig{2023_TH_Q9-s1.pdf}

\medskip
\noindent\textbf{System B}

For the second system, we follow a generally similar scheme, but we see that
the minimum is flat. $t_{\mathrm{fall}}$ is the time of the descent from the
baseline brightness to the lowest point, i.e., between contacts I and II, while
$t_{\mathrm{flat}}$ is the time between contacts II and III during which the
lightcurve remains at its lowest point:

\ioaafig{2023_TH_Q9-s2.pdf}

From this, we obtain
$t_{\mathrm{flat}}/t_{\mathrm{fall}} = 2(R_2 - R_1)/2R_1$ and, after
rearranging, $R_2/R_1 = t_{\mathrm{flat}}/t_{\mathrm{fall}} + 1$.

Here $R_1$ is the radius of the smaller star and there are two solutions --- it
could either be (partially) \emph{eclipsing} or (entirely) \emph{eclipsed}.
Let us first assume that it was eclipsed. After measuring $m_{\min}$ for the
eclipse and $m_{1,2}$ for the baseline we again (just like for system A) assume
$m_{\min} = m_1$ and write down: $L_1/L_2 = 10^{(m_1 - m_{1,2})/2.5} - 1$.

Reading off $m_{1,2}$ and $m_{\min}$ from the plot, we arrive at $L_2/L_1 = 4$,
$R_2/R_1 = 2$.\\
$L_i = S_i \pi R_i^2$, where $S_i$ stands for the surface brightness of each
star.\\
$\implies S_1 = S_2 = S$ (or
$T_{\mathrm{eff},1} = T_{\mathrm{eff},2} = T_{\mathrm{eff}}$).

In such a case, \emph{it does not matter} if star 2 is eclipsing or eclipsed.
During the baseline, we always see two disks with a total brightness of
$S(\pi R_1^2 + \pi R_2^2)$, while during the eclipse, we see one disk of radius
$R_2$ and total brightness $S\pi R_2^2$.

\textbf{Note about marking:} This solution can be independently verified by
assuming that the smaller component was eclipsing and cutting out a disk of
surface $\pi R_1^2$ and brightness contribution $S_1 \pi R_1^2$ in the star of
surface $\pi R_2^2$ and surface brightness $S_2$. Instead of
$m_{\min} = m_1$ of star 1, we would have measured
$m_{\min} = m_{\mathrm{ecl},1}$ of a combination of the two, with a brightness
$L_{ecl1} = S_1 \pi R_1^2 + S_2 \pi (R_2^2 - R_1^2)$. The same conclusion
$S_1 = S_2 = S$ will follow. Any solution correctly arriving at this conclusion
and the correct figure --- no matter the order of steps taken --- should be
scored equally. For the maximum number of points, however, the student should
make a convincing point that this is the \emph{only} solution, i.e.\ lifting
the previously made assumption (or exploring both cases).

To draw the second eclipse, one must simply copy the shape of the first one:

\ioaafig{2023_TH_Q9-s3.pdf}

\solhead{8.17}{Binary Hardening}
% source id: 2024_TH_T8
\begin{description}
\item[(a)] Let $v_p$ be the closest approach velocity of the star. From
  conservation of mechanical energy:
  \[ -\frac{G M_T m}{r_p} + \frac{1}{2} m v_p^2 = \frac{1}{2} m v^2 \]
  \hfill(2.0)

  where $M_T = 2M$ is the total mass of the binary. From conservation of
  angular momentum, we obtain $v_p$:
  \[ v_p = \frac{b}{r_p}\,v \]
  \hfill(2.0)

  setting $r_p \approx a/2$ and inserting $v_p$ into the first equation, we
  solve for $b$:
  \begin{align*}
    b &= a\sqrt{\frac{1}{4} + \frac{2GM}{a v^2}} \\
    b &\approx \frac{\sqrt{2GMa}}{v}
  \end{align*}
  \hfill(1.0)

  Where we have used $v^2 \ll \frac{GM}{a}$ in the last step.

  \textbf{Note:} The student can also solve it by exploring the geometry of
  the star's orbit, which is hyperbolic, by writing
  $r_p = a_{orbit}(e - 1)$ and finding $a_{orbit}$ and $e$ by writing out
  equations for the angular momentum and mechanical energy.

\item[(b)] The speed of the black hole in circular orbit can be found by
  equating the gravitational force to the centripetal net force:
  \[ \frac{GM^2}{a^2} = \frac{M V^2}{\frac{a}{2}}
     \qquad \therefore\ V = \sqrt{\frac{GM}{2a}} \]
  \hfill(1.0)

  For the rest of the problem, we propose two different solutions:

  \textbf{Solution I:}
  Additionally, let $v_f$ be the ejection speed of the star, and $V'$ the
  final speed of the black hole. To find $v_f$, we use conservation of linear
  momentum and conservation of mechanical energy for the two bodies:
  \[ M V = M V' + m v_f \]
  \hfill(1.5)
  \[ \frac{1}{2} M V^2 = \frac{1}{2} m v_f^2 + \frac{1}{2} M V'^2 \]
  \hfill(1.5)
  \begin{align*}
    \frac{1}{2} M V^2
      &= \frac{1}{2} m v_f^2
       + \frac{1}{2} M \left( V - \frac{m}{M} v_f \right)^2 \\
    0 &= v_f^2 - 2 V v_f + \frac{m}{M} v_f^2
  \end{align*}
  Since $m \ll M$ \hfill(1.0)\\
  we neglect the last term, yielding:
  \[ v_f = 2V = \sqrt{\frac{2GM}{a}} \]
  \hfill(1.0)

  Because we are considering two moments where the distance between the star
  and the black hole is very large, we disregard potential energy terms.

  \textbf{Solution II:}
  In the reference frame of the component, the star approaches the component
  with a velocity $-\vec{V}$. \hfill(1.5)

  And, since mechanical energy and linear momentum are conserved, it can be
  shown that the relative speed between the components is the same before and
  after the interaction --- that is, restitution coefficient equal to unity.
  Furthermore, since $M \gg m$, the velocity of the component can be
  considered unchanged (notice this is not valid in Solution I, as there we
  deal with second order terms arising from energy considerations) in the
  interaction. Therefore, after the interaction, the star is ejected back with
  a velocity $\vec{v}\,'_f = \vec{V}$ in the opposite direction. \hfill(1.5)

  To return to the CM frame, we write $\vec{v}_f = \vec{v}\,'_f + \vec{V}$,
  such that
  \[ v_f = 2V \]
  \hfill(2.0)

  As we have previously found.

\item[(c)] First, we know the kinetic energy acquired by a star during an
  encounter:
  \[ \Delta K_* = \frac{1}{2} m v_f^2 - \frac{1}{2} m v^2
     \approx \frac{GMm}{a} \]
  \hfill(1.0)

  Where we used that $v^2$ is negligible with respect to $v_f^2$, since
  $v_f^2 = \frac{2GM}{a}$, and $v^2 \ll \frac{GM}{a}$. This is also the
  energy lost by the binary during an encounter.

  In order to find the rate of energy extraction, we must first obtain the
  rate of encounters, which can be estimated as follows: imagine the binary
  travelling with velocity $v$ in the reference frame of a far away star of
  impact parameter $b$. During a small interval $dt$, all the stars included
  inside the annular cylinder of length $L = v\,dt$ bounded by radii of $b$
  and $b + db$ will interact with the star. Call this number of stars $dN$.
  It is given by:
  \begin{align*}
    dN &= n \cdot V_{\mathrm{cylinder}} \\
       &= n \cdot A_{\mathrm{annulus}} \cdot L \\
       &= n \cdot \pi \cdot \left[ (b + db)^2 - b^2 \right] \cdot v \cdot dt \\
       &\approx 2\pi n\, b\, v\, db\, dt
  \end{align*}
  \hfill(4.0)

  Defining $d\varphi \equiv \frac{dN}{dt}$ to be the rate of encounters with
  stars of impact parameter between $b$ and $b + db$ and using $b$ from part
  (a),
  \begin{align*}
    d\varphi &= 2\pi n\, b\, v\, db \\
             &= 2\pi n \cdot \frac{\sqrt{2GMa}}{v} \cdot v \cdot db \\
    d\varphi &= 2\pi n \sqrt{2GMa} \cdot db
  \end{align*}
  \hfill(1.0)

  To find the total rate of encounters ($\varphi = \sum d\varphi$), we sum
  over all impact parameters in the interval
  $\frac{1}{2} b_0 \le b \le \frac{3}{2} b_0$. Notice
  $\sum db = \frac{3}{2} b_0 - \frac{1}{2} b_0 = b_0$, so that:
  \begin{align*}
    \varphi &= 2\pi n \sqrt{2GMa} \cdot \sum db \\
            &= 2\pi n \sqrt{2GMa} \cdot \frac{\sqrt{2GMa}}{v_0} \\
            &= \frac{4\pi n G M a}{v_0}
  \end{align*}
  \hfill(3.0)

  Since each encounter takes away $\Delta K_*$ and there are $\varphi$
  encounters per unit time, the rate of change of energy for the binary is:
  \[ \frac{dE_{\mathrm{bin}}}{dt} = (-\Delta K_*) \cdot \varphi
     = -\frac{4\pi n G^2 M^2 m}{v_0} \]
  \hfill(2.0)

  A change in the binary total energy can also be written as:
  \[ \frac{dE_{\mathrm{bin}}}{dt}
     = \frac{d}{dt}\left( -\frac{G M^2}{2a} \right)
     = -\frac{G M^2}{2} \frac{d}{dt}\left( \frac{1}{a} \right) \]
  \hfill(2.0)

  Therefore, equating the two:
  \[ \frac{d}{dt}\left( \frac{1}{a} \right) = \frac{8\pi G n m}{v_0} \]
  \[ \frac{d}{dt}\left( \frac{1}{a} \right) = 8\pi\,\frac{G\rho}{v_0} \]
  \hfill(1.0)

  Thus, $H = 8\pi$.
\end{description}

\solhead{8.18}{Stellar masses in a visual binary system}
% source id: 2008_DA_D2
\begin{description}
  \item[1.] \ioaafig{2008_DA_D2-s1.pdf}
  \hfill(20\%)

  \item[2.] \ioaafig{2008_DA_D2-s2.pdf}

  \begin{itemize}
  \item For $\alpha$ Cen.\ A., its colour index is given by $B-V = 0.65$.
    Then, by taking the line that intersecting the above curve we find
    $BC = 0.10$. \hfill(10\%)
  \item For $\alpha$ Cen.\ B., its colour index is given by $B-V = 0.85$.
    Then, by taking the line that intersecting the above curve we find
    $BC = 0.25$. \hfill(10\%)
  \end{itemize}
  From the relation $m_{\mathrm{v}} - m_{\mathrm{bol}} = BC$ it is then
  obtained,
  \begin{itemize}
  \item For $\alpha$ Cen.\ A, $m_{\mathrm{v}} = -0.01$, we get
    $m_{\mathrm{bol}} = -0.01 - 0.10 = -0.11$ \hfill(10\%)
  \item For $\alpha$ Cen.\ B, $m_{\mathrm{v}} = 1.34$, we get
    $m_{\mathrm{bol}} = 1.34 - 0.25 = 1.09$ \hfill(10\%)
  \end{itemize}

\item[3.] Periodic orbit of $\alpha$ Centauri star is $P = 79.24$ years and
  its angle of semi major axis is $\alpha = 17.59''$. From the result 2a we
  find that the apparent bolometric magnitude of $\alpha$ Centauri A is
  $m_{\mathrm{bol}A} = -0.11$, while the apparent bolometric magnitude of
  $\alpha$ Centauri B is $m_{\mathrm{bol}B} = 1.09$ (two decimal places).

  \emph{1st iteration:}
  \begin{description}
  \item[Step 1:] Let $\mathcal{M}_A$ and $\mathcal{M}_B$ be the mass of
    Centauri A and Centauri B stars, respectively. Then, we firstly take
    $\mathcal{M}_A + \mathcal{M}_B = 2$.
  \item[Step 2:] Inserting the values of $\mathcal{M}_A + \mathcal{M}_B$,
    $P$, and $\alpha$ to the third Kepler equation namely,
    \[ \left( \frac{\alpha}{p} \right)^3
       = (\mathcal{M}_A + \mathcal{M}_B) P^2 \]
    One gets $p$, that is
    \[ p = \left[ \frac{\alpha^3}{(\mathcal{M}_A + \mathcal{M}_B) P^2}
           \right]^{1/3}
         = \left[ \frac{(17.59)^3}{2 (79.24)^2} \right]^{1/3}
         = \left[ 0.4334 \right]^{1/3} = 0.7568'' \]
  \item[Step 3:] We determine the magnitude of absolute bolometric using
    Pogson equation,
    \[ m_{bol} - M_{bol} = -5 - 5\log p
       \;\Longrightarrow\; M_{bol} = m_{bol} + 5 + 5\log p \]
    For $\alpha$ Centauri A:
    \begin{align*}
      M_{bolA} = m_{bolA} + 5 + 5\log p &= -0.11 + 5 + 5\log(0.7568) \\
        &= -0.11 + 5 - 0.6051 = 4.2849
    \end{align*}
    For $\alpha$ Centauri B:
    \begin{align*}
      M_{bolB} = m_{bolB} + 5 + 5\log p &= 1.09 + 5 + 5\log(0.7568) \\
        &= 1.09 + 5 - 0.6051 = 5.4849
    \end{align*}
  \item[Step 4:] We determine the mass of both stars using luminosity mass
    relation,
    \[ M_{bol} = -10.2 \log\left( \frac{\mathcal{M}}{\mathcal{M}_\odot}
       \right) + 4.9
       \;\Longrightarrow\;
       \mathcal{M} = 10^{\left( \frac{4.9 - M_{bol}}{10.2} \right)}
       \mathcal{M}_\odot \]
    For $\alpha$ Centauri A:
    \[ \mathcal{M}_A = 10^{\left( \frac{4.9 - M_{bolA}}{10.2} \right)}
       \mathcal{M}_\odot
       = 10^{\left( \frac{4.9 - 4.2849}{10.2} \right)} \mathcal{M}_\odot
       = 1.1490\,\mathcal{M}_\odot \]
    For $\alpha$ Centauri B:
    \[ \mathcal{M}_B = 10^{\left( \frac{4.9 - M_{bolB}}{10.2} \right)}
       \mathcal{M}_\odot
       = 10^{\left( \frac{4.9 - 5.4949}{10.2} \right)} \mathcal{M}_\odot
       = 0.8763\,\mathcal{M}_\odot \]
    % sic — the exponent shows 5.4949 where Step 3 gave M_bolB = 5.4849;
    % the printed result 0.8763 follows from 5.4849.
  \end{description}
  \hfill(30\%)

  \emph{2nd iteration:} Repeating the step 1 until 4 and using the previous
  resulting mass in the first iteration:
  \begin{description}
  \item[Step 1:] $\mathcal{M}_A + \mathcal{M}_B = 1.1490 + 0.8763 = 2.0253$
  \item[Step 2:]
    \[ p = \left[ \frac{\alpha^3}{(\mathcal{M}_A + \mathcal{M}_B) P^2}
           \right]^{1/3}
         = \left[ \frac{(17.59)^3}{(2.0253)(79.24)^2} \right]^{1/3}
         = 0.7536 \]
  \item[Step 3:] For $\alpha$ Centauri A:
    \[ M_{bolA} = m_{bolA} + 5 + 5\log p
       = -0.11 + 5 + 5\log(0.7536) = 4.2757 \]
    For $\alpha$ Centauri B:
    \[ M_{bolB} = m_{bolB} + 5 + 5\log p
       = 1.09 + 5 + 5\log(0.7536) = 5.4757 \]
  \item[Step 4:] For $\alpha$ Centauri A:
    \[ \mathcal{M}_A
       = 10^{\left( \frac{4.9 - 4.2757}{10.2} \right)} \mathcal{M}_\odot
       = 1.1513\,\mathcal{M}_\odot \]
    For $\alpha$ Centauri B:
    \[ \mathcal{M}_B
       = 10^{\left( \frac{4.9 - 5.4757}{10.2} \right)} \mathcal{M}_\odot
       = 0.8781\,\mathcal{M}_\odot \]
  \end{description}
  \hfill(10\%)

  \emph{3rd iteration:} Repeating the step 1 until 4 and using the previous
  resulting mass in the second iteration:
  \begin{description}
  \item[Step 1:] $\mathcal{M}_A + \mathcal{M}_B = 1.1513 + 0.8781 = 2.0294$
  \item[Step 2:]
    \[ p = \left[ \frac{(17.59)^3}{(2.0294)(79.24)^2} \right]^{1/3}
         = 0.75314 \]
  \item[Step 3:] For $\alpha$ Centauri A:
    \[ M_{bolA} = -0.11 + 5 + 5\log(0.7531) = 4.2751 \]
    For $\alpha$ Centauri B:
    \[ M_{bolB} = 1.09 + 5 + 5\log(0.7531) = 5.4751 \]
  \item[Step 4:] For $\alpha$ Centauri A:
    \[ \mathcal{M}_A
       = 10^{\left( \frac{4.9 - 4.2751}{10.2} \right)} \mathcal{M}_\odot
       = 1.1515\,\mathcal{M}_\odot \]
    For $\alpha$ Centauri B:
    \[ \mathcal{M}_B
       = 10^{\left( \frac{4.9 - 5.4751}{10.2} \right)} \mathcal{M}_\odot
       = 0.8782\,\mathcal{M}_\odot \]
  \end{description}

  \emph{4th iteration:} Repeating the step 1 until 4 and using the previous
  resulting mass in the third iteration:
  \begin{description}
  \item[Step 1:] $\mathcal{M}_A + \mathcal{M}_B = 1.1515 + 0.8782 = 2.0297$
  \item[Step 2:]
    \[ p = \left[ \frac{(17.59)^3}{(2.0297)(79.24)^2} \right]^{1/3}
         = 0.7531 \]
  \item[Step 3:] For $\alpha$ Centauri A:
    \[ M_{bolA} = -0.11 + 5 + 5\log(0.7531) = 4.2753 \]
    For $\alpha$ Centauri B:
    \[ M_{bolB} = 1.09 + 5 + 5\log(0.7531) = 5.4743 \]
  \item[Step 4:] For $\alpha$ Centauri A:
    \[ \mathcal{M}_A
       = 10^{\left( \frac{4.9 - 4.2753}{10.2} \right)} \mathcal{M}_\odot
       = 1.1517\,\mathcal{M}_\odot \]
    For $\alpha$ Centauri B:
    \[ \mathcal{M}_B
       = 10^{\left( \frac{4.9 - 5.4743}{10.2} \right)} \mathcal{M}_\odot
       = 0.8784\,\mathcal{M}_\odot \]
  \end{description}

  \emph{5th iteration:} Repeating the step 1 until 4 and using the previous
  resulting mass in the fourth iteration:
  \begin{description}
  \item[Step 1:] $\mathcal{M}_A + \mathcal{M}_B = 1.1517 + 0.8784 = 2.0301$
  \item[Step 2:]
    \[ p = \left[ \frac{(17.59)^3}{(2.0301)(79.24)^2} \right]^{1/3}
         = 0.7530'' \]
  \item[Step 3:] For $\alpha$ Centauri A:
    \[ M_{bolA} = -0.11 + 5 + 5\log(0.7530) = 4.2740 \]
    For $\alpha$ Centauri B:
    \[ M_{bolB} = 1.10 + 5 + 5\log(0.7530) = 5.4740 \]
    % sic — the source writes 1.10 here where every other iteration uses
    % m_bolB = 1.09.
  \item[Step 4:] For $\alpha$ Centauri A:
    \[ \mathcal{M}_A
       = 10^{\left( \frac{4.9 - 4.2740}{10.2} \right)} \mathcal{M}_\odot
       = 1.1518\,\mathcal{M}_\odot \]
    For $\alpha$ Centauri B:
    \[ \mathcal{M}_B
       = 10^{\left( \frac{4.9 - 5.4740}{10.2} \right)} \mathcal{M}_\odot
       = 0.8785\,\mathcal{M}_\odot \]
  \end{description}

  The results of the fifth iteration equal the results of the fourth
  iteration, so
  \begin{itemize}
  \item The mass of $\alpha$ Centauri A:
    $\mathcal{M}_A = 1.1518\,\mathcal{M}_\odot$
  \item The mass of $\alpha$ Centauri B:
    $\mathcal{M}_B = 0.8785\,\mathcal{M}_\odot$
  \end{itemize}
\end{description}

\solhead{8.19}{Light curves of stars}
% source id: 2010_DA_D2
\begin{description}
  \item[1)] \ioaafig{2010_DA_D2-s1.pdf}
  \ioaafig{2010_DA_D2-s2.pdf}
  \hfill(6p)

\item[2)]
  \[ \langle \Delta V \rangle = \frac{1}{n} \sum_{i=1}^{n} \Delta V_i
     = -0.248\,\mathrm{mag} \]
  \hfill(4p)
  \[ \langle \Delta R \rangle = \frac{1}{n} \sum_{i=1}^{n} \Delta R_i
     = 0.127\,\mathrm{mag} \]
  \hfill(4p)

\item[3)]
  \[ \sigma_{\Delta V} = \sqrt{ \frac{1}{n-1} \sum_{i=1}^{n}
     \left( \Delta V_i - \langle \Delta V \rangle \right)^2 }
     = 0.083\,\mathrm{mag} \]
  \hfill(4p)
  \[ \sigma_{\Delta R} = \sqrt{ \frac{1}{n-1} \sum_{i=1}^{n}
     \left( \Delta R_i - \langle \Delta R \rangle \right)^2 }
     = 0.011\,\mathrm{mag} \]
  \hfill(4p)

\item[4)] measured from the differences of times at the maximum values of the
  fits of the two peaks in V and R, respectively: 0.06 days, 0.06 days.
  \hfill(4p)

\item[5)] measured from the differences of magnitudes at the maximum values of
  the fits of the two peaks in V and R, respectively: 0.79 mag, 0.49 mag.
  \hfill(4p)

\item[6)] measured from the differences of times at the maximum values of the
  fits of the first peaks in V and R: $0(\pm 0.025)$ P. \hfill(5p)
\end{description}

\solhead{8.20}{Analysis of times of minima}
% source id: 2011_DA_D1
\begin{description}
\item[1)] The star is a W UMa type, so both primary and secondary minima were
  observed. If the depth of the minimum isn't given, there is no way of
  deciding which (primary or secondary it is).

\item[2)] From the data, the can be seen that on four nights more than one
  minimum was observed. This allows us to estimate half the period based on
  the difference between pairs of observations:

  \begin{center}
  \begin{tabular}{ccc}
    \hline
    minima & $\Delta T$ & average \\
    \hline
    2--1   & .1367 & $\Sigma/5 = .1367$ \\
    4--3   & .1372 & Period $P =$ \\
    10--9  & .1364 & obtained average \\
    11--10 & .1373 & $\times\,2 = \underline{0.2734}$ \\
    12--11 & .1359 & \\
    \hline
  \end{tabular}
  \end{center}

\item[3)] For a first approximation of the elements, we chose an initial epoch
  ($M_0$), for example the first observed minimum 2454092.4111. The
  approximate elements are then:
  \[ M_{\mathrm{calc}} = 2454092.4111 + 0.2734 \cdot E \]

\item[4)] For an exact determination of the the elements we draw an O--C
  diagram. The data for making the O--C diagram are calculates as follows:

  \begin{center}
  \begin{tabular}{cccccc}
    \hline
    No & $(M_{\mathrm{obs}})$ & $E = (M_{\mathrm{obs}} - M_0)/P$
       & $E$ (rounded & $M_{\mathrm{calc}}$
       & O--C $=$ \\
       & $2454000+$ & & to the nearest & $2454000.0+$
       & $M_{\mathrm{obs}} - M_{\mathrm{calc}}$ \\
       & & & whole or half & & \\
       & & & integer) & & \\
    \hline
    1  & 092.4111 & 0         & 0.0    & 092.4111 & 0.0000 \\
    2  & 092.5478 & 0.5       & 0.5    & 092.5478 & 0.0000 \\
    3  & 367.3284 & 1005.5497 & 1005.5 & 367.3148 & 0.0136 \\
    4  & 367.4656 & 1006.0516 & 1006.0 & 367.4515 & 0.0141 \\
    5  & 388.5175 & 1083.0519 & 1083.0 & 388.5033 & 0.0142 \\
    6  & 388.6539 & 1083.5508 & 1083.5 & 388.6400 & 0.0139 \\
    7  & 704.8561 & 2240.1061 & 2240.0 & 704.8271 & 0.0290 \\
    8  & 776.4901 & 2502.1178 & 2502.0 & 776.4579 & 0.0322 \\
    9  & 835.2734 & 2717.1262 & 2717.0 & 835.2389 & 0.0345 \\
    10 & 847.3039 & 2761.1295 & 2761.0 & 847.2685 & 0.0354 \\
    11 & 847.4412 & 2761.6316 & 2761.5 & 847.4052 & 0.0360 \\
    12 & 847.5771 & 2762.1287 & 2762.0 & 847.5419 & 0.0352 \\
    \hline
  \end{tabular}
  \end{center}

\item[5)] The O--C diagram should be plotted on graph paper.
  \textit{[Official 2011 DA solution (p.82, scan) contains no drawn O-C diagram; it only instructs plotting on graph paper]}

\item[6)] From the diagram, it can be seen that over 2762 epochs the value of
  (O--C) increased by 0.0355 days, therefore over one cycle (one epoch) the
  period was too short by:
  \[ 0.0355/2762 = 0.0000129 \]
  and the estimated period $P$ should be increased by this value. Taking into
  account that the graph can be read to 3 significant figures, the correction
  should also be to 3 significant figures.
  \[ P' = 0.2734 + 0.0000129 = 0.2734129 \]
  thus the corrected light elements of V1107 Cas are:
  \[ M_{\mathrm{calc}} = 2454092.4111 + 0.2734129 \cdot E \]

\item[7)] The correction of the period could also be determined directly from
  the table, using the mean value of O--C for the initial epoch ($E = 0$; 0.5)
  and the final epochs ($E = 2761$; 2761.5; 2762).

\item[8)] To calculate the ephemeris, we need to know the JD of a given day.
  Since dates and JD numbers are given in the observations, we can calculate
  that, for example on 1 January 2008 at 12:00 the JD was 2454467.0, while on
  1 January 2011 at 12:00 it was 2455563.0. Since 1 September 2011 is the
  243-rd day of the year, thus at 12:00 on that day JD 2455806.0 begins.

\item[9)] On 1 September 2011 at 18:00 (JD 2455806.25) there have been
  1713.8389 days or 6268.3176 epochs since the inital epoch $M_0$
  (2454092.4111). Thus the nearest (secondary) minimum which agrees with
  these elements will be:
  \[ \mathrm{JD}\ 2454092.4111 + 0.2734129 \cdot 6268.5 = 2455806.2999
     = \underline{19^{\mathrm{h}}\,12^{\mathrm{m}}\ \mathrm{UT}} \]
  The next (primary) minimum will be:
  \[ \mathrm{JD}\ 2454092.4111 + 0.2734129 \cdot 6269 = 2455806.4366
     = \underline{22^{\mathrm{h}}\,29^{\mathrm{m}}\ \mathrm{UT}} \]
  The next (secondary) minimum will be on the same night, at:
  \[ \mathrm{JD}\ 2455806.4111 + 0.2734129 \cdot 6269.5 = 2457520.5733 \]
  % sic — so printed in the source; the first term should read
  % 2454092.4111, giving JD 2455806.5733.
  that is on 2 September at
  $\underline{1^{\mathrm{h}}\,46^{\mathrm{m}}\ \mathrm{UT}}$.

  Earlier and later minima can be determined by adding or subtracting half a
  period, or $3^{\mathrm{h}}17^{\mathrm{m}}$.
\end{description}

\solhead{8.21}{Cepheids}
% source id: 2012_DA_D2
% Official solution (scan), 2012 DA solutions p. 1-2. The "Pontuation"
% point-allocation list is omitted per transcription policy.

\begin{description}
\item[Items 1, 2:] Figure 1.
  \ioaafig[0.5]{2012_DA_D2-s1.pdf}

\item[3:] \hfill
  \begin{description}
  \item[Star 1:] $m = 25.35$, $P = 48$ days, $M = -6.06 \pm 0.05$,
    $d = 19.1 \pm 0.4$ Mpc
  \item[Star 2:] $m = 26.25$, $P = 22$ days, $M = -5.08 \pm 0.05$,
    $d = 18.4 \pm 0.4$ Mpc
  \end{description}

\item[4:] $\Delta d = 0.67$ Mpc. If one considers that a typical galaxy has a
  100 Kpc radius, then it is possible that they are in the same galaxy,
  considering the 0.4 Mpc uncertainty of each distance.
\end{description}

\solhead{8.22}{Cepheid HV2257}
% source id: 2015_DA_D1
\begin{description}
\item[a.] The light curve. \hfill(10 points)
  \ioaafig{2015_DA_D1-s1.pdf}

\item[b.] The color curve. \hfill(10 points)
  \ioaafig{2015_DA_D1-s2.pdf}

\item[c.] The radial velocity curve. \hfill(10 points)
  \ioaafig{2015_DA_D1-s3.pdf}
\end{description}

\noindent\textbf{Solar absolute bolometric magnitude calculation}
\hfill(5 points)

$F_\odot$ is the solar flux at the distance 10 pc which corresponds to the
solar bolomatric absolute magnitude.
\[ F_\odot = \frac{L_\odot}{4\pi \cdot (10\,pc)^2}
   = \frac{3.96 \cdot 10^{26}}
          {4 \cdot 3.14 \cdot 100 \cdot 3.086^2 \cdot 10^{32}}
   = 3.31 \cdot 10^{-10}\,\mathrm{W/m^2} \]

\medskip
\noindent\textbf{Formula derivation} \hfill(15 points)

From Stefan--Boltzmann equation we have luminosity of star:
\[ L = 4\pi R^2 \sigma T^4, \]
here $R$ is radius of the star, $\sigma$ is the Stefan--Boltzmann constant,
and $T$ is effective temperature of star. Star flux at distance $d$ will be
equal:
\[ F = \frac{\sigma R^2 T^4}{d^2} \]
From Pogson definition we have:
\[ m_1 - m_2 = -2.5\log\frac{F_1}{F_2} \]
Or using the Sun as reference, a star's observed bolometric flux will be:
\[ F = F_\odot\,10^{-\frac{m_{bol} - M_{\odot bol}}{2.5}} \]

Consider two moments, say $t_1$ and $t_2$. It is better to choose the phase
$t_1$ and $t_2$ during which the star's expansion acceleration close to
constant, and the difference in magnitude and color as large as possible. At
the moment $t_1$, with measured temperature $T_1$ and radius $R_1$, absolute
bolometric flux will be:
\[ F_1 = \frac{\sigma R_1^2 T_1^4}{d^2}
   \tag*{(1a)} \]
Later, at moment $t_2$:
\[ F_2 = \frac{\sigma R_2^2 T_2^4}{d^2}
   \tag*{(1b)} \]
During this time, the star's atmosphere has expanded from $R_1$ to $R_2$:
\[ R_2 = R_1 + \Delta R \]
or:
\[ \frac{R_2}{R_1} = 1 + \frac{\Delta R}{R_1} \]
Reminding:
\[ \frac{F_2}{F_1}
   = \left( \frac{T_2}{T_1} \right)^4 \left( \frac{R_2}{R_1} \right)^2 \]
We have
\[ \frac{F_2}{F_1}
   = \left( \frac{T_2}{T_1} \right)^4
     \left( \frac{\Delta R}{R_1} + 1 \right)^2
   \tag*{(2)} \]
From Pogson definition, and (2) we can find:
\[ R_1 = \frac{\Delta R}
   {\left( \frac{T_1}{T_2} \right)^2 \cdot 10^{-\frac{m_2 - m_1}{5}} - 1} \]

\medskip
\noindent\textbf{Calculation of $\Delta R$ from the radial velocity graph}
\hfill(20 points)

For finding $\Delta R$, we can use radial velocity curve, taking two moments,
$t_1$ and $t_2$, between which the expansion acceleration can be assumed
constant,
\[ \Delta R = \frac{(v_2 - v_1)\,\Delta\varphi\,P}{2} \]
here $P$ is pulsation period of star and $\Delta\varphi$ is phase difference
in moment between $t_1$ and $t_2$.

\begin{quote}
Alternative method using integral / sum:
\[ \Delta R = \int_{t_1}^{t_2} v(t)\,dt \]
We can calculate the integral by drawing lines connecting two adjacent points,
calculating the area of the trapesium under the line segment and sum up for
all line segment between moment $t_1$ and $t_2$.
\end{quote}

So, from radial velocity curve we choose the part of the graph which is close
to linear:
\ioaafig{2015_DA_D1-s4.pdf}
\begin{align*}
  t_1 &= 0.75 \\
  t_2 &= 0.85
\end{align*}
Which corresponds to
\begin{align*}
  v(t_1) &= 272\ \mathrm{km/s} \\
  v(t_2) &= 230\ \mathrm{km/s}
\end{align*}
Now we can calculate $\Delta R$:
\[ \Delta R = \frac{(230000 - 272000)(0.85 - 0.75) \cdot 39.294 \cdot 86400}
   {2} = -7.13 \cdot 10^9\,\mathrm{m} \]

\medskip
\noindent\textbf{Temperatures and magnitudes from photometric data}
\hfill(10 points)

For moments $t_1$ and $t_2$: from light curve we have the magnitude $V$:
\begin{align*}
  V_1 &= 13.58 \\
  V_2 &= 12.55
\end{align*}
From color curve we have the color $V-R$:
\begin{align*}
  (V-R)_1 &= 0.81 \\
  (V-R)_2 &= 0.53
\end{align*}
From Fig.~1 we can find the temperatures:
\begin{align*}
  T_1 &= 4000\,\mathrm{K} \\
  T_2 &= 4750\,\mathrm{K}
\end{align*}
From table 4 we have bolometric corrections for these moments (by using
Table 4 with linear interpolation):
\begin{align*}
  BC_1 &= -1.23 \\
  BC_2 &= -0.52
\end{align*}
Now we can calculate bolometric magnitudes:
\begin{align*}
  m_{bol}(t_1) &= 13.58 - 1.23 = 12.35 \\
  m_{bol}(t_2) &= 12.55 - 0.52 = 12.03
\end{align*}

\medskip
\noindent\textbf{Calculate $R_1$} \hfill(10 points)

First we calculate star's observed flux:
\[ F_1 = 2.913 \cdot 10^{-10} \cdot 10^{-\frac{12.35 - 4.72}{2.5}}
   = 2.584 \cdot 10^{-13}\,\mathrm{W/m^2} \]
% sic — the source uses F_sun = 2.913e-10 here although the value computed
% above was 3.31e-10 W/m^2.
Then the radius of the star at the moment $t_1$ will be:
\[ R_1 = \frac{\Delta R}
   {\left( \frac{T_1}{T_2} \right)^2 \cdot 10^{-\frac{m_2 - m_1}{5}} - 1}
   = \frac{-7.13 \cdot 10^9}
   {\left( \frac{4000}{4750} \right)^2
    \cdot 10^{-\frac{12.03 - 12.35}{5}} - 1}
   = 4.0 \cdot 10^{10}\,\mathrm{m} \]

\medskip
\noindent\textbf{Calculate distance} \hfill(10 points)

\[ d = \sqrt{ \frac{\sigma \cdot R_1^2 \cdot T_1^4}{F_1} }
   = 3.0 \cdot 10^{20}\,\mathrm{m}
   = \frac{3.0 \cdot 10^{20}}{3.086 \cdot 10^{16}} = 9.72\ \mathrm{kpc} \]

Answer: 9.72 kpc

\solhead{8.23}{Binary Pulsar}
% source id: 2016_DA_D1
% Official solution, 2016 DA examination pp. 1-8. The per-item marking
% bullet lists and the half/full-credit tolerance tables are omitted per
% transcription policy.

\begin{description}
\item[(D1.1)] Graph Number: D1.1.
  \ioaafig{2016_DA_D1-s1.pdf}

\item[(D1.2)] See above (the best-fit ellipse is drawn on graph D1.1).

\item[(D1.3)] Values are determined from lengths of axes of ellipse and
  mid-point of $P$-axis. Error margins may be determined by estimating
  extreme ellipses covering the points with errorbars. Any reasonable method
  of estimating error margins will be accepted.
  \begin{align*}
    2P_t &= (1.34 \pm 0.04)\,\mu\mathrm{s}
      \quad [(13.4 \pm 0.4)\ \mathrm{cm\ on\ graph}] \\
    P_t &= (0.67 \pm 0.02)\,\mu\mathrm{s}
  \end{align*}
  \hfill(2.0)
  \[ P_0 = (7588.48 \pm 0.02)\,\mu\mathrm{s} \]
  \hfill(2.0)
  \begin{align*}
    2a_t &= (3.42 \pm 0.12)\ \mathrm{m\,s^{-2}}
      \quad [(17.1 \pm 0.6)\ \mathrm{cm\ on\ graph}] \\
    a_t &= (1.71 \pm 0.06)\ \mathrm{m\,s^{-2}}
  \end{align*}
  \hfill(3.0)

\item[(D1.4)] We can easily recover the orbital period ($P_B$) and the radius
  of the orbit ($r$) in a circular orbit:
  \[ a_t = \frac{4\pi^2}{P_B^2}\,r
     \qquad \therefore\ r = \frac{P_B^2 a_t}{4\pi^2} \]
  \hfill(1.0)
  \[ P_t = \frac{2\pi P_0}{P_B} \times \frac{r}{c}
     = \frac{2\pi P_0}{P_B c} \times \frac{P_B^2 a_t}{4\pi^2}
     = \frac{P_0 P_B a_t}{2\pi c} \]
  \[ \therefore\ P_B = \frac{P_t}{P_0}\,\frac{2\pi c}{a_t} \]
  \hfill(1.0)
  \[ r = \frac{a_t}{4\pi^2}
     \left( \frac{P_t}{P_0}\,\frac{2\pi c}{a_t} \right)^2 \]
  \hfill(1.0)
  \[ \therefore\ r = \left( \frac{P_t}{P_0} \right)^2 \frac{c^2}{a_t} \]
  \hfill(1.0)

  Alternative algebraic routes accepted.

\item[(D1.5)]
  \[ P_B = \frac{P_t}{P_0}\,\frac{2\pi c}{a_t}
     = \frac{0.67}{7588.48} \times
       \frac{2\pi \times 2.998 \times 10^8}{1.71}\ \mathrm{s}
     = 96\,260\ \mathrm{s} = 1.125\,70\ \mathrm{d} \]
  \hfill(1.0)
  \begin{align*}
    \Delta P_B &= P_B \sqrt{
      \left( \frac{\Delta P_t}{P_t} \right)^2
      + \left( \frac{\Delta P_0}{P_0} \right)^2
      + \left( \frac{\Delta a_t}{a_t} \right)^2 } \\
    &= 1.12570 \times \sqrt{
      \left( \frac{0.02}{0.67} \right)^2
      + \left( \frac{0.02}{7588.48} \right)^2
      + \left( \frac{0.06}{1.71} \right)^2 }\ \mathrm{d} \\
    &= 1.12570 \times 0.0461 \simeq 0.052\ \mathrm{d}
  \end{align*}
  \hfill(1.0)
  \[ \therefore\ P_B = (1.13 \pm 0.05)\ \mathrm{d} \]
  \hfill(1.0)
  \[ r = \left( \frac{P_t}{P_0} \right)^2 \frac{c^2}{a_t}
     = \left( \frac{0.67}{7588.48} \right)^2 \times
       \frac{\left( 2.998 \times 10^8 \right)^2}{1.71}\ \mathrm{m} \]
  \[ \therefore\ r = 4.097\,39 \times 10^8\ \mathrm{m}
     = 2.738\,90 \times 10^{-3}\ \mathrm{AU} \]
  \hfill(1.0)
  \begin{align*}
    \Delta r &= r \sqrt{
      \left( \frac{2\Delta P_t}{P_t} \right)^2
      + \left( \frac{2\Delta P_0}{P_0} \right)^2
      + \left( \frac{\Delta a_t}{a_t} \right)^2 } \\
    &= 2.738\,90 \times 10^{-3} \times \sqrt{
      \left( \frac{2 \times 0.02}{0.67} \right)^2
      + \left( \frac{2 \times 0.02}{7588.48} \right)^2
      + \left( \frac{0.06}{1.71} \right)^2 }\ \mathrm{AU} \\
    &= 2.738\,90 \times 10^{-3} \times 0.069\,25\ \mathrm{AU}
      \simeq 0.19 \times 10^{-3}\ \mathrm{AU}
  \end{align*}
  \hfill(1.0)
  \[ r = (2.74 \pm 0.19) \times 10^{-3}\ \mathrm{AU} \]
  \hfill(1.0)

  Errors in $P_B$ and $r$ can be also estimated as maximum possible (worst
  case) error. In such case, errors would be about 1.5 times the standard
  error calculated above ($\delta P_B = 0.07$, $\delta r = 0.25$).

\item[(D1.6)] Using these newly determined orbital parameters, we can
  calculate the angular orbital phase for each data point, i.e., for each
  pair of acceleration and period measured $(P, a)$.
  \[ \varphi = \tan^{-1}\left( \frac{-a}{a_t}\,
     \frac{P_t}{P - P_0} \right) \]
  Care has to be taken to choose the value of the phase from among $\varphi$,
  $\pi \pm \varphi$, $2\pi - \varphi$, depending on the sign of $\cos\varphi$
  and $\sin\varphi$.

  \begin{center}
  \begin{tabular}{ccccc}
    \hline
    Sr. no. & $T$ (tMJD) & $P$ ($\mu$s) & $a$ ($\mathrm{m\,s^{-2}}$)
      & $\varphi$ \\
    \hline
    1 & 5740.654 & 7587.8889 & $-0.92$ & $148.62^\circ$ \\
    4 & 5746.675 & 7588.5810 & $+1.67$ & $278.77^\circ$ \\
    6 & 5983.932 & 7587.8552 & $-0.44$ & $164.57^\circ$ \\
    8 & 6040.857 & 7589.1350 & $+0.00$ & $0.00^\circ$ \\
    9 & 6335.904 & 7589.1358 & $+0.00$ & $0.00^\circ$ \\
    \hline
  \end{tabular}
  \end{center}

\item[(D1.7a)]
  \[ \frac{T_1 - T_0}{P_B} = \frac{\varphi_1}{2\pi}
     \;\Rightarrow\; T_0 = T_1 - \frac{\varphi_1}{2\pi}\,P_B \]
  \hfill(1.0)
  \begin{align*}
    T_0 &= 5740.654 - \frac{148.62^\circ}{360^\circ}
      \times 1.12570\ \mathrm{tMJD} \\
    T_0 &= 5740.189\ \mathrm{tMJD}
  \end{align*}
  \hfill(1.0)

\item[(D1.7b)]
  \[ T_{\mathrm{calc}} = T_0
     + \left( n + \frac{\varphi}{360^\circ} \right) P_B \]
  where $n$ = Integer part of $[(T - T_0)/P_B]$.

  \begin{center}
  \begin{tabular}{cccccc}
    \hline
    Sr. No. & $T$ (tMJD) & $\varphi$ & $n$ & $T_{\mathrm{calc}}$ (MJD)
      & $T_{\mathrm{O-C}}$ (days) \\
    \hline
    1 & 5740.654 & $148.62^\circ$ & 0   & 5740.654 & $0.000$ \\
    4 & 5746.675 & $278.77^\circ$ & 5   & 5746.689 & $-0.014$ \\
    6 & 5983.932 & $164.57^\circ$ & 216 & 5983.855 & $0.077$ \\
    8 & 6040.857 & $0.00^\circ$   & 267 & 6040.751 & $0.106$ \\
    9 & 6335.904 & $0.00^\circ$   & 529 & 6335.684 & $0.220$ \\
    \hline
  \end{tabular}
  \end{center}

\item[(D1.7c)] Graph Number: D1.7.
  \ioaafig{2016_DA_D1-s2.pdf}

\item[(D1.7d)] A linear fit to the plot of $T_{\mathrm{O-C}}$ vs $n$ gives
  the offset of period per cycle (slope) and the shift in the zero-phase
  point (intercept). \hfill(2.0)

  From a linear fit,
  \begin{align*}
    \mathrm{Slope} &= 0.000\,43\ \mathrm{d}/n \\
    \mathrm{Intercept} &= -0.010\ \mathrm{d}
  \end{align*}
  \hfill(3.0)
  \[ T_{0,r} = 5740.189 - 0.010 = 5740.179\ \mathrm{tMJD} \]
  \hfill(1.0)
  \begin{align*}
    P_{B,r} &= (1.12570 + 0.00043)\ \mathrm{d} \\
            &= 1.126\,13\ \mathrm{d} \\
    P_B &= 1.1261\ \mathrm{d}
  \end{align*}
  \hfill(1.0)
\end{description}

\solhead{8.24}{Compact Object in a Binary System}
% source id: 2018_DA_D2
% Official solution, 2018 DA solutions pp. 10-14. The grading-standard
% paragraphs for the figures and per-item mark breakdowns are omitted per
% transcription policy; the source's equation numbers (15)-(35) are kept.

\noindent\textbf{2.1. APOGEE RV Measurements} \hfill[15 marks]

\begin{description}
\item[(D2.1.1)] \hfill[6 marks]\\
  Table 3 gives 3 epochs of RV measurements from MJD 56204 to 56233, and the
  maximum acceleration occurred during the first and the second observation:
  \[ a_{\max} = \frac{RV_2 - RV_1}{t_2 - t_1}
     = 2.89\ \mathrm{km/(s \cdot d)} \tag*{(15)} \]

\item[(D2.1.2)] \hfill[9 marks]\\
  An easy approach to estimate the companion mass is consider the force
  exerted on the star:
  \[ F = M_1 a = \frac{G M_1 M_2}{r^2}
     \;\Rightarrow\; M_2 = \frac{a r^2}{G} \tag*{(16)} \]
  Therefore, we need to estimate $a$ and $r$ properly. Here we can assume the
  acceleration of the primary star as $a = a_{\max}$, and the distance
  between the two star $r = (v_2 - v_1) \cdot (t_2 - t_1)/2$, and we can
  derive:
  \[ M_2 \sim \frac{(v_2 - v_1)^3 (t_2 - t_1)}{4G}
     = 1.5\ M_\odot \tag*{(17)} \]
  (Note: This is only a rough estimate. If examinees derive a $M_2$ within
  the range of $0.5\,M_\odot \sim 3\,M_\odot$ with clear and proper
  assumption, then his/her solution should be considered correct.)
\end{description}

\medskip
\noindent\textbf{2.2. TRES RV Measurements} \hfill[25 marks]

\begin{description}
\item[(D2.2.1)] \hfill[14 marks]\\
  Given the data in Table 4, we can plot out all the TRES RV measurements,
  with a sinusoidal function connecting all the points:
  \ioaafig{2018_DA_D2-s1.pdf}

  The best-fit parameters of period an RV semi-amplitude are:
  \[ P_{\mathrm{orb}} = 83.26 \pm 0.02\ \mathrm{d};\quad
     K = 45.02 \pm 0.03\ \mathrm{km/s} \tag*{(18)} \]
  (error bars are reported here for benefit of jury. They do not carry any
  points)

  Answers within the range of $P_{\mathrm{orb}} \in [81, 85]$ d and
  $K \in [44, 46]$ km/s could be considered correct.

\item[(D2.2.2)] \hfill[4 marks]\\
  The orbital radius satisfies that:
  \[ 2\pi R_{\mathrm{orb}} = P_{\mathrm{orb}} \cdot
     \frac{K}{\sin i_{\mathrm{orb}}}
     \;\Rightarrow\;
     R_{\mathrm{orb}} > \frac{P_{\mathrm{orb}} K}{2\pi}
     = 74\ R_\odot = 0.34\ \mathrm{AU} \tag*{(19)} \]

\item[(D2.2.3)] \hfill[7 marks]\\
  The mass function of a binary system is defined as:
  \[ f(M_1, M_2) = \frac{M_2^3 \sin^3 i_{\mathrm{orb}}}
     {(M_1 + M_2)^2} \tag*{(20)} \]
  This mass function is actually derived from Newton's law (force balance):
  \[ \frac{G M_1 M_2}{D^2}
     = M_1 \left( \frac{2\pi}{P} \right)^2 R_{\mathrm{orb}} \tag*{(21)} \]
  And the relation between orbital radius and binary separation ($D$):
  \[ R_{\mathrm{orb}} = \frac{M_2}{M_1 + M_2}\,D
     = \frac{P_{\mathrm{orb}} K}{2\pi \sin i_{\mathrm{orb}}} \tag*{(22)} \]
  Therefore, Equation 20 could be rewritten as:
  \[ f(M_1, M_2) = \frac{M_2^3 \sin^3 i_{\mathrm{orb}}}{(M_1 + M_2)^2}
     = \frac{K^3 P_{\mathrm{orb}}}{2\pi G}
     \approx 0.79\ M_\odot \tag*{(23)} \]
\end{description}

\medskip
\noindent\textbf{2.3. Stellar Parameters} \hfill[35 marks]

\begin{description}
\item[(D2.3.1)] \hfill[16 marks]\\
  The first (and the most direct) thing we could solve from the given
  parameter table is the distance of this system, which is:
  \[ d = \frac{1}{\pi} = 3.68\ \mathrm{kpc} \tag*{(24)} \]
  And the uncertainty of $d$ should be:
  \[ \frac{\Delta d}{d}
     = \sqrt{ \left( -\frac{\Delta\pi}{\pi} \right)^2 }
     \;\Rightarrow\;
     \Delta d = \frac{\Delta\pi}{\pi} \cdot d
     = 0.66\ \mathrm{kpc} \tag*{(25)} \]
  Therefore, the luminosity of the star should be:
  \[ L_1 = 4\pi d^2 \cdot F = 464\ L_\odot \tag*{(26)} \]
  And the uncertainty is:
  \[ \frac{\Delta L_1}{L_1}
     = \sqrt{ \left( \frac{\Delta F}{F} \right)^2
       + \left( \frac{2\Delta d}{d} \right)^2 }
     \;\Rightarrow\; \Delta L_1 = 173\ L_\odot \tag*{(27)} \]
  By using Stefan-Boltzmann law, we can calculate the radius of the star:
  \[ R_1 = \sqrt{ \frac{L_1}{\sigma T^4 \cdot 4\pi} }
     = 30\ R_\odot \tag*{(28)} \]
  With an uncertainty of:
  \[ \Delta R_1 = R_1 \cdot \sqrt{
       \left( -2\,\frac{\Delta T}{T} \right)^2
       + \left( \frac{1}{2}\,\frac{\Delta L_1}{L_1} \right)^2 }
     = R_1 \cdot \sqrt{
       \left( -2\,\frac{\Delta T}{T} \right)^2
       + \left( \frac{1}{2}\,\frac{\Delta F}{F} \right)^2
       + \left( \frac{\Delta d}{d} \right)^2 }
     = 6\ R_\odot \tag*{(29)} \]
  Combined with the observed rotational velocity, we can calculate the sine
  of inclination angle:
  \[ \sin i = \frac{P \cdot v\sin i}{2\pi R_1} = 0.77 \tag*{(30)} \]
  And the uncertainty should be:
  \[ \Delta \sin i = \sin i \cdot \sqrt{
       \left( -\frac{\Delta R_1}{R_1} \right)^2
       + \left( \frac{\Delta (v\sin i)}{v\sin i} \right)^2 }
     = 0.15 \tag*{(31)} \]
  Now we can calculate the mass $M_1$. First of all, we need to translate the
  surface gravity from log scale to linear scale:
  \[ g = 10^{\log g} = 1.58\ \mathrm{m/s^2} \tag*{(32)} \]
  With an uncertainty of:
  \[ \Delta g = g \ln 10 \cdot \Delta(\log g)
     = 0.36\ \mathrm{m/s^2} \tag*{(33)} \]
  Therefore, $M_1$ should be:
  \[ M_1 = \frac{g R_1^2}{G} = 5.2\ M_\odot \tag*{(34)} \]
  And the uncertainty is:
  \[ \Delta M_1 = M_1 \sqrt{
       \left( \frac{\Delta g}{g} \right)^2
       + \left( 2\,\frac{\Delta R_1}{R_1} \right)^2 }
     = 2.3\ M_\odot \tag*{(35)} \]

\item[(D2.3.2)] \hfill[4 marks]\\
  From the temperature and the luminosity, we can find that the primary star
  should be a (3) Red Giant.

\item[(D2.3.3)] \hfill[10 marks]\\
  Since we have derived the binary system mass function, with a given range
  of $M_1$, we can calculate the corresponding range of $M_2$ and hence make
  a plot of $M_2 - M_1$ (Figure 5).
  \ioaafig{2018_DA_D2-s2.pdf}

\item[(D2.3.4)] \hfill[5 marks]\\
  The vertical shadowed zone should start at $2.9\,M_\odot$ and end at
  $7.5\,M_\odot$, while the maximum mass of white dwarf is $1.4\,M_\odot$ and
  the maximum mass of neutron star should be around $2.4 \sim 2.6\,M_\odot$
  (here we adopt $2.5\,M_\odot$ on the plot). From Figure 5, we can readout
  that the possible mass of the companion is roughly
  $6^{+4}_{-2.5}\,M_\odot$. Since we didn't receive any optical radiation
  from the companion, as well as $M_2$ has already exceeded the maximum mass
  of neutron star, the invisible companion in this binary system could be a
  \textbf{stellar-mass black hole}, or some unknown compact object that
  current theory hasn't predicted.

  (Note: Examinee's answer is not required to give an uncertainty range of
  $M_2$: if the answer is within the possible range of
  $6^{+4}_{-2.5}\,M_\odot$, the answer will be considered correct. In
  addition to this, examinee don't need to include ``unknown compact
  object'' in his/her answer while addressing the possible type of the
  companion: an answer like ``black hole'' or ``stellar-mass black hole'' is
  enough for receiving full grades.)
\end{description}

\solhead{8.25}{Photometry and spectroscopy of Nova Del 2013}
% source id: 2019_DA_1
% Official solution, 2019 DA solutions pp. 1-8; the source's equation
% numbers (1.1)-(1.15) are kept.

\begin{description}
\item[a)] From the light curve plot the Modified Julian Dates can be read
  with an error of about 0.5. The peak of the maximum brightness is very
  narrow and obviously placed at $\mathrm{MJD}_0 = 56\,520.5$ with a value of
  $m_0 = 4.5^m$, therefore:
  \[ \mathrm{MJD}_0 = 56\,520.5 \pm 0.5 \]
  \hfill(2 p)
  \[ m_0 = 4.5^m \pm 0.05^m \]
  \hfill(1 p)

\item[b)] The brightness values of $2^m$ and $3^m$ decline are
  $6.5^m \pm 0.05^m$ and $7.5^m \pm 0.05^m$, respectively. \hfill(2 p)

  The corresponding Modified Julian Dates are $\mathrm{MJD}_2$ and
  $\mathrm{MJD}_3$. Because of the poorly defined slopes on the light curve
  around these dates, their acceptable error is larger than in other parts of
  the light curve, let say it is about $1^d$, so:
  \[ \mathrm{MJD}_2 = 56\,531.5 \pm 1, \qquad
     \mathrm{MJD}_3 = 56\,543.5 \pm 1 \]
  \hfill(2 p)
  \[ t_2 = 11^d \pm 1^d, \qquad t_3 = 23^d \pm 1^d \]
  \hfill(2 p)

\item[c)] The text of this part does not ask for calculating the individual
  errors of the formulae, but it is worth estimating them here, just for the
  sake of completeness. (Students won't do it.)
  \begin{description}
  \item[(a)] The form of the function is
    \[ M = a + b \arctan\frac{c - \log t_2}{d}, \qquad
       a = -7.92,\ b = -0.81,\ c = 1.32,\ d = 0.23,
       \tag*{(1.1)} \]
    so its derivative:
    \[ M' = -\frac{b}{d \log(10)
       \left[ 1 + \left( \frac{c - \log t_2}{d} \right)^2 \right]}
       \frac{1}{t_2}
       \;\rightarrow\;
       \Delta M = -\frac{b}{d \log(10)
       \left[ 1 + \left( \frac{c - \log t_2}{d} \right)^2 \right]}
       \frac{\Delta t_2}{t_2}
       \tag*{(1.2)} \]
    The value and its error calculated from the formulae above (error is not
    necessary):
    \[ M_0^{(a)} = -8.63^m \pm 0.06^m \]
    \hfill(1 p)
  \item[(b)] The form of the function is
    \[ M = a + b \log t_2, \qquad a = -11.32,\ b = 2.55
       \tag*{(1.3)} \]
    so its derivative:
    \[ M' = \frac{b}{t_2 \log(10)}
       \;\rightarrow\;
       \Delta M = \frac{b}{\log(10)} \frac{\Delta t_2}{t_2}
       \tag*{(1.4)} \]
    The value and its error calculated from the formulae above (error is not
    necessary):
    \[ M_0^{(b)} = -8.66^m \pm 0.10^m \]
    \hfill(1 p)
  \item[(c)] The form of the function is
    \[ M = a + b \log t_3, \qquad a = -11.99,\ b = 2.54
       \tag*{(1.5)} \]
    so its derivative:
    \[ M' = \frac{b}{t_3 \log(10)}
       \;\rightarrow\;
       \Delta M = \frac{b}{\log(10)} \frac{\Delta t_3}{t_3}
       \tag*{(1.6)} \]
    The value and its error calculated from the formulae above (error is not
    necessary):
    \[ M_0^{(c)} = -8.53^m \pm 0.05^m \]
    \hfill(1 p)
  \end{description}

  The standard deviation of a dataset can be calculated as
  \[ \sigma = \sqrt{ \frac{\sum\limits_{i=1}^{n} (x_i - \bar{x})^2}{n-1} },
     \tag*{(1.7)} \]
  where $n$ is the number of data points, $x_i$ is the $i$th individual data
  value and $\bar{x}$ is their mean,
  $\bar{x} = (x_1 + x_2 + \ldots + x_n)/n$.

  The mean and the standard deviation of the three absolute maximum
  brightness values:
  \[ M_0 = -8.61^m \pm 0.07^m \]
  \hfill(2 p)

\item[d)] The color excess $E(B-V)$ is the difference between the observed
  color index of the star and the intrinsic color index predicted from its
  spectral type:
  \[ E(B-V) = (B-V) - (B-V)_0 = A_B - A_V
     \tag*{(1.8)} \]
  The total extinction is quantified by $A_V$ (at 5550 \AA). The ratio of
  total-to-selective extinction:
  \[ R = \frac{A_V}{E(B-V)}
     \;\rightarrow\; A_V = R\,E(B-V), \text{ where } R = 3.1
     \tag*{(1.9)} \]
  \hfill(2 p)

  With the given value and error of $E(B-V)$:
  \[ A_V = 3.1 \times E(B-V) = 3.1 \times (0.184^m \pm 0.035^m)
     = 0.57^m \pm 0.11^m \]
  \hfill(2 p)

\item[e)] According to the formula for the distance modulus:
  \[ m_V - M_V = -5 + 5\log d + A_V \;\rightarrow\;
     \tag*{(1.10)} \]
  \hfill(1 p)
  \[ \log d = \frac{m_V - M_V + 5 - A_V}{5} \;\rightarrow\;
     \tag*{(1.11)} \]
  \[ d = 10^{(m_V - M_V + 5 - A_V)/5},
     \tag*{(1.12)} \]
  \hfill(2 p)\\
  where the distance $d$ is in parsecs.

  Since $\Delta a^x / \Delta x = a^x \ln a$, therefore the error of the
  distance $d$:
  \[ \Delta d = 10^{(m_V - M_V + 5 - A_V)/5} \times \ln 10 \times
     \Delta\left( (m_V - M_V + 5 - A_V)/5 \right) \]
  \hfill(2 p)

  The error of $(m_V - M_V + 5 - A_V)/5$ can be estimated with the sum of
  the errors of $m_V$, $M_V$ and $A_V$, so:
  \[ \Delta\left( (m_V - M_V + 5 - A_V)/5 \right) \approx 0.05 \]
  \hfill(2 p)

  With the data:
  \[ d = 3220\ \mathrm{pc} \quad\text{and}\quad \Delta d = 338\ \mathrm{pc}, \]
  \hfill(2 p)\\
  so the distance to the nova:
  \[ d \approx 3.2 \pm 0.3\ \mathrm{kpc} \]
  \hfill(2 p)

\item[f)] The well known Doppler formula between the wavelength displacement
  and radial velocity:
  \[ \frac{\lambda - \lambda_0}{\lambda_0} = \frac{\Delta\lambda}{\lambda_0}
     = \frac{v_r}{c}
     \;\rightarrow\; v_r = \frac{\Delta\lambda}{\lambda_0}\,c,
     \tag*{(1.13)} \]
  \hfill(1 p)\\
  where $\lambda$ is the measured wavelength of the line feature, $\lambda_0$
  is the rest wavelength of the line, $c$ is the speed of light, and $v_r$ is
  the radial velocity to be calculated.

  The wavelength of the P Cygni absorption peak should be extracted from the
  figure with an error of about 1 \AA. The main part of this error is coming
  from the ``definition'' of the peak position of the Gaussian-like profile.
  This will result in inaccuracy of about
  $\Delta v_r = \pm 50\ \mathrm{km\,s^{-1}}$ in radial velocities. \hfill(1 p)

  The wavelengths and radial velocities should be something like these:
  \begin{center}
  \begin{tabular}{ccc}
    \hline
    MJD & WL & RV \\
    \hline
    56518.986 & 6527 & $-1636$ \\
    56519.813 & 6531 & $-1454$ \\
    56520.843 & 6534 & $-1317$ \\
    56521.835 & 6537 & $-1179$ \\
    56522.829 & 6542 & $-951$ \\
    56523.827 & 6544 & $-860$ \\
    \hline
  \end{tabular}
  \end{center}
  \hfill(12 p)

  Radial velocities within the range of $\pm 50\ \mathrm{km\,s^{-1}}$ of the
  RV values listed in the table should be given full marks, but velocities in
  the range of $\pm 100\ \mathrm{km\,s^{-1}}$ are still acceptable with half
  marks.

\item[g)] See the attached figure as an example for the acceptable solution.
  The plotted data are taken from the table above. For the sake of simplicity
  the absolute values of the radial velocities have been used for making the
  graph. \hfill(6 p)
  \ioaafig{2019_DA_1-s3.pdf}

\item[h)] It is obvious from the plot, that the radial velocities lie along a
  straight line. To estimate the size of the expanding envelope we need to
  calculate the area below the $t - v_r$ graph between the first and last
  date. Hence the graph is a straight line, this is very simple: we have to
  determine the area of the hatched region which is a trapezoid.

  If the two bases and the height of the trapezoid are $a$, $b$, and $m$,
  respectively, then the area of the trapezoid is:
  \[ T = \frac{a + b}{2}\,m
     \tag*{(1.14)} \]
  \hfill(3 p)

  In our case $a = v_{r1}$, $b = v_{r6}$, and $m = t_6 - t_1$. \hfill(1 p)

  We could use the fitted line (dashed) as the upper side of the trapezoid,
  but this would be a bit complicated --- because of the difficulties of the
  fitting process ---, and not necessary at all. Instead of this we use the
  line connecting the first and last radial velocity points as this runs very
  close to the fitted line.
  \[ R = \int_{t_1}^{t_6} |v_r|\,dt
     \approx \frac{(v_{r1} + v_{r6})}{2}\,(t_6 - t_1)
     \approx 3.5\ \mathrm{AU} \]
  The result: $R \approx 3.5$ AU \hfill(3 p)

\item[i)] The apparent angular diameter of the spherical envelope seen from
  the Earth:
  \[ \vartheta = 2 \times \arctan\left( \frac{R}{d} \right)
     \approx 2\,\frac{R}{d}
     \tag*{(1.15)} \]
  \hfill(2 p)

  Using the values of $d \approx 3.2$ kpc and $R \approx 3.5$ AU, 5 days
  after the discovery the angular diameter of the envelope is:
  \[ \vartheta = 0.0022'' = 2.2\ \mathrm{mas} \]
  \hfill(2 p)

  A less formal solution:
  \begin{itemize}
  \item By definition a parsec (1 pc) is the distance from the Sun to an
    astronomical object that has a parallax angle of one arcsecond, i.e.\ it
    represents the distance at which the radius of Earth's orbit (1 AU)
    subtends an angle of one arcsecond.
  \item Because of the very small angles the distance is a linear function of
    the parallax. This means that the radius $R \approx 3.5$ AU of the
    envelope subtends one arcsecond viewing from a distance of
    $d \approx 3.5$ pc, and one milliarcsecond from a distance of
    $1000d \approx 3500\ \mathrm{pc} = 3.5\ \mathrm{kpc}$.
  \item Since this value is close to the distance of Nova Del 2013 determined
    earlier, we can conclude that the apparent angular diameter of the
    spherical shape envelope 5 days after the discovery was about 2
    milliarcseconds.
  \end{itemize}

\ioaafig{2019_DA_1-s1.pdf}
\ioaafig{2019_DA_1-s2.pdf}
\end{description}

\solhead{8.26}{Triply eclipsing hierarchical triple stellar system}
% source id: 2019_DA_2
% Official solution, 2019 DA solutions pp. 8-15; the source's equation
% numbers (2.1)-(2.55) are kept. Photometric phases are written varphi_{1,2}
% and position angles phi_{1,2}, as in the source.

\begin{description}
\item[i)] \hfill
  \begin{description}
  \item[a)] See the following table. \hfill(10 p)

    \begin{center}
    \begin{tabular}{cccccc}
      \hline
      event no. & contacts & components & BJD & $\varphi_1$ & $\varphi_2$ \\
      \hline
      1 & I   & A, B & 2455476.1096 & 0.1226 & 0.9747 \\
        & II  & A, C & 2455476.4245 & 0.4703 & 0.9816 \\
        & III & A, B & 2455477.9677 & 0.1743 & 0.0155 \\
        & IV  & A, B & 2455478.4722 & 0.7313 & 0.0266 \\
      2 & I   & A, B & 2455521.5217 & 0.2643 & 0.9734 \\
      3 & III & A, C & 2455568.9434 & 0.6248 & 0.0163 \\
      4 & I   & A, C & 2455612.4733 & 0.6882 & 0.9736 \\
        & III & A, C & 2455614.3571 & 0.7682 & 0.0150 \\
      5 & III & A, B & 2455659.9241 & 0.0808 & 0.0171 \\
        & IV  & A, C & 2455660.2422 & 0.4320 & 0.0241 \\
      \hline
    \end{tabular}
    \end{center}

  \item[b)] See the following table. \hfill(5 p)

    \begin{center}
    \begin{tabular}{cc}
      \hline
      event no. & closer component \\
      \hline
      1 & A \\
      2 & A \\
      3 & A \\
      4 & A \\
      5 & A \\
      \hline
    \end{tabular}
    \end{center}

    All these events occur close to $\varphi_2 = 0$ (or $\varphi_2 = 1$)
    which means by definition that star A eclipses stars B and C. Therefore,
    star A is closer to the observer.

  \item[c)] In the moments of the 1st and last (4th) contacts:
    \[ R_A + R_{B,C} = d_{\mathrm{A-B,C}},
       \tag*{(2.1)} \]
    \hfill(1 p)\\
    while in the case of the 2nd and 3rd (inner) contacts:
    \[ R_A - R_{B,C} = d_{\mathrm{A-B,C}},
       \tag*{(2.2)} \]
    \hfill(1 p)\\
    where $d_{\mathrm{A-B,C}}$ stands for the sky-projected distance of the
    disks of the occulting star A and occulted component B or C.

    Let $\vec{r}_1$ the radius vector directed from star B toward star C,
    and $\vec{r}_2$ another radius vector directed from the centre of mass
    of stars B and C toward star A. By the use of these two (Jacobian)
    vectors, the position vectors connecting stars B and C with star A, can
    be written as:
    \[ \vec{r}_{\mathrm{BA}} = \vec{r}_A - \vec{r}_B
       = \vec{r}_2 + \frac{m_C}{m_{BC}}\,\vec{r}_1,
       \tag*{(2.3)} \]
    \hfill(1 p)
    \[ \vec{r}_{\mathrm{CA}} = \vec{r}_A - \vec{r}_C
       = \vec{r}_2 - \frac{m_B}{m_{BC}}\,\vec{r}_1.
       \tag*{(2.4)} \]
    \hfill(1 p)

    Using the facts that ($i_{\mathrm{rel}} = 0^\circ$) and, furthermore,
    $i_1 = i_2 = 90^\circ$, it is worthy to introduce the orbital plane as
    reference frame. Let the origin of our frame of reference be the centre
    of mass of the close binary (denoted as $O_2$ in figure above).
    \ioaafig{2019_DA_2-s1.pdf}
    Furthermore, let axis $X$ directed toward the line of sight. Therefore,
    axis $Y$ is located in the tangential plane of the sky.

    In this frame of reference, the components of vectors $\vec{r}_1$ and
    $\vec{r}_2$ can simply be written as
    \[ \vec{r}_1 = a_1 \left[ \cos(\phi_1); \sin(\phi_1) \right],
       \tag*{(2.5)} \]
    \hfill(1 p)
    \[ \vec{r}_2 = a_2 \left[ \cos(\phi_2); \sin(\phi_2) \right],
       \tag*{(2.6)} \]
    \hfill(1 p)\\
    as $\phi_1$ takes the values of $2k\pi$ in the moments, when star C
    eclipses star B, and similarly, $\phi_2 = 2k\pi$ when star A
    ``eclipses'' the centre of mass of the inner binary.

    One can also notice that there is a simple relation between these
    position angles and the photometric phases defined above as
    \[ \phi_{1,2} = 2\pi\varphi_{1,2}
       \tag*{(2.7)} \]
    \hfill(1 p)

    We are interested in the sky-projected distances of the stellar disks,
    which, in this frame of reference, are equal to the $y$ components of
    the vector equations (2.3 and 2.4). Accordingly,
    \[ d_{\mathrm{A-B}} = \left| a_2 \left[ \sin(2\pi\varphi_2)
       + \frac{m_C}{m_{BC}}\,\frac{a_1}{a_2}\,\sin(2\pi\varphi_1)
       \right] \right|,
       \tag*{(2.8)} \]
    \hfill(1 p)
    \[ d_{\mathrm{A-C}} = \left| a_2 \left[ \sin(2\pi\varphi_2)
       - \frac{m_B}{m_{BC}}\,\frac{a_1}{a_2}\,\sin(2\pi\varphi_1)
       \right] \right|.
       \tag*{(2.9)} \]
    \hfill(1 p)

    \emph{From this point we can use different combinations of eclipsing
    events given in the table. Therefore, here we use one possible order
    only for illustration.}

    We notice that both outer contacts (i.e.\ contacts I and IV) of the
    first event stars A and B are involved. Therefore, we can write the same
    equation for both contacts as follows:
    \[ \frac{R_A + R_B}{a_2} = \left| \sin(2\pi\varphi_2)
       + \frac{m_C}{m_{BC}}\,\frac{a_1}{a_2}\,\sin(2\pi\varphi_1) \right|.
       \tag*{(2.10)} \]
    \hfill(1 p)\\
    Therefore, we have two independent equations and two unknown variables,
    namely
    \[ \frac{R_A + R_B}{a_2}, \quad\text{and}\quad
       \frac{m_C}{m_{BC}}\,\frac{a_1}{a_2} = \frac{a_B}{a_2}.
       \tag*{(2.11)} \]
    \hfill(2 p)

    The computation is as follows: for the 1st contact
    $\sin 2\pi\varphi_2 = -0.158\,296$,
    $\sin 2\pi\varphi_1 = 0.696\,364$, while for the 4th one
    $\sin 2\pi\varphi_2 = 0.166\,356$,
    $\sin 2\pi\varphi_1 = -0.993\,105$. Consequently
    \[ 0.158\,296 - 0.696\,364\,\frac{a_B}{a_2}
       = 0.166\,356 - 0.993\,105\,\frac{a_B}{a_2},
       \tag*{(2.12)} \]
    \hfill(1 p)
    \[ \frac{a_B}{a_2} = 0.027\,161
       \tag*{(2.13)} \]
    \hfill(1 p)\\
    and, therefore,
    \[ \frac{R_A + R_B}{a_2} = 0.1394.
       \tag*{(2.14)} \]
    \hfill(1 p)

    Now, considering the also available 3rd contact moment of the same event
    (in which similarly, stars A and B play the roles), one can get
    \[ \frac{R_A - R_B}{a_2}
       = \left| \sin(2\pi \times 0.0155)
       + 0.027\,161 \times \sin(2\pi \times 0.1743) \right|
       \tag*{(2.15)} \]
    \hfill(1 p)
    \[ \frac{R_A - R_B}{a_2} = 0.1214
       \tag*{(2.16)} \]
    \hfill(1 p)

    Combining these results, we obtain that
    \[ \frac{R_A}{a_2} = \frac{1}{2}\left( \frac{R_A + R_B}{a_2}
       + \frac{R_A - R_B}{a_2} \right) = 0.1304
       \tag*{(2.17)} \]
    \hfill(2 p)
    \[ \frac{R_B}{a_2} = \frac{1}{2}\left( \frac{R_A + R_B}{a_2}
       - \frac{R_A - R_B}{a_2} \right) = 0.0090
       \tag*{(2.18)} \]
    \hfill(2 p)

    We repeat a very similar calculation for the 3rd contacts of events 3
    and 4, in which cases star B is substituted with star C. For the 3rd
    contact of the third event $\sin 2\pi\varphi_2 = 0.102\,237$,
    $\sin 2\pi\varphi_1 = -0.706\,218$, while the 3rd contact of the fourth
    event $\sin 2\pi\varphi_2 = 0.094\,108$,
    $\sin 2\pi\varphi_1 = -0.993\,469$. Therefore,
    \[ 0.102\,237 + 0.706\,218\,\frac{a_C}{a_2}
       = 0.094\,108 + 0.993\,469\,\frac{a_C}{a_2}
       \tag*{(2.19)} \]
    \hfill(1 p)
    \[ \frac{a_C}{a_2} = 0.028\,298
       \tag*{(2.20)} \]
    \hfill(1 p)\\
    and, accordingly,
    \[ \frac{R_A - R_C}{a_2} = 0.1222
       \tag*{(2.21)} \]
    \hfill(1 p)

    Taking into account that $R_A/a_2$ is already known from the previous
    stage, the dimensionless relative radius of star C can be obtained
    simply as
    \[ \frac{R_C}{a_2} = \frac{R_A}{a_2} - \frac{R_A - R_C}{a_2}
       = 0.1304 - 0.1222 = 0.0082
       \tag*{(2.22)} \]
    \hfill(2 p)

    The other possibility is, however, that e.g.\ from the 1st contact of
    event 4 we calculate the sum $(R_A + R_C)/a_2$, as
    \[ \frac{R_A + R_C}{a_2}
       = \left| \sin(2\pi \times 0.9736)
       - 0.028\,298 \times \sin(2\pi \times 0.6882) \right|
       \tag*{(2.23)} \]
    \[ \frac{R_A + R_C}{a_2} = 0.1389,
       \tag*{(2.24)} \]
    and then we obtain that
    \[ \frac{R_A}{a_2} = \frac{1}{2}\left( \frac{R_A + R_C}{a_2}
       + \frac{R_A - R_C}{a_2} \right) = 0.1306,
       \tag*{(2.25)} \]
    \[ \frac{R_C}{a_2} = \frac{1}{2}\left( \frac{R_A + R_C}{a_2}
       - \frac{R_A - R_C}{a_2} \right) = 0.0084.
       \tag*{(2.26)} \]

    \begin{quote}
    \emph{Several additional, equivalently acceptable scenarios.}

    We can write three equations for $(R_A + R_B)/a_2$:
    \begin{align*}
      \frac{R_A + R_B}{a_2}
        &= 0.158\,296 - 0.696\,364\,\frac{a_B}{a_2}
        && \text{(1st event Ist contact)} \\
      \frac{R_A + R_B}{a_2}
        &= 0.166\,356 - 0.993\,105\,\frac{a_B}{a_2}
        && \text{(1IV)} \\
      \frac{R_A + R_B}{a_2}
        &= 0.166\,433 - 0.995\,966\,\frac{a_B}{a_2}
        && \text{(2I)}
    \end{align*}
    Combining them one can get:
    \begin{center}
    \begin{tabular}{ccc}
      \hline
      contacts & $a_B/a_2$ & $(R_A + R_B)/a_2$ \\
      \hline
      1I--1IV & 0.027\,161 & 0.1394 \\
      1I--2I  & 0.027\,159 & 0.1394 \\
      1IV--2I & 0.026\,988 & 0.1396 \\
      \hline
    \end{tabular}
    \end{center}
    Similarly, for $(R_A - R_B)/a_2$:
    \begin{align*}
      \frac{R_A - R_B}{a_2}
        &= 0.097\,235 + 0.889\,001\,\frac{a_B}{a_2}
        && \text{(1III)} \\
      \frac{R_A - R_B}{a_2}
        &= 0.107\,236 + 0.486\,152\,\frac{a_B}{a_2}
        && \text{(5III)}
    \end{align*}
    Combining them one can get:
    \begin{center}
    \begin{tabular}{ccc}
      \hline
      contacts & $a_B/a_2$ & $(R_A - R_B)/a_2$ \\
      \hline
      1III--5III & 0.024\,826 & 0.1193 \\
      \hline
    \end{tabular}
    \end{center}
    Now for $(R_A + R_C)/a_2$:
    \begin{align*}
      \frac{R_A + R_C}{a_2}
        &= 0.165\,116 - 0.925\,553\,\frac{a_C}{a_2}
        && \text{(4I)} \\
      \frac{R_A + R_C}{a_2}
        &= 0.150\,847 - 0.414\,376\,\frac{a_C}{a_2}
        && \text{(5IV)}
    \end{align*}
    The only combination gives:
    \begin{center}
    \begin{tabular}{ccc}
      \hline
      contacts & $a_C/a_2$ & $(R_A + R_C)/a_2$ \\
      \hline
      4I--5IV & 0.027\,914 & 0.1393 \\
      \hline
    \end{tabular}
    \end{center}
    Finally, for $(R_A - R_C)/a_2$:
    \begin{align*}
      \frac{R_A - R_C}{a_2}
        &= 0.115\,353 + 0.185\,529\,\frac{a_C}{a_2}
        && \text{(1II)} \\
      \frac{R_A - R_C}{a_2}
        &= 0.102\,237 + 0.706\,218\,\frac{a_C}{a_2}
        && \text{(3III)} \\
      \frac{R_A - R_C}{a_2}
        &= 0.094\,108 + 0.993\,469\,\frac{a_C}{a_2}
        && \text{(4III)}
    \end{align*}
    The combinations are:
    \begin{center}
    \begin{tabular}{ccc}
      \hline
      contacts & $a_C/a_2$ & $(R_A - R_C)/a_2$ \\
      \hline
      1II--3III  & 0.025\,190 & 0.1200 \\
      1II--4III  & 0.026\,295 & 0.1202 \\
      3III--4III & 0.028\,299 & 0.1222 \\
      \hline
    \end{tabular}
    \end{center}
    \emph{End additional, equivalently acceptable scenarios.}
    \end{quote}

    Now, we are in the position to calculate the inner mass ratio $q_1$, as
    follows:
    \[ q_1 = \frac{a_B/a_2}{a_C/a_2}
       = \frac{0.027\,161}{0.028\,298} = 0.9598
       \tag*{(2.27)} \]
    \hfill(2 p)

    Furthermore, we can also get the ratio of the semi-major axes as
    \[ \frac{a_1}{a_2} = \frac{a_B}{a_2} + \frac{a_C}{a_2}
       = 0.027\,161 + 0.028\,298 = 0.055\,459
       \tag*{(2.28)} \]
    \hfill(2 p)

    The requested results with the range of the full marks (summary):
    \begin{align*}
      0.1150 &\le \frac{R_A}{a_2} \le 0.1450 \\
      0.0075 &\le \frac{R_B}{a_2} \le 0.0105 \\
      0.0070 &\le \frac{R_C}{a_2} \le 0.0100 \\
      0.045 &\le \frac{a_1}{a_2} \le 0.065 \\
      0.80 &\le q_1 \le 1.25
    \end{align*}

  \item[d)] The quickest method to obtain $q_2$ comes with the double use of
    Kepler's third law. Writing this law for both the inner and outer
    orbits, and dividing them one can get, that
    \[ \left( \frac{a_1}{a_2} \right)^3 \left( \frac{P_2}{P_1} \right)^2
       = \frac{m_{BC}}{m_{ABC}} = \frac{q_2}{1 + q_2}.
       \tag*{(2.29)} \]
    \hfill(4 p)\\
    Therefore
    \[ q_2 = \frac{\left( \frac{a_1}{a_2} \right)^3
       \left( \frac{P_2}{P_1} \right)^2}
       {1 - \left( \frac{a_1}{a_2} \right)^3
       \left( \frac{P_2}{P_1} \right)^2}
       \tag*{(2.30)} \]
    \hfill(2 p)
    \[ q_2 \approx \frac{0.055\,459^3 \times 50.206\,763^2}
       {1 - 0.055\,459^3 \times 50.206\,763^2} \approx 0.7543
       \tag*{(2.31)} \]
    \hfill(2 p)
  \end{description}

\item[ii)] The circular velocity of star A is
  \[ v_A = \frac{2\pi}{P_2}\,r_A.
     \tag*{(2.32)} \]
  \hfill(1 p)\\
  The amplitude of the RV curve is equal to maximum value of the
  line-of-sight component of that velocity, i.e.
  \[ K_A = v_A \sin i_2.
     \tag*{(2.33)} \]
  \hfill(1 p)

  The sinusoidal component in the occurrence of the eclipsing minima times
  comes from the light-travel time effect caused by the revolution of stars
  B and C around the centre of mass of the whole triple system. In this
  regard the motion of components B and C can simply be substituted by the
  movement of their centre of mass along the outer orbit. Therefore, the
  total variation of the inner binary's distance to the Earth is the
  line-of-sight component of the diameter of the orbit of the centre of mass
  of the inner binary around the centre of mass of the complete triple
  system, i.e.
  \[ \Delta z_{BC} = 2 r_{BC} \sin i_2.
     \tag*{(2.34)} \]
  \hfill(2 p)\\
  Therefore, the amplitude of the light-travel time sine wave is
  \[ A_{\mathrm{ETV}} = \frac{\Delta z}{2c}
     = \frac{r_{BC} \sin i_2}{c}
     \;\rightarrow\;
     r_{BC} = \frac{c A_{\mathrm{ETV}}}{\sin i_2}.
     \tag*{(2.35)} \]
  \hfill(2 p)

  Such a way, using that
  \[ q_2 = \frac{m_{BC}}{m_A} = \frac{r_A}{r_{BC}},
     \tag*{(2.36)} \]
  \hfill(1 p)\\
  one can obtain, that
  \[ q_2 = \frac{P_2 K_A}{2\pi c A_{\mathrm{ETV}}}.
     \tag*{(2.37)} \]
  \hfill(1 p)\\
  which gives a second, independent determination of the outer mass ratio
  $q_2$.

  The uncertainty can be estimated either as
  \[ \Delta q_2 = q_2 \sqrt{
       \left( \frac{\Delta P_2}{P_2} \right)^2
       + \left( \frac{\Delta K_A}{K_A} \right)^2
       + \left( \frac{\Delta A_{\mathrm{ETV}}}{A_{\mathrm{ETV}}} \right)^2 }
     \tag*{(2.38)} \]
  or
  \[ \Delta q_2 = q_2 \left(
       \left| \frac{\Delta P_2}{P_2} \right|
       + \left| \frac{\Delta K_A}{K_A} \right|
       + \left| \frac{\Delta A_{\mathrm{ETV}}}{A_{\mathrm{ETV}}} \right|
       \right)
     \tag*{(2.39)} \]
  \hfill(2 p)

  With numerical values:
  \[ q_2 = 0.621 \pm 0.047 \quad\text{or}\quad q_2 = 0.621 \pm 0.048 \]
  \hfill(1 p)

  Then, the semi-major axis of the outer orbit can be calculated in the
  following alternative ways:
  \[ a_2 = r_A\,\frac{1 + q_2}{q_2}
     = \frac{P_2}{2\pi}\,\frac{K_A}{\sin i_2}\,\frac{1 + q_2}{q_2}
     \tag*{(2.40)} \]
  \[ a_2 = r_{BC}\,(1 + q_2)
     = \frac{c A_{\mathrm{ETV}}}{\sin i_2}\,(1 + q_2)
     \tag*{(2.41)} \]
  \[ a_2 = r_A + r_{BC}
     = \frac{P_2}{2\pi}\,\frac{K_A}{\sin i_2}
     + \frac{c A_{\mathrm{ETV}}}{\sin i_2},
     \tag*{(2.42)} \]
  One of the three equations above: \hfill(2 p)

  We assumed that $i_2 = 90^\circ$, therefore $\sin i_2 = 1$. Then
  \[ a_2 = (60.712 \pm 2.851) \times 10^6\ \mathrm{km}
     = (87.293 \pm 4.099)\ R_\odot
     = (0.406 \pm 0.019)\ \mathrm{AU} \]
  \hfill(2 p)

  The uncertainties from the three different formulae above:
  \[ \Delta a_2 = a_2 \sqrt{
       \left( \frac{\Delta P_2}{P_2} \right)^2
       + \left( \frac{\Delta K_A}{K_A} \right)^2
       + \left( \frac{\Delta q_2}{q_2 (1 + q_2)} \right)^2 }
     \tag*{(2.43)} \]
  \[ \Delta a_2 \approx 2.851 \times 10^6\ \mathrm{km}
     \approx 4.10\ R_\odot \approx 0.019\ \mathrm{AU} \]
  \[ \Delta a_2 = a_2 \sqrt{
       \left( \frac{\Delta A_{\mathrm{ETV}}}{A_{\mathrm{ETV}}} \right)^2
       + \left( \frac{\Delta q_2}{1 + q_2} \right)^2 }
     \tag*{(2.44)} \]
  \[ \Delta a_2 \approx 4.946 \times 10^6\ \mathrm{km}
     \approx 7.11\ R_\odot \approx 0.033\ \mathrm{AU} \]
  \[ \Delta a_2 = \sqrt{
       \left( \frac{\Delta P_2}{2\pi}\,K_A \right)^2
       + \left( \frac{P_2}{2\pi}\,\Delta K_A \right)^2
       + \left( c\,\Delta A_{\mathrm{ETV}} \right)^2 }
     \tag*{(2.45)} \]
  \[ \Delta a_2 \approx 2.849 \times 10^6\ \mathrm{km}
     \approx 4.10\ R_\odot \approx 0.019\ \mathrm{AU} \]
  One of the three formulae above with its uncertainty: \hfill(3 p)

  In what follows, we use the smallest uncertainties (first row above) but
  the others, and also the nonquadratic ones are acceptable, too.

  The total mass of the triple now can be calculated directly from Kepler's
  third law or, as an alternative way, one can obtain it from the equations
  on the centripetal accelerations.
  \[ r_A\,\frac{4\pi^2}{P_2^2} = \frac{G m_{BC}}{a_2^2}
     \tag*{(2.46)} \]
  \[ r_{BC}\,\frac{4\pi^2}{P_2^2} = \frac{G m_A}{a_2^2},
     \tag*{(2.47)} \]
  where the only unknown quantities are the masses, so they can be
  calculated separately from the two equations or, summing the masses we can
  get back (formally) Kepler's third law:
  \[ m_{ABC} = \frac{4\pi^2}{G}\,\frac{a_2^3}{P_2^2}
     = (4.312 \pm 0.607)\ M_\odot,
     \tag*{(2.48)} \]
  \hfill(2 p)\\
  from which
  \[ m_A = \frac{m_{ABC}}{1 + q_2} = (2.660 \pm 0.383)\ M_\odot
     \tag*{(2.49)} \]
  \hfill(1 p)
  \[ m_{BC} = m_{ABC}\,\frac{q_2}{1 + q_2} = (1.652 \pm 0.245)\ M_\odot
     \tag*{(2.50)} \]
  \hfill(1 p)

\item[iii)] We obtained the total mass ($m_{BC}$) of the inner binary in the
  second part of the problem. Combining this with the inner mass ratio
  ($q_1$) obtained in the first part, one can get that
  \[ m_B = \frac{m_{BC}}{1 + q_1}
     = \frac{1.652\ M_\odot}{1 + 0.960} = 0.843\ M_\odot
     \tag*{(2.51)} \]
  \hfill(3 p)\\
  and
  \[ m_C = q_1 \times m_B = 0.960 \times 0.843\ M_\odot = 0.809\ M_\odot
     \tag*{(2.52)} \]
  \hfill(3 p)

  The dimensionless radius of each star relative to the semi-major axis was
  calculated in the first part, while the semi-major axis of the outer orbit
  ($a_2$) was obtained in the second part. Their multiplication gives the
  physical radii of the stars:
  \[ R_A = 0.1304 \times 87.29\ R_\odot = 11.381\ R_\odot
     \tag*{(2.53)} \]
  \hfill(3 p)
  \[ R_B = 0.0090 \times 87.29\ R_\odot = 0.786\ R_\odot
     \tag*{(2.54)} \]
  \hfill(3 p)
  \[ R_C = 0.0082 \times 87.29\ R_\odot = 0.712\ R_\odot
     \tag*{(2.55)} \]
  \hfill(3 p)
\end{description}

\solhead{8.27}{Colours of the binary epsilon Trianguli Australis}
% source id: 2012_OBS-TEL_3
% Official solution: 2012 observational exam, 1st night, task 3
% ($IOAA_PAPERS_DIR/observation/2012_OBS_sol.pdf, page 1).

\begin{description}
\item[Brighter:] White/blue (\quad) \qquad Yellow (X) \qquad
  Red (\quad) \hfill (K0III)
\item[Dimmer:] White/blue (X) \qquad Yellow (\quad) \qquad
  Red (\quad) \hfill (B9IV)
\end{description}

If the student is not able to point to the star the evaluator will do so, but
student loses half of the points.
